\documentclass{report}
\newif\ifsolo

\solofalse

\usepackage[T1]{fontenc} 
\usepackage[utf8]{inputenc}
\usepackage[french]{babel}
\usepackage[a4paper, margin=2cm]{geometry}
\usepackage[activate={true,nocompatibility},final,kerning=true,spacing=true,factor=1100,stretch=10,shrink=10]{microtype}
\usepackage{lmodern}
% \usepackage{mlmodern} 
\usepackage{amsmath}
\usepackage{bbm}
\usepackage{bm}
\usepackage{amsthm}
\usepackage{amssymb}
\usepackage[cal=pxtx]{mathalpha}
\usepackage{mathtools}
\usepackage{centernot}
\usepackage{changepage}
\usepackage{enumitem}
\usepackage{stmaryrd}
\usepackage{xcolor}
\usepackage{mdframed}
\usepackage{etoolbox}
\usepackage{xparse}
\usepackage{parskip}
\usepackage{tikz,tkz-tab}
\usepackage{pgfplots}
\usepackage[bottom,stable]{footmisc}
\usepackage{titleps}
\ifsolo
\else
\usepackage{lastpage}
\usepackage[nohints,french]{minitoc}
\fi
\usepackage{tocloft}
\usepackage{enumitem}
\usepackage{needspace}
\usepackage[noautomatic]{imakeidx}
\usepackage[normalem]{ulem}
\usepackage[hidelinks]{hyperref}

\makeindex[intoc,columns=2]%,options={-s indexstyle.ist -o build/\jobname.ind}]
\makeatletter
\def\@idxitem{\par\hangindent 0pt}
\makeatother

\AtBeginDocument{\def\labelitemi{$\bullet$}}
\setlist[itemize]{leftmargin=.5cm}
\setlist[enumerate]{leftmargin=.7cm}

\newif\ifishypconc
\ishypconcfalse
\newcommand{\newthm}[2]{
    \newenvironment{#1}[1][]
    {\par\begin{minipage}{\textwidth}\noindent\normalfont\ifstrempty{##1}{\textbf{#2.}}{\textbf{#2} (##1).}\par\nopagebreak%
    \begin{mdframed}[
        linewidth=1pt,
        linecolor=black,
        bottomline=false,topline=false,rightline=false,
        innerrightmargin=0pt,innertopmargin=0pt,innerbottommargin=0pt,
        innerleftmargin=1em,
        skipabove=.4\baselineskip
    ]}
    {\ifishypconc \end{adjustwidth} \fi\end{mdframed}\end{minipage}\ishypconcfalse}
}


\theoremstyle{definition}
\newthm{thm}{Théorème}
\newthm{thmdef}{Théorème -- Définition}
\newthm{dfn}{Définition}
\newthm{defprop}{Définition -- Proposition}
\newthm{lmm}{Lemme}
\newthm{rem}{Remarque}
\newthm{res}{Résultat}
\newthm{prop}{Proposition}
\newthm{cor}{Corollaire}
\newthm{exo}{Exercice}
\newthm{notation}{Notation}
\newtheorem*{ex}{Exemple}
\newtheorem{csq}{Conséquence}[subsection]
\newtheorem*{csq*}{Conséquence}
\renewcommand{\thecsq}{\arabic{csq}}

\AtBeginEnvironment{thmdef}{\Needspace{5\baselineskip}}
\AtBeginEnvironment{defprop}{\Needspace{5\baselineskip}}

\renewcommand{\thesection}{\arabic{section}}
\renewcommand{\thesubsection}{\emph{\alph{subsection}})}
\renewcommand{\thesubsubsection}{\emph{\roman{subsubsection}})}


\setlength{\cftsecnumwidth}{0.6cm}
\renewcommand{\cftsecpresnum}{}
\renewcommand{\cftsecaftersnum}{}


\ifsolo
\addto\captionsfrench{
  \renewcommand{\contentsname}%
    {Plan du cours}%
}
\else
\renewcommand{\thechapter}{\Roman{chapter}}
\mtcsetoffset{minitoc}{-2em}
\renewcommand{\cftsecfont}{\small\bfseries}
\setcounter{tocdepth}{1}
\setlength{\cftchapnumwidth}{1.4cm}
\fi

\newpagestyle{main}{
  \setheadrule{1pt}
  \ifsolo\else \setfootrule{1pt}\fi
  \sethead{\thesection ~--~\textsc{\sectiontitle}}{}{Page \thepage\ifsolo\else~/~\pageref{LastPage}\fi}
  \setfoot{\ifsolo\else\textbf{Chapitre \thechapter} : \textsc{\chaptertitle}\fi}{}{}
}

\newif\ifmrmspace
\NewDocumentCommand\genmrm{O{#2}m}{
    \expandafter\NewDocumentCommand\csname#1\endcsname{m}{%
        \mrmspacetrue%
        % don't put extra space if followed by ()_^
        % to get Ker(u), Ker u, (u+id), Sp_K, ...
        \ifstrequal{##1}{(}{\mrmspacefalse}{}%
        \ifstrequal{##1}{)}{\mrmspacefalse}{}%
        \ifstrequal{##1}{_}{\mrmspacefalse}{}%
        \ifstrequal{##1}{^}{\mrmspacefalse}{}%
        \mathrm{#2}%
        \ifmrmspace\;\fi##1%
    }%
}
\genmrm{Ker}
\genmrm{dom}
\genmrm{Vect}
\genmrm{Sp}
\genmrm[Img]{Im}
\genmrm{id}
\genmrm{rg}
\genmrm{ord}
\genmrm{Tr}
\genmrm{Com}
\genmrm{Cov}
\genmrm{Conv}
\genmrm{diag}
\genmrm{Hom}
\genmrm{sh}
\genmrm{ch}
\genmrm{Aut}
\NewDocumentCommand\CPM{O{0}}{\mathcal{C}^{#1}_{{P\!M}}}
\newcommand\diff{\mathrm{d}}
\newcommand\nop{\|\hspace{-1pt}|}
\newcommand\defeq{\;\underset{\text{def}}=\;}
\newcommand\1{\mathbbm{1}}
\renewcommand\Im{\mathfrak{Im}}
\renewcommand\Re{\mathfrak{Re}}
\renewcommand\mathring[1]{\overset\circ{#1}}
\renewcommand\bar[1]{\overline{#1}}
\renewcommand\leq\leqslant
\renewcommand\epsilon\varepsilon
\renewcommand\emptyset\varnothing
\renewcommand\geq\geqslant
\newcommand\floor[1]{\left\lfloor#1\right\rfloor}
\newcommand\ceil[1]{\left\lceil#1\right\rceil}
\DeclareDocumentCommand\P{}{\mathbf P}
\DeclareDocumentCommand\R{}{\mathbf R}
\DeclareDocumentCommand\K{}{\mathbf K}
\DeclareDocumentCommand\N{}{\mathbf N}
\DeclareDocumentCommand\Z{}{\mathbf Z}
\DeclareDocumentCommand\Q{}{\mathbf Q}
\DeclareDocumentCommand\C{}{\mathbf C}
\DeclareDocumentCommand\E{}{\mathbf E}
\newcommand\indep{\protect\mathpalette{\protect\indepT}{\perp}}
\def\indepT#1#2{\mathrel{\rlap{$#1#2$}\mkern2mu{#1#2}}}
\newcommand\suit{\hookrightarrow}

\makeatletter
\newcommand*{\xnrightarrow}[2][]{%
  \ext@arrow 0359\nrightarrowfill@{#1}{#2}%
}
\newcommand*{\nrightarrowfill@}{%
    \narrowfill@\relbar\relbar\rightarrow{\not \relbar}
}
\newcommand*{\narrowfill@}[5]{%
  $\m@th\thickmuskip0mu\medmuskip\thickmuskip\thinmuskip\thickmuskip
  \relax#5#1\mkern-8mu%
  \cleaders\hbox{$#5\mkern-2mu#2\mkern-2mu$}\hfill
  \mkern-5mu %
  #4%
  \mkern-5mu %
  \cleaders\hbox{$#5\mkern-2mu#2\mkern-2mu$}\hfill
  \mkern-7mu#3$%
}
\makeatother

\newcommand{\todo}[1]{{\color{red}À faire: #1}}

\newcommand*\circled[1]{\tikz[baseline=(char.base)]{
            \node[shape=circle,draw,inner sep=2pt] (char) {#1};}~}

\newcommand{\hyp}{\circled{H}}
\newcommand{\conc}{\circled{C}}

\newcommand{\Hyp}{\ishypconctrue\begin{adjustwidth}{3em}{}~\\[.1cm] \hspace*{-3em}\hyp\hspace{1em}}
\newcommand{\Conc}{~\\[.1cm] \hspace*{-3em}\conc\hspace{1em}}
\newenvironment{concenum}{
    ~\\[0.1cm]\hspace*{-3em}\begin{tabular}{@{}lp{\textwidth-4em}} \conc & \vspace{-.45cm}
\begin{enumerate}[leftmargin=.5cm]
}
{
\end{enumerate}
\end{tabular}
}

\newcommand*{\colorboxed}{}
\def\colorboxed#1#{%
  \colorboxedAux{#1}%
}
\newcommand*{\colorboxedAux}[3]{%
  % #1: optional argument for color model
  % #2: color specification
  % #3: formula
  \begingroup
    \colorlet{cb@saved}{.}%
    \color#1{#2}%
    \boxed{%
      \color{cb@saved}%
      #3%
    }%
  \endgroup
}

\tocloftpagestyle{empty}
\pagestyle{main}

\pgfplotsset{compat=1.16}

\NewDocumentCommand\endchapter{}{%
\vfill
\[\star \qquad \quad \star \qquad \quad \star \qquad \quad \star\]
\[\star \qquad \quad \star \qquad \quad \star \qquad \quad \star \qquad \quad \star\]
\vfill
}


\begin{document}

\pagestyle{empty}
\begin{center}

    ~ \vfill

    \textsc{\LARGE Lycée Faidherbe}\\[1.5cm]
    \textsc{\Large Classe de MP$^\star$}\\[0.5cm]

    \rule{\linewidth}{0.5mm} \\[0.4cm]
    {\huge \bfseries Cours de Mathématiques }\\[0.1cm]
    \rule{\linewidth}{0.5mm} \\[0.4cm]
    \textsc{\large Programme officiel et compléments}\\[1.5cm]

    % \begin{minipage}{0.4\textwidth}
    %     \begin{flushleft} \large
    %         \emph{Professeur}\\
    %         Serge \textsc{Varjabédian}
    %     \end{flushleft}
    %     \end{minipage}\hspace{1cm}\begin{minipage}{0.4\textwidth}
    %     \begin{flushright} \large
    %         \emph{Prise de notes} \\
    %         Théophile \textsc{Cailliau}
    %     \end{flushright}
    % \end{minipage}\\[2cm]
    \vfill
    \includegraphics[width=5cm]{faidherbe.pdf}\\[1cm]
    \vfill
    {\large Année 2019-2020}\\
\end{center}

\dominitoc
\tableofcontents
\faketableofcontents

\begin{minipage}{\textwidth}
Dans tout le cours, les éléments explicitement au programme sont précédés par \textbf{Théorème}, \textbf{Définition}, \textbf{Propriété}, \textbf{Proposition}. Les remarques, corollaires et conséquences sont généralement importants. Les \textbf{Résultats} ou éléments démontrés sans indication quant à leur nature sont des compléments de cours (hors programme). Tous les compléments de cours ont leur intérêt, mais il n'est pas nécessaire de tous les connaître.

Les démonstrations manquantes sont jugées suffisamment facile à retrouver rapidement. Lorsqu'une démonstration est admise, elle apparaît toujours sous la forme :

\begin{proof}
    Admis(e).
\end{proof}
\end{minipage}

\newpage

\pagestyle{main}

\chapter{Analyse asymptotique}

\ifsolo
    ~

    \vspace{1cm}

    \begin{center}
        \textbf{\LARGE Analyse Asymptotique} \\[1em]
    \end{center}
    \tableofcontents
\else
    \minitoc
\fi
\thispagestyle{empty}

\ifsolo \newpage \setcounter{page}{1} \fi
\section{Développements limités -- Rappels}

Cette section ne contient que des rappels de sup.

\begin{dfn}[Développement limité]\index{développement limité}
    Un développement limité de $f:V\longrightarrow \mathbb R$ à l'ordre $n$ au voisinage $V$ de $a$ (qu'on abrègera $\mathrm{DL}_n(a)$ de $f$) est une écriture du type: \[
        f(x)=\underbrace{a_0+a_1(x-a)+\cdots +a_n(x-a)^n}_{\text{partie polynômiale}}+o_a((x-a)^n)
    \]
\end{dfn}

\begin{rem}
    Lorsqu'un $\mathrm{DL}_n(a)$ existe, il y a unicité de la partie polynômiale.
\end{rem}

\begin{rem}
    Les DL usuels sont à connaître.

\begin{align*}
    e^x &= 1 + x + \frac{x^2}{2} +\ldots+ \frac{x^n}{n!} &+ o_0(x^n) \\
%              = \sum_{k=0}^n \frac{x^n}{n !} + x^{n+1}\varepsilon(x) \\
    \sin x &= x - \frac{x^3}{3!}+\frac{x^5}{5!} +\ldots+(-1)^p\frac{x^{2p+1}}{(2p+1)!}
           & +o_0(x^{2p+1} ) \\
%              = \sum_{k=0}^p (-1)^p\frac{x^{2p+1}}{(2p+1)!} + x^{2p+2}\varepsilon(x) \\
    \cos x &= 1 - \frac{x^2}{2!}+\frac{x^4}{4!} +\ldots+(-1)^p\frac{x^{2p}}{(2p)!}
           & +o_0(x^{2p}) \\
    \frac{1}{1-x} &= 1 + x + x^2 + + x^3 + \ldots + x^n
        & + o_0(x^n) \\
    \ln (1-x) &= -x - \frac{x^2}{2} - \frac{x^3}{3} - \frac{x^4}{4} - \ldots - \frac{x^n}{n}
        & + o_0(x^n) \\
    (1+x)^\alpha &= 1 + \alpha x + \alpha (\alpha-1) \frac{x^2}{2!} + \ldots \\
        & \qquad\qquad +\alpha(\alpha-1)\ldots(\alpha-n+1)\frac{x^n}{n!}
        & + o_0(x^n) \\
\end{align*}
Les règles de calcul sont aussi à connaître (vues en sup: composition, troncature, multiplication)
\end{rem}

\begin{dfn}[Voisinage\index{voisinage}]
    Un voisinage de $a\in\mathbb R$ est un intervalle du type $]a-\varepsilon, a+\varepsilon[$. Si $V$ est un voisinage de $a$, on note $V\in\mathcal V(a)$.
\end{dfn}

Si $V\in\mathcal V(0)$ et si pour $x\in V\setminus\{0\}$, $f(x)=a_0+a_1x +a_px^p+o_p(x^p)$, $a_p\neq 0$ et $p\geq 2$ alors \begin{itemize}
    \item $f$ est prolongeable en $0$ et $f(0)=a_0$
    \item $f$ est dérivable en $0$ et $f'(0)=a_1$
    \item $y=a_0+a_1x$ est l'équation cartésienne de la tangente au graphe de $f$ en $0$
    \item $f(x)-a_0-a_1x$ a le même signe que $a_px^p$ au voisinage de $0$.
\end{itemize}

\begin{rem}
    $f$ admet un $\mathrm{DL}_1(a)$ si et seulement si $f$ est dérivable en $a$. Ça n'est pas vrai pour les ordres supérieurs de dérivation.
\end{rem}

\begin{ex}
    On pose \[
        f: x\in]-\pi;\pi[\setminus\{0\} \longmapsto \frac 1x-\frac1{\sin x}
    \]
    et on a (en utilisant les DL usuels et les règles de calcul) \[
        \forall x\in D_f,\quad f(x)=\frac1x-\frac1{x-\frac{x^3}3+o_0(x^3)}=\frac1x \left( 1- \left( 1+\frac {x^2}6+o_0(x^2) \right)\right)=-\frac16x+o_0(x)
    \]
    donc $f$ est prolongeable en $0$ avec $f(0)=0$, $f'(0)=-\frac 16$.

    La fonction $f$ est $\mathcal C^1$ sur son domaine de définition, l'est-elle en $0$ ? On a \[
        \forall x\in D_f, \qquad f'(x)=-\frac1{x^2}+\frac{\cos x}{\sin^2x}=\frac1{x^2} \left( -1+\frac{1-\frac{x^2}2+o_0(x^2)}{(1-\frac{x^2}6+o_0(x^2))^2} \right)=-\frac16+o_0(1)\xrightarrow[x\to 0]{}-\frac16
    \]
    donc $f$ est $\mathcal C^1$ en $0$
\end{ex}

\section{Taylor-Young}

\begin{thm}[Taylor-Young\index{Taylor-Young (théorème de -- )}]
    \Hyp $f\in\mathcal C^n(V\in\mathcal V(0), \mathbb R)$
    \Conc \[
            \forall x\in V, \qquad f(x)=\sum_{k=0}^n\frac{f^{(k)}(0)}{k!}x^k+o_0(x^n)
    \]
\end{thm}

\begin{rem}
    On n'a en fait besoin que de $f\in\mathcal C^{n-1}(V, \mathbb R)$ et $f^{(n-1)}$ dérivable en $0$
\end{rem}

\begin{rem}
    On ne peut en général pas dériver termes à termes. Si $f$ est $\mathcal C^n$, on peut dériver termes à termes pour obtenir le $\mathrm{DL}_{n-1}(0)$ de $f'$.
\end{rem}

\begin{ex}[$a$]
    Pour obtenir un $\mathrm{DL}_3(1)$ de $\arctan$, on note $f=\arctan$ et \[
        f(1)=\frac\pi4\qquad f'(1)=\frac12\qquad f''(1)=-\frac12\qquad f^{(3)}(1)=\frac1{2}
    \]
    donc \[
        \forall x\in \left]-\frac\pi2;\frac\pi2  \right[, \qquad f(x)=\frac\pi4+\frac{x-1}2-\frac{(x-1)^2}4+\frac{(x-1)^3}{12}+o_1((x-1)^3)
    \]
\end{ex}

\begin{ex}[$b$]
    On note $f\in\mathcal C^{\infty}([-1; 1], \mathbb R)$ telle que $\forall n\in\mathbb N^\star$, \begin{align*}
        f \left( \frac1n \right) &= \frac{\sqrt n}{\sqrt{n+1}+\sqrt n}\\
                                 &= \frac1{1+\sqrt{1+\frac1n}}\\
                                 &= \frac{\sqrt{1+\frac1n}-1}{\frac 1n}\\
                                 &= n \left( \sum_{k=1}^{2020}\binom {1/2}{k}\frac1{n^k}+o_{+\infty}(\frac{1}{n^{2020}}) \right) \\
                                 &= \sum_{k=0}^{2019}\frac{f^{(k)}(0)}{k!}\frac1{n^k}+o_{+\infty}(\frac1{n^{2019}})
    \end{align*}
    Puis $n\to +\infty$ donne \[
        \binom{1/2}1=f(0)
    \]
    et $nf(1/n)$ avec $n\to+\infty$ donne \[
        \binom{1/2}2=\frac{f'(0)}{1!}
    \]
    et en itérant on obtient \[
        \binom{1/2}{2020}=\frac{f^{(2019)}(0)}{2019!}
    \]
\end{ex}

\section{Utilisation standard des développements limités}

\subsection{Calcul de limite}

On va calculer \[
    \lim_{x\to 0^+}\frac{x^{\sh x}-(\sh x)^x}{(\sin x)^x-x^{\sin x}}
\]
Et on a \[
    \sh x=x+\frac{x^3}6+o_0(x^3)\qquad\qquad \sin x=x-\frac{x^3}6+o_0(x^3)
\]
donc, après calculs, \[
    \frac{x^{\sh x}-(\sh x)^x}{(\sin x)^x-x^{\sin x}} = \frac{\frac{x^3}6\ln x+o_0(x^3\ln x)-\frac{x^3}6+o_0(x^3)}{\frac{-x^3}6+o_0(x^3)+\frac{x^3}6\ln x+o_0(x^3\ln x)}\underset{0+}\sim -1
\]

\subsection{Développement limité d'une réciproque}

\begin{exo}
    Soit $x\in\mathbb R$. Montrer $\exists!\psi(x)\defeq y\in\mathbb R, e^{x+y}+y=1$. Montrer que $\psi$ admet un $\mathrm{DL}_2(0)$ que l'on donnera
\end{exo}

\begin{proof}[Résolution (Idée 1)]
    On étudie $y\longmapsto e^{x+y}+y-1$ à $x$ fixé. Problème: comment trouver un $\mathrm{DL}_2$ en $x$ ?
\end{proof}

\begin{proof}[Résolution (Idée 2)]
    «On perturbe le problème.» \[
        e^{x+y}+y=1\iff e^{x+y}+x+y=1+x \iff \varphi(x+y)=1+x
    \]
    avec $\varphi: t\longmapsto e^t+t\in\mathcal C^{\infty}(\mathbb R, \mathbb R)$ strictement croissante de dérivée jamais nulle. C'est donc une bijection de $\mathbb R$ dans $\mathbb R$ de bijection réciproque dérivable: \[
        \varphi^{-1'}=\frac1{\varphi'\circ \varphi^{-1}}\in\mathcal C^1
    \]
    donc $\varphi^{-1}\in\mathcal C^2$ (en fait $\mathcal C^{\infty}$) donc $\varphi^{-1}$ admet un $\mathrm{DL}_2(a)$ en $a\in\mathbb R$.

    On a $\psi(x)=-x+\varphi^{-1}(1+x)$ donc $\psi$ a un $\mathrm{DL}_2(0)$. Donc on écrit \[
        \psi(x)=a_0+a_1x+a_2x^2+o_0(x^2)
    \]
    et \[\psi(\varphi(x)-1)=-\varphi(x)+1+x=-x-\frac{x^2}2+o_0(x^2)=\psi \left(2x+\frac{x^2}2+o_0(x^2)\right)\]
    donc \[
        -x+\frac{x^2}2+o_0(x^2)=a_0+2a_1x+\left(\frac{a_1}2+4a_2\right)x^2+o_0(x^2)
    \]
    et on identifie les coefficients
\end{proof}

\section{Utilisation non standard des développements limités}

\subsection{Calul de primitive}

On veut calculer une primitive de \[
    F:x\longmapsto \frac1{x^{2019}(1+x^2)}.
\]
On a \[
    F(x)=\sum_{k=0}^{1010}(-1)^kx^{2k-2019}+o_0(x)
\]
donc \[
    F(x)-\sum_{k=0}^{1009}(-1)^kx^{2k-2019}=o_0(1)
\]
est une fraction rationnelle sans partie polaire en $0$, donc \[
    \sum_{k=0}^{1009}(-1)^kx^{2k-2019}
\]
est la partie polaire de $F$ en $0$ et \[
    F(x)=\sum_{k=0}^{1009}(-1)^kx^{2k-2019}+\frac{\alpha x+\beta}{1+x^2}
\]
et en injectant $i$, $\alpha i+\beta=i$ d'où $\alpha=1, \beta=0$. Puis on peut intégrer termes à termes

\subsection{Développement limité d'un terme général}

On note $f$ une fonction $\mathcal C^2$ sur un voisinage $V$ de $0$ telle que $f(0)=0$. On définit pour $n$ suffisament grand, \[
    u_n=\sum_{k=1}^nf\left(\frac{k}{n^2}\right)
\]
On va faire l'étude de $u_n$. \[
    \forall k\in\llbracket 1, n\rrbracket, \quad f \left( \frac{k}{n^2} \right)=\frac{f'(0)k}{n^2}+\frac{f''(0)}{2}\cdot \frac {k^2}{n^4}+\frac {k^2}{n^4}\cdot \varepsilon \left( \frac{k}{n^2} \right), \qquad \qquad \varepsilon (x)\xrightarrow[x\to0]{}0
\]
donc \[
    u_n=\frac{f'(0)}{n^2}\sum_{k=1}^nk+\underbrace{\frac{f''(0)}{2n^4}\sum_{k=1}^nk^2}_{\longrightarrow 0}+\underbrace{\frac1{n^4}\sum_{k=1}^nk^2\varepsilon \left( \frac k{n^2} \right)}_{\longrightarrow 0\text{ car }\varepsilon\text{ bornée }}\xrightarrow[n\to+\infty]{}\frac{f'(0)}2
\]

\subsection{Polynôme sous contrainte}

\begin{exo}
Trouver un polynôme $P_n$ tel que $P_n^2(x)=1+x+o_0(x^n)$. En déduire une racine carrée de \[
    A= \begin{pmatrix}
        4&1\\-9&-2
    \end{pmatrix}
\]
\end{exo}
\begin{proof}
    \[
        P_n=\sum_{k=1}^n\binom{1/2}{k}x^k
    \]
    convient et \[
        A=I_n+ \underbrace{\begin{pmatrix}
            3&1\\-9&-3
    \end{pmatrix}}_{\defeq N},\qquad \qquad N^2=0
    \]
    donc $P_1^2(N)=I_n+N+\frac{N^2}4=I_n+N=A$ donc $P_1(N)$ convient.
\end{proof}

\subsection{Formule combinatoire}

On va montrer que si $0\leq m\leq n$ alors \[
    \sum_{k=0}^n(-1)^{n-k}\binom nkk^m=\delta_{n,m}n!
\]
On écrit \[
    (e^x-1)^n=\sum_{k=0}^n\binom nk (-1)^{n-k}e^{kx}=(x+o(x))^n=x^n+o_0(x^n)
\]
et pour $m<n$, on identifie le coefficient de $x^m$ (il est nul). Sinon, $n=m$ et le résultat est le bon.

\section{Majorer, minorer, encadrer}

\subsection{Rappels}

\[
    |x|\leq a\iff -a\leq x\leq a
\]

\[
    \begin{cases}
        x\leq x'\\ y\leq y'
    \end{cases}
    \implies \begin{cases}
        x+y\leq x'+y'
    \end{cases}
\]
\[
    \begin{cases}
        0\leq x\leq x'\\ 0\leq y\leq y'
    \end{cases}
    \implies xy\leq x'y'
\]
\[
    ||x|-|y||\leq |x+y|\leq |x|+|y|
\]

\subsection{Sommation terme à terme}

Si $a_1\leq b_1, \cdots, a_n\leq b_n$ alors \[
    a_1+\cdots +a_n\leq b_1+\cdots b_n
\]
avec égalité ssi $a_1=b_1,\cdots a_n=b_n$. Le résultat se généralise identiquement aux séries convergentes.

\begin{ex}
    pour $n\geq 3$, \[
        \frac1{n^3}\leq \frac1{n(n-1)(n-2)}=\frac1{2n}-\frac1{n-1}+\frac1{2(n-2)}
    \]
    et en sommant, \[
        \zeta(3)=\sum_{k=1}^{+\infty}\frac1{k^3}< \frac{11}8
    \]
\end{ex}

\subsection{Inégalité de Cauchy-Schwartz}

\index{Cauchy-Schwartz (inégalité de -- )}

On va montrer l'inégalité de Cauchy-Schwartz de deux manières différentes


\textbf{Identité de Lagrange.} On note $a_1, \cdots, a_n, b_1, \cdots, b_n\in\mathbb R$. \begin{align*}
        \left( \sum_{i=1}^na_i^2 \right) \cdot \left( \sum_{i=1}^n b_i^2 \right)- \left( \sum_{i=1}^na_ib_i \right)^2 &= \sum_{i=1}^n\sum_{j=1}^n \underbrace{a_i^2b_j^2-a_ib_ia_jb_j}_{=0\text{ si }i=j} \\
                                                                                                                      &= \sum_{i<j}(a_i^2b_j^2+a_j^2b_i^2)-2a_ib_ia_jb_j \\ 
                                                                                                                      &= \sum_{i<j}(a_ib_j-a_jb_i)^2\geq 0
    \end{align*}

\textbf{Utilisation de l'homogénéité.} On suppose les $b_i$ non tous nuls. On pose \[
    f:(a, b)\longrightarrow \left( \sum_{i=1}^na_i^2 \right) \cdot \left( \sum_{i=1}^n b_i^2 \right)- \left( \sum_{i=1}^na_ib_i \right)^2
\]
On a, pour $\lambda, \mu\geq 0$, $f(\lambda a, \mu b)=\lambda^2\mu^2f(a, b)$ de même signe que $f(a, b)$. Il existe $\lambda, \mu\geq 0$ tels que \[
    \sum_{i=1}^n \lambda^2a_i^2=\sum_{j=1}^nb_j^2\mu^2=1
\]
de sorte qu'on peut supposer sans perte de généralité \[
    \sum_{i=1}^na_i^2=\sum_{j=1}^n b_j=1
\]
Puis $|a_ib_i|\leq \frac{a_i^2+b_i^2}2$ donc \[
    \sum_{i=1}^n|a_ib_i|\leq \frac{1+1}2=1
\]
d'où \conc

\subsection{Inégalités fonctionnelles}

Montrer que si $0<p<1$ et $\theta\in \left[ 0; \frac\pi2 \right]$ alors \[
    \cos^p\theta\leq \cos(p\theta)
\]
On pose \[
    f_p:\theta\in \left[ 0;\frac\pi2 \right]\longmapsto -\cos^p\theta+\cos(p\theta)
\]
de sorte que \[
    f_p'(\theta)=-p\sin(p\theta)+p\sin\theta\cos^{p-1}\theta
\]
or $\sin(p\theta)\leq \sin\theta\leq \sin\theta\cos^{p-1}\theta$ donc $f_p'(\theta)\geq 0$ et $f_p$ croissante, ce qui conclut puisque $f_p(0)=0$

\subsection{Inégalité de réordonnement}

On note \[
    \begin{cases}
        a_1\leq \cdots \leq a_n\\ b_1\leq \cdots \leq b_n
    \end{cases}
\]
et $\sigma\in\mathfrak S_n$. Puis, on définit \[
    S_\sigma\defeq\sum_{i=1}^na_ib_{\sigma(i)}.
\]
$S_\sigma$ est maximal pour $\sigma=\id$, et minimal pour $\sigma=(n\;\;n-1\;\;\cdots \;\; 1)$. On suppose que $\sigma$ maximise $S_\sigma$. On note $\sigma'=(1\;\;\sigma(1))\circ \sigma$ donc \[ S_{\sigma'}-S_\sigma=a_1b_1+a_{\sigma^{-1}(1)}b_{\sigma(1)}-a_1b_{\sigma(1)}-a_{\sigma^{-1}(1)}b_1 =(a_1-a_{\sigma^{-1}(1)})(b_1-b_{\sigma(1)})\geq 0\]
donc on peut supposer $\sigma(1)=1$. En itérant, $\sigma=\id$ est maximisante. On procède de la même manière pour la permutation minimisante.

\begin{ex}[Application 1]
    On note $f:\mathbb N^\star\to\mathbb N^\star$ bijective. On cherche la nature de $\sum\dfrac1{nf(n)}$.
    On note $a_1<\cdots <a_n$ les éléments $f(1), \cdots, f(n)$ ordonnés. Alors \[
        \sum_{k=1}^n\frac1{kf(k)}\leq \frac1{na_n}+\cdots \frac1{a_1}\leq \frac1{n^2}+\cdots +\frac1{1^2}\leq\zeta(2)
    \]
    donc la série converge.
\end{ex}

\begin{ex}[Inégalité AM-GM]
    On note $x_1, \cdots, x_n\geq 0$, et on note \[
        f(x_1, \cdots, x_n)=\frac1n\sum x_i-\sqrt[n]{\prod x_i}
    \]
    de sorte que $f(\lambda x_1, \cdots, \lambda x_n)=\lambda f(x_1, \cdots, x_n)$ pour $\lambda\geq 0$ donc on peut supposer $\prod x_i=1$. En notant \[
        a_k=\prod_{i=1}^kx_i\qquad\qquad\qquad b_i=\frac1{a_i}
    \]
    on a \[
        n=a_1b_1+\cdots a_nb_n\leq \underbrace{a_1b_n}_{x_1}+\underbrace{a_2b_1}_{x_2}+\cdots + a_nb_{n-1}=x_1+\cdots+x_n
    \]
    ce qui conclut.
\end{ex}

\subsection{Échelles de références}

\begin{dfn}[Rappels]
    Si $f$ et $g$ sont définies sur un voisinage de $a\in\bar{\mathbb R}$ alors \[
        f(x)=o_a(g(x))\iff \exists V'\subset V, V'\in\mathcal V(a), \exists \varepsilon:V'\to\mathbb R, \forall x\in V', f(x)=\varepsilon(x)g(x)\text{ et }\varepsilon(x)\xrightarrow[x\to a]{}0
    \]
    \[
        f(x)=O_a(g(x))\iff \exists M>0, \exists V'\subset V, V\in\mathcal V(a), \forall x\in V', |f(x)|\leq M|g(x)|
    \]
    \[
        f(x)\underset a\sim g(x)\iff f(x)=g(x)+o_a(g(x))
    \]
\end{dfn}

\begin{ex}
    \[ x^\alpha|\ln x|^\beta=o_0(x^{\alpha'}|\ln x|^{\beta'})\iff \alpha>\alpha'\quad \text{ ou }\quad \alpha=\alpha'\text{ et }\beta<\beta'\]
    \[x^\alpha (\ln x)^\beta=o_{+\infty}(x^{\alpha'}(\ln x)^{\beta'})\iff \alpha<\alpha'\quad\text{ ou }\quad \alpha=\alpha'\text{ et }\beta<\beta'\underset{\text{def}}\iff (\alpha, \beta)\prec (\alpha', \beta')\]
\end{ex}

\begin{csq}
    $(x\in\mathbb R_+^\star\longmapsto x^\alpha |\ln x|^\beta)_{\alpha,\beta}$ est une famille libre (pour une famille finie, on ordonne \\$(\alpha_1, \beta_1)\prec~\cdots~\prec~(\alpha_n,\beta_n)$ et on prend la limite en $+\infty$)
\end{csq}

\section{Suites numériques}

\subsection{Caractérisation de la borne supérieure}

\begin{thm}
    \Hyp $\emptyset\neq A\subset\mathbb R$, $M\in\mathbb R$.
    \Conc Il y a équivalence entre \begin{enumerate}[label=(\alph{enumi}),left=1.2cm]
        \item $M=\sup A$
        \item $M$ majore $A$ et $\exists a=(a_n)_n\in A^{\mathbb N}, \quad a_n \longrightarrow M$
    \end{enumerate}
\end{thm}

\begin{proof}
    ~ \begin{enumerate}
        \item $(a\implies b)$ $M$ majore $A$ par définition et pour tout $n$, il existe $a_n\in A$ tel que \[
                M-\frac1n<a_n
            \]
            donc $a_n\longrightarrow M$
        \item $(b\implies a)$ Supposons qu'il existe un majorant $M'<M$ de $A$. Alors toute suite convergente de $A$ a une limite $\ell\leq M'<M$ donc $a_n\longrightarrow M\implies M\leq M'$ absurde. Donc $M$ est le plus petit majorant.
    \end{enumerate}
\end{proof}

\subsection{Limite supérieure et inférieure}

On se donne $u$ une suite réelle bornée. On pose pour $n\in\mathbb N$, \[
    u_n^+=\sup\{u_k, \quad k\geq n\}\qquad \qquad u_n^-=\inf \{u_k,\quad k\geq n\}
\]
de sorte que $u_0^-\leq u_n^-\leq u_n\leq u_n^+\leq u_0^+$ et les deux suites $u^+, u^-$ étant monotones bornées, elles sont convergentes.

\begin{dfn}
    Pour une suite réelle $u$, on appelle limite supérieure et limite inférieure les quantités \[
        \limsup_{n\to+\infty}u_n\defeq\lim_{n\to+\infty}u_n^+\qquad\qquad\qquad\qquad \liminf_{n\to+\infty}u_n\defeq\lim_{n\to+\infty}u_n^-
    \]
\end{dfn}

On va construire une extractice $\varphi$ de la manière suivante: \begin{itemize}
    \item $\varphi(0)=0$. On a $u_1^+-1<u_1+$ donc $\exists k_1\geq 1, u_1^+-1<u_{k_2}\leq u_1^+$
    \item $\varphi(1)=k_1$. On a $u_2^+-\frac12<u_2+$ donc $\exists k_2\geq k_1, u_2^+-\frac12<u_{k_2}\leq u_2^+$
    \item $\cdots$
\end{itemize}
On construit ainsi une sous-suite convergente (qui converge vers $\limsup u$).

\begin{rem}
    On a montré le théorème de Bolzano-Weierstrass\index{Bolzano-Weierstrass (théorème de -- )}
\end{rem}


\chapter{Suites et séries numériques}

\ifsolo
    ~

    \vspace{1cm}

    \begin{center}
        \textbf{\LARGE Suites et séries numériques} \\[1em]
    \end{center}
    \tableofcontents
\else
    \chapter{Suites et séries numériques}

    \minitoc
\fi
\thispagestyle{empty}

\ifsolo \newpage \setcounter{page}{1} \fi

\section{Rappels sur les séries numériques}

\begin{dfn}
    Soit $u\in\R^{\N}$. On appelle \textbf{série de terme général} \index{série numérique}$u_n$ la suite $(S_n(u))_n$ définie par \[
        S_n(u)\defeq \sum_{k=1}^nu_k
    \]
    On dira que cette série converge si $(S_n(u))$ converge, et on écrira \[
        \lim_{n\to+\infty}S_n(u)=\sum_{k=1}^{+\infty}u_k
    \]
    Dans ce cas, on pose \[
        R_n(u)\defeq\sum_{k=1}^{+\infty}u_k-S_n(u)
    \]
    qu'on appelle reste d'ordre $n$ de la série. La série de terme général $u_k$ sera notée $\sum u_k$
\end{dfn}


\begin{rem}
    Le terme général d'une séries convergente tend vers $0$. Si $u_n\longrightarrow \ell\in\bar{ \R}\setminus\{0\}$ alors la série diverge grossièrement. Si $\sum u_n$ converge alors $R_n(u)\longrightarrow0$ et on peut noter \[
        R_n(u)=\sum_{k=n+1}^{+\infty}u_k
    \]
\end{rem}

\begin{rem}
    On peut remplacer $\R$ par $\C$ ou n'importe quel espace vectoriel normé (traité dans un chapitre ultérieur)
\end{rem}

\begin{ex}
    Si $u_n=a^n$ pour un complexe $a$, alors \[
        S_n(u)= \begin{cases}
            n+1 &\text{ si }a=1\\[1em]
            \dfrac{1-a^{n+1}}{1-a}&\text{ sinon}
        \end{cases}
    \]
    d'où \[
        \sum u_n\text{ CV }\iff |a|<1
    \]
    et \[
        \sum_{k=0}^{+\infty}u_k=\frac1{1-a}
    \]
\end{ex}

\begin{ex}
    On a \[
        \ln \left( 1+\frac2{n(n+3)} \right)=\ln(n+1)+\ln(n+2)-\ln(n)-\ln(n+3)
    \]
    donc \[
        \sum_{k=1}^n\ln \left( 1+\frac2{k(k+3)} \right) =\ln3+\ln \left( \frac{n+1}{n+3} \right)\xrightarrow[n\to+\infty]{}\ln 3
    \]
\end{ex}

\section{Opérations}

\begin{prop}
    \Hyp $\sum u_n, \sum v_n$ deux séries numériques, $\lambda\in\C$
    \begin{concenum}
    \item $\lambda\sum u_n+\sum v_n=\sum(\lambda u_n+v_n)$
    \item Si les séries convergent, alors $\sum(\lambda u_n+v_n)$ converge (+ égalité des limites)
    \end{concenum}
\end{prop}

\begin{rem}
    On en déduit
    \begin{center}
    \begin{tabular}{ccccc}
        CV & + & CV & = & CV\\
        DV & + & CV & = & DV\\
        DV & + & DV & = & ?
    \end{tabular}
    \end{center}
\end{rem}

\begin{prop}
    \Hyp $u\in\C^{\N}$
    \Conc \[
        (u_n)\text{ CV }\iff \sum (u_{n+1}-u_n)\text{ CV }
    \]
\end{prop}

\begin{ex}[Séries de Riemann\index{série de Riemann}]
    On veut étudier $\sum\frac1{n^\alpha}$.
    \begin{itemize}
        \item Si $\alpha\leq0$ alors la série diverge grossièrement
        \item Si $\alpha> 0$ alors \[
                \frac1{(n+1)^\alpha}-\frac1{n^\alpha}=\frac1{n^\alpha} \left( -\frac\alpha n+O \left( \frac1n \right) \right)\sim -\frac\alpha{n^\alpha}
            \]
            or $(\frac1{n^{\alpha}})_n$ converge vers $0$ donc $\sum \frac1{(n+1)^\alpha}-\frac1{n^\alpha}$ aussi et donc, par équivalence\footnote{c'est une propriété qui sera vue dans la prochaine partie}, $\sum\frac1{n^{\alpha+1}}$ aussi.

        \item Pour $\alpha<0$, on trouve alors que $\sum\frac1{n^{\alpha+1}}$ diverge
    \end{itemize}
    Il ne reste que le cas $\alpha=1$ à traiter: la série diverge. On a donc \[
    \sum\frac1{n^\alpha}\text{ CV }\iff \alpha>1
\]
\end{ex}

On va maintenant s'intéresser à la série harmonique\index{série harmonique} \[
    H_n=\sum_{k=1}^n\frac1k.
\]
On a \[
    \ln(n+1)-\ln (n)\sim \frac1n \implies \sum\frac1k \text{ DV }
\]
et par équivalence des sommes partielles\footnote{vu plus tard} (ou bien par comparaison série-intégrale), \[
    H_n\sim\ln(n).
\]
On note $u_n=H_n-\ln n$. On a \[
    u_{n+1}-u_n=\frac1{n+1}-\ln \left( 1+\frac1n \right)=\frac1n\times \frac1{1+\frac1n}-\frac1n-\frac1{2n^2}+o \left( \frac1{n^2} \right)=-\frac2{n^2}+o \left( \frac1{n^2} \right)
\]
donc $\sum u_{n+1}-u_n$ converge donc $(u_n)$ converge vers un réel $\gamma$. On note \[
    v_n=H_n-\ln n-\gamma
\]
et \[
    v_{n+1}-v_n=-\frac1{2n^2}+o \left( \frac1{n^2} \right)
\]
donc $(v_n)$ converge. Par équivalence des restes (série CV), \[
    \sum_{k=n}^{+\infty}v_{k+1}-v_k\sim\sum_{k=n} -\frac12 \left( \frac1n-\frac1{n+1} \right)
\]
d'où $v_n\sim \frac1{2n}$.
On pose \[
    \omega_n=H_n-\ln n-\gamma-\frac1{2n}
\]
Après calculs (vérifiez-le), \[
    \omega_{n+1}-\omega_n=\frac1{6n^3}+o \left( \frac1{n^3} \right)\sim\frac1{6n^3}\sim\frac1{12}\cdot \left( \frac 1{n^2}-\frac1{(n+1)^2} \right)
\]
d'où $\omega_n\sim\frac1{12n^2}$

\textbf{Bilan.} (Question X) \[
    H_n=\ln n+\gamma+\frac1{2n}-\frac1{12n^2}+o \left( \frac1{n^2} \right)
\]

\begin{exo}[Mines-Ponts]
    On note \[
        k_n=\min\{k\in\N, \quad H_k>n\}
    \]
    Calculer \[
        \lim_{n\to+\infty}\frac{k_{n+1}}{k_n}
    \]
\end{exo}

\section{Séries à termes positifs}

\begin{prop}
    \Hyp $u$ est une suite positive
    \Conc $\sum u_n$ converge si et seulement si $(S_n(u))$ est majorée
\end{prop}

\begin{prop}
    \Hyp $u, v$ des suites positives
    \begin{concenum}
    \item Si $u_n\leq v_n$ à partir d'un certain rang et si $\sum v_n$ converge alors $\sum u_n$ aussi
    \item Si $u_n=O(v_n)$ et $\sum v_n$ converge alors $\sum u_n$ aussi
    \item Si $u_n\sim v_n$ alors $\sum u_n$ et $\sum v_n$ ont la même nature
    \end{concenum}
\end{prop}

\begin{proof}~
    \begin{enumerate}
        \item Si $u_n\leq v_n$ à partir du rang $N$, \[
                S_n(u)\leq \sum_{k=0}^Nu_k+\sum_{k=N+1}^{+\infty}v_k
            \]
            donc $(S_n(u))$ croissante est majorée donc converge.
        \item Idem à constante multiplicative près
        \item $u_n\sim v_n$ donne $u_n=O(v_n)$ et $v_n=O(u_n)$ et on applique $2$.
    \end{enumerate}
\end{proof}

\begin{exo}
    On note $(a_n)_{n\geq 1}$ une suite positive telle que $\sum a_n$ converge. Donner la nature de \[
        \sum a_n\sin(a_n), \qquad \sum\frac{a_n}{1+a_n^2}, \qquad \sum\frac{\sqrt{a_n}}{n}
    \]
\end{exo}

\begin{proof}[Résolution]
    À partir d'un certain rang, $0\leq a_n\sin(a_n)\leq a_n$ puis \[
        \frac{a_n}{1+a_n^2}\sim a_n
    \]
    et enfin (AM-GM) \[
        0\leq \frac{\sqrt{a_n}}n\leq\frac12 \left( a_n+\frac1{n^2} \right)
    \]
    donc les trois séries convergent.
\end{proof}

\begin{exo}[X]
    Soit $x$ une suite positive. Montrer que si $\sum x_n$ converge, alors \[
        \sum x_n^{\frac n{n+1}}
    \]
    aussi
\end{exo}

\begin{proof}[Résolution]~
    Idée: On fait un découpage. On pose \[
        I=\{n\in\N, \quad x_n\geq 0 \text{ et }x_n^{-\frac1{n+1}}>2\}
    \]
    Si $n\in I$ alors \[
        \frac12>x_n^{\frac1{n+1}}\iff \frac1{2^{n+1}}>x_n\implies \frac1{2^n}>y_n
    \]
    Sinon, \[
        x_n=0\implies y_n=0\leq 2x_n
    \]
    et \[
        x_n^{-\frac1{n+1}}\leq 2\implies y_n\leq 2x_n
    \]
    donc dans tous les cas, \[
        0\leq y_n\leq \frac1{2^n}+2x_n
    \]
\end{proof}

\section{Absolue convergence, semi-convergence}

\begin{dfn}
    Soit $a$ une suite réelle. \begin{enumerate}
        \item On dira que $\sum a_n$ est absolument convergnete\index{absolue convergence} si $\sum |a_n|$ converge
        \item On dira que $\sum a_n$ est \index{semi-convergence}semi convergente si elle est convergente mais pas absolument convergente
    \end{enumerate}
\end{dfn}

\begin{thm}
    L'absolue convergence entraîne la convergence
\end{thm}

\begin{proof}
    On écrit $a_n=a_n^+-a_n^-$ avec $a_n^+=\max(a_n, 0)$ et $a_n^-=\max(-a_n, 0)$ de sorte que $|a_n|=a_n^++a_n^-$
\end{proof}

\begin{rem}
    \begin{center}
    \begin{tabular}{ccccc}
        ACV & + & ACV & = & ACV\\
        ACV & + & SCV & = & SCV\\
        SCV & + & SCV & = & ?
    \end{tabular}
    \end{center}
\end{rem}

\section{Comparaison logarithmique}

\begin{prop}
    \Hyp $u, v$ deux suites strictement positives telles que APCR $N$ \[
    \frac{u_{n+1}}{u_n}\leq \frac{v_{n+1}}{v_n}
    \]
    \Conc $u_n=O(v_n)$
\end{prop}

\begin{proof}
    \[
        \forall n>N, \prod_{k=N}^{n-1}\frac{u_{k+1}}{u_k}=\frac{u_n}{u_N}\leq \frac{v_n}{v_N}\implies u_n\leq v_n\frac{u_N}{v_N}=O(v_n)
    \]
\end{proof}

\begin{thm}[Règle de D'Alembert\index{D'Alembert (règle de -- )}]
    \Hyp $a$ une suite positive telle que \[
        \frac{a_{n+1}}{a_n}\xrightarrow[n\to+\infty]{}\ell\in\bar{\R}
    \]
    \begin{concenum}
    \item Si $0\leq \ell <1$ alors $\sum a_n$ converge
    \item Si $\ell>1$ alors $\sum a_n$ diverge
    \end{concenum}
\end{thm}

\begin{proof}~
    \begin{itemize}
        \item $\exists \epsilon>0, \quad 0<\ell'=\ell+\epsilon<1$ et on applique la proposition précédente à $u=a$, $v=(\ell'^n)_n$.
        \item $\exists \epsilon > 0, \quad 1<\ell'=\ell-\epsilon$ et on fait pareil.
    \end{itemize}
\end{proof}

\begin{rem}
    Il faut faire attention à avoir $(a_n)$ jamais nulle à partir d'un certain rang. Le critère est intéressant lorsque $\frac{a_{n+1}}{a_n}$ a une expression simple.
\end{rem}

\begin{exo}
    Donner la nature de $\sum\frac{n^n}{n!^\alpha}$ en fonction de $\alpha$
\end{exo}

\begin{proof}[Résolution]
    On a \[
        \frac{a_{n+1}}{a_n}=\frac{(n+1)^{n+1}}{n^n}\times \frac1{(n+1)^\alpha}=\frac{ \left( 1+\frac1n \right)^n }{(n+1)^{\alpha-1}}\xrightarrow[n\to+\infty]{} \begin{cases}
            0&\text{ si }\alpha>1 \\ e &\text{ si }\alpha=1\\ +\infty& \text{ si }\alpha<1
        \end{cases}
    \]
    et on conclut avec la règle de d'Alembert
\end{proof}

\begin{exo}
    On note $a$ une suite réelle telle que \[
        |a_n|^{\frac1n}\xrightarrow[n\to+\infty]{}\ell\in\R\cup\{+\infty\}
    \]
    Donner la nature de $\sum |a_n|$ en fonction de $\ell$
\end{exo}

\begin{proof}[Résolution]
    Si $\ell>1$ alors $\exists \epsilon>0, 1<\ell'=\ell-\epsilon\leq |a_n|^{\frac1n}$ donc $\ell'^n\leq |a_n|$ et la série diverge.

    Si $\ell<1$ alors par un raisonnement similaire, la série converge.

    Si $\ell=1$ alors on ne peut pas conclure en général: \[
        \left| \frac1{n^2} \right|^{\frac1n}\xrightarrow[n\to+\infty]{}1\qquad\text{ et }\qquad \sum\frac1{n^2}\text{ CV }
    \]
    \[
        \left| \left( 1+\frac1n\right)^n \right|^{\frac1n}\xrightarrow[n\to+\infty]{}1\qquad \text{ et }\qquad \sum \left( 1+\frac1n \right)^{n}\text{ DV grossièrement }
    \]
\end{proof}

\section{Sommation des relations de comparaion}

\begin{thm}
    \Hyp $u, v$ sont des suites réelles positives et $\sum v_n$ converge.
    \begin{concenum}
    \item Si $u_n=O(v_n)$ alors $\sum u_n$ converge et \[
            R_n(u)=O(R_n(v))
        \]
    \item Si $u_n=o(v_n)$ alors $\sum u_n$ converge et \[
            R_n(u)=o(R_n(v))
        \]
    \item Si $u_n\sim v_n$ alors $\sum u_n$ converge et \[
            R_n(u)\sim R_n(v)
        \]
    \end{concenum}
\end{thm}

\begin{proof}
    La convergence est déjà vue, on ne montre que les relations de comparaison

    \begin{enumerate}
        \item Il existe un rang $N$ et une constante $C>0$ tels que $\forall n\geq N, u_n\leq Cv_n$. On a alors \[
    \forall m\geq n\geq N, \qquad \sum_{k=n}^Nu_k\leq C\sum_{k=n}^mv_k
            \]
            puis $m\to+\infty$ donnc $R_n(u)\leq CR_n(v)$ ie \conc
        \item Soit $\epsilon>0$. Il existe un rang $N$ après lequel $u_n\leq \epsilon v_n$. Par le même raisonnement, $R_n(u)\leq \epsilon R_n(v)$ après $N$ donc \conc
        \item $|u_n-v_n|=o(v_n)$ donc \[
                |R_n(u)-R_n(v)|=|R_n(u-v)|\leq R_n(|u-v|)=o(R_n(v))
            \]
    \end{enumerate}
\end{proof}

\begin{thm}
    \Hyp $u, v$ suites réelles positives et $\sum u_n$ diverge.
    \item Si $u_n=O(v_n)$ alors $\sum v_n$ diverge et \[
            S_n(u)=O(S_n(v))
        \]
    \item Si $u_n=o(v_n)$ alors $\sum v_n$ diverge et \[
            S_n(u)=o(S_n(v))
        \]
    \item Si $u_n\sim v_n$ alors $\sum v_n$ diverge et \[
            S_n(u)\sim S_n(v)
        \]
\end{thm}

\begin{proof}
    La divergence est déjà vue, on ne montre que le reste.
    \begin{enumerate}
        \item Il existe $C>0$, et un rang $N$ à partir duquel $u_n\leq Cv_n$. \[
                \forall n\geq N, \quad \underbrace{S_n(u)}_{\longrightarrow+\infty}=\sum_{k=0}^{N-1}u_k+\sum_{k=N}^nu_k\leq \sum_{k=0}^{N-1}+C\underbrace{\sum_{k=0}^nv_k}_{S_n(v)\longrightarrow+\infty}-C\sum_{k=0}^{N-1}v_k
            \]
            donc à partir d'un certain rang, \[
                \frac{S_n(u)}{S_n(v)}\leq C+\frac{\sum_{k=0}^{N-1}u_k-C\sum_{k=0}^{N-1}v_k}{S_n(v)}\xrightarrow[n\to+\infty]{}C
            \]
            d'où \conc
        \item Idem
        \item $S_n(|u-v|)=o(S_n(v))$ (attention aux hypothèses, si $\sum |u-v|$ converge ok sinon on applique 2.)
    \end{enumerate}
\end{proof}

\begin{ex}
    On va montrer \[
        \sum_{k=n}^{+\infty}\frac{\ln(k+1)}{k^2}\sim\frac{\ln n}n
    \]
\end{ex}

\begin{proof}
    La série converge puisque $\frac{\ln n}{n^2}=o \left( \frac1{n^{3/2}} \right)$. Puis,
    on a \[
        \frac{\ln(n+1)}{n+1}-\frac{\ln n}n=\frac{\ln \left( \frac{n+1}n \right)}{n+1}+\frac{\ln n}{n+1}-\frac{\ln n}n=\ln n \left( \frac1{n+1}-\frac1n \right)+\ln \left( 1+\frac1n \right)\frac1{n+1}\sim -\frac{\ln n}{n^2}.
    \]
    d'où la conclusion.
\end{proof}

\begin{exo}
    Donner un développement asymptotique de \[
        R_n=\sum_{k=n}^{+\infty}\frac1{k^2}
    \]
\end{exo}

\begin{proof}[Résolution]
    \[
        \frac1{n^2}\sim\frac1n-\frac1{n+1}
    \]
    donc \[
        R_n\sim \frac1n
    \]
    puis \[
        R_n-\frac1n=\sum_{k=n}^{+\infty} \left( \frac1{k^2}-\frac1k+\frac1{k+1} \right)=\sum_{k=n}^{+\infty}\frac1{k^2(k+1)}
    \]
    or \[
        \frac1{n^3}\sim \frac1{2n^2}-\frac1{2(n+1)^2}
    \]
    donc \[
        R_n=\frac1n+\frac1{2n^2}+o \left( \frac1{n^2} \right)
    \]
\end{proof}

\begin{ex}[Équivalent de Stirling\index{Stirling (formule de -- )}]
    On va trouver un équivalent de $u_n=n!$. On a \[
        \ln(u_n)=\sum_{k=1}^n\ln k
    \]
    or \[
        (k+1)\ln(k+1)-k\ln k\sim \ln k
    \]
    qui est le terme général positif d'une série divergente donc \[
        \ln (u_n)\sim\sum_{k=1}^n((k+1)\ln(k+1)-k\ln k)\sim n\ln n.
    \]
    On note $v_n=\ln u_n-n\ln n$ puis \[
        v_{n+1}-v_n=\ln(n+1)-(n+1)\ln(n+1)+n\ln n=-n\ln \left( 1+\frac1n \right)\sim -1
    \]
    qui est le terme général négatif (mais de signe constant) d'une série divergente donc \[
        v_n\sim -n
    \]
    On note $w_n=v_n+n$ et \[
        w_{n+1}-w_n=-n \left( 1+\frac1n \right)+1=-n \left( \frac1n-\frac1{2n^2}+o \left( \frac1{n^2} \right) \right)+1=\frac1{2n}+o \left( \frac1n \right)\sim\frac1{2n}\sim \frac12 (\ln(n+1)-\ln n)
    \]
    donc \[
        w_n\sim \frac{\ln n}2
    \]
    On pose $y_n=w_n-\frac12\ln n$ et \[
        y_{n+1}-y_n=-n\ln \left( 1+\frac1n \right)+1-\frac12\ln \left( 1+\frac1n \right)\sim\frac1{6n^2}
    \]
    qui est le terme général d'une séries ACV donc $y_n$ converge vers un réel $\ell$.
    On a donc \[
        \ln(u_n)=n\ln n-n+\frac12\ln n+\ell+o(1)
    \]
    donc \[
        u_n\sim e^\ell\sqrt n\; n^ne^{-n}
    \]
    Dans le chapitre d'intégration \ifsolo\else(\textbf{\nameref{sec:wallis}} p.\pageref{sec:wallis})\fi{} on montrera $e^\ell=\sqrt{2\pi}$
\end{ex}

\begin{ex}[Raabe-Duhamel\index{Raabe-Duhamel (règle de -- )}]
    On note $u$ une suite réelle jamais nulle à partir d'un certain rang telle que \[
        \frac{u_{n+1}}{u_n}=1-\frac\lambda n+o \left( \frac1n \right), \qquad \lambda>0
    \]
    On va chercher la nature de $\sum u_n$ en fonction de $\lambda$

    On a \[
        \frac{u_{n+1}}{u_n}\sim 1
    \]
    donc $u_n$ est de signe constant à partir d'un certain rang, on suppose spdg que c'est positif. \[
        \ln u_{n+1}-\ln u_n\sim-\frac\lambda n
    \]
    qui est le terme général d'une série divergente (signe constant) donc \[
        \ln u_n\sim -\lambda H_n\sim -\lambda\ln n
    \]
    d'où $\ln u_n=-\lambda \ln n + o(\ln n)$ donc $u_n=n^{-\lambda+o(1)}=\dfrac1{n^{\lambda+o(1)}}$.

    Si $\lambda>1$ alors $\sum u_n$ converge, si $\lambda <1$ alors $\sum u_n$ diverge et si $\lambda=1$ on ne peut pas conclure (e.g. $u_n=1/n$ et $u_n=1/(n\ln^2n)$).
\end{ex}

\begin{rem}
    Si $\ell\neq 0$, $u_n\sim \ell$ (tg série DV) donc $\sum_{k=1}^nu_k\sim n\ell$ d'où le théorème de Cesàro. Si $\ell=0$ alors $u_n=o(1)$ et $\sum_{k=1}^nu_k=o(n)$ et on retrouve aussi Cesàro.

    Plus généralement, si \[
        \frac{a_{n+1}-a_n}{b_{n+1}-b_n}\longrightarrow\ell\neq 0
    \]
    avec $b_n$ croissante qui tend vers $+\infty$, alors \[
        a_{n+1}-a_n\sim \ell (b_{n+1}-b_n)
    \]
    qui est le terme général d'une série divergente donc \[
        \sum_{k=1}^{n-1}a_{n+1}-a_n=a_n-a_1\sim \ell (b_n-b_1) \implies a_n\sim \ell b_n\implies \frac{a_n}{b_n}\longrightarrow \ell
    \]
    d'où le théorème de Cesàro-Stolz (dans le cas $\ell\in\R^\star$)
\end{rem}

\section{Comparaison série-intégrale\texorpdfstring{\footnote{Épisode 1}}{}}

On note $f$ décroissante continue positive au voisinage de $+\infty$ (pour simplifier on prend $\R_+$ comme voisinage). On a \[
    f(k+1)\leq \int_k^{k+1}f\leq f(k)
\]

\begin{ex}
    On trouve de cette manière \[
        H_n\sim\ln n\qquad \text{ et }\qquad H_{2n}-H_n\sim \ln 2.
    \]
\end{ex}

\begin{thm}
    \Hyp $f$ est continue décroissante positive sur $\R_+$, $u_n=\displaystyle\int_{n-1}^n f-f(n)$
    \Conc $\sum u_n$ converge et $u_n\geq 0$
\end{thm}

\begin{proof}
    $u_n\leq f(n-1)-f(n)$ tg série CV car $f$ converge en $+\infty$
\end{proof}

\begin{exo}
    On note $u$ une suite positive, $S_n=\sum_{k=0}^nu_k$ et quand c'est possible, $R_n=\sum_{k=n}^{+\infty}u_k$
    \begin{enumerate}
        \item Montrer que si $\sum u_n$ diverge alors $\sum \frac{u_n}{S_n^\alpha}$ converge ssi $\alpha>1$
        \item Montrer que si $\sum u_n$ converge alors $\sum \frac{u_n}{R_n^\alpha}$ converge ssi $\alpha<1$
    \end{enumerate}
\end{exo}

\begin{proof}[Résolution (uniquement 1.)]~
    \begin{enumerate}
        \item On a $S_n\longrightarrow+\infty$, $(S_n)$ croissante. On suppose $\alpha>1$. \[
                \frac{u_n}{S_n^\alpha}=\frac{S_n-S_{n-1}}{S_n^\alpha}\leq \int_{S_{n-1}}^{S_n}\frac{\diff t}{t^\alpha}\implies \sum_{k=1}^n\frac{u_k}{S_k^\alpha}\leq \frac1{\alpha-1} \left( \frac1{S_0^{\alpha-1}}-\frac1{S_n^{\alpha-1}} \right)\leq \frac1{\alpha-1}S_0^{\1-\alpha} \text{ donc CV }
            \]
            Si $\alpha\geq 0$ alors APCR, $\frac{u_n}{S_n^\alpha}\geq u_n$ donc la série diverge.

            Si $\alpha=1$, on a \[
                \sum_{k=n}^p\frac{u_k}{S_k}\geq \frac1{S_p}\sum_{k=n}^pu_k=\frac{S_p-S_{n-1}}{S_p}\geq 1-\frac{S_{n-1}}{S_p}
            \]
            Si $\sum\frac{u_k}{S_k}$ converge alors $p\longrightarrow+\infty$ donne $\sum_{k=n}^{+\infty}\frac{u_k}{S_k}\geq 1$ absurde, donc la série diverge.
            Pour $\alpha\in ]0, 1]$, \[
                \frac{u_k}{S_k^\alpha}\underset{\text{APCR}}\geq \frac{u_k}{S_k}
            \]
            donc la série diverge.
    \end{enumerate}
\end{proof}

\section{Théorème des séries alternées}

\begin{thm}
    \Hyp $a$ une suite décroissante positive qui tend vers $0$.
    \begin{concenum}
    \item $\sum a_n(-1)^n$ converge
    \item \[
            \sum_{k=n+1}^{+\infty}(-1)^ka_k
        \]
        a le même signe que $(-1)^{n+1}a_{n+1}$ (le premier terme de la somme)
    \item \[
            \left|  \sum_{k=n+1}^{+\infty}(-1)^ka_k \right|\leq a_{n+1}
    \]
    \end{concenum}
\end{thm}

\begin{proof}~
    \begin{enumerate}
        \item Les suites $S_{2n}$ et $S_{2n+1}$ sont adjacentes donc convergentes de limite $\ell$.
        \item Vu le calcul fait en 1., \[
                \sum_{k=2p}^{+\infty}(-1)^ka_k=\ell-S_{2p-1}\geq 0
            \]
            du même signe que son premier terme. C'est la même chose pour les restes d'indices impair.
        \item \[ 0\leq \ell-S_{2p-1}\leq S_{2p}-S_{2p-1}\implies 0\leq \sum_{k=2n}(-1)^ka_k\leq a_{2n}\]
            et on fait dans l'autre sens pour les impairs.
    \end{enumerate}
\end{proof}

\section{Transformation d'Abel}

On part de deux suites $(a_n)$ et $(b_n)$ et d'une somme $\displaystyle\sum_{k=0}^na_kb_k$. On pose \[
    S_n=\sum_{k=0}^na_k
\]
On a $a_n=S_n-S_{n-1}$ pour $n\geq 0$ avec $S_{-1}=0$. Alors, \[
    \sum_{k=0}^na_kb_k=\sum_{k=0}^n(S_n-S_{n-1})b_k=\sum_{k=0}^{n-1}S_k(b_k-b_{k+1})+S_n b_n
\]

\begin{ex}
    On suppose $|S_n|$ majorée par $M$ et $(b_n)$ décroissante de limite nulle. \[
        b_nS_n\longrightarrow0
    \]
    et \[
        |S_k(b_k-b_{k+1})|\leq M(b_k-b_{k+1})
    \]
    donc \[
        \sum S_k(b_k-b_{k+1})\text{ ACV }
    \]
    donc $\sum a_kb_k$ converge.
\end{ex}

\begin{ex}
    $\sum\frac{\sin n}n$ converge car $S_n(\sin n)$ est bornée (on peut faire le calcul) et $\frac1n$ est décroissante de limite nulle.
\end{ex}

\begin{ex}
    On va chercher la nature de $\sum \frac{\sin \sqrt n}n$. On note $S_n$ les sommes partielles de sorte que \[
        \forall n\geq 1, \qquad \sin \sqrt n=S_{n}-S_{n-1}
    \]
    et \[
        \sum_{k=1}^n\frac{\sin \sqrt n}{n}=\sum_{k=1}^n\frac{S_k-S_{k-1}}k=\sum_{k=1}^{n-1}S_k \left( \frac1k-\frac1{k+1} \right)+\frac{S_n}n.
    \]
    Pour $x\in [k, k+1]$, il existe $c_x\in[k, k+1]$ tel que \[
        \sin\sqrt x-\sin \sqrt k=(x-k)\frac{\cos \sqrt {c_x}}{2\sqrt {c_x}}
    \]
    donc \[
        |\sin \sqrt x-\sin \sqrt k|\leq \frac1{2\sqrt k}
    \]
    et \[
        \left| \sin k-\int_k^{k+1}\sin\sqrt t\;\diff t \right|\leq \int_k^{k+1}|sin\sqrt k-\sin \sqrt t|\diff t\leq \frac1{2\sqrt k}
    \]
    donc \[
        \left| \sum_{k=1}^n\left(\sin \sqrt k-\int_k^{k+1}\sin\sqrt t\;\diff t\right) \right|= \left| S_n-\int_1^{n+1}\sin\sqrt t\;\diff t \right| \leq \sum_{k=1}^n\frac1{2\sqrt k}
    \]
    or \[
        \int_1^{n+1}\sin\sqrt t\;\diff t=\int_1^{\sqrt{n+1}}2u\sin u\diff u=2[-u\cos u+\sin u]_1^{\sqrt{n+1}}=O(\sqrt n).
    \]
    Puis, \[
        \frac1{2\sqrt k}\sim\sqrt {k+1}-\sqrt k \implies \sum_{k=1}^n\frac1{2\sqrt k}\sim \sqrt n
    \]
    On a donc \[
        |S_n|- \left| \int_1^{n+1}\sin\sqrt t\;\diff t \right|\leq \sum_{k=1}^n\frac1{2\sqrt k}=O(\sqrt n)\implies S_n=O(\sqrt n)
    \]
    donc \[
        S_k \left( \frac1k-\frac1{k+1} \right)=O\left(\frac1{n^{3/2}}\right)
    \]
    et la série converge.
\end{ex}

\section{Utilisation des paquets}

\begin{rem}
    Si $\sum a_n$ converge alors pour toute suite d'entiers $(u_n)$ telle que $u_n\geq n$, on a \[
        \sum_{k=n}^{u_n}a_k=S_{u_n}(a)-S_{n-1}(a)\xrightarrow[n\to+\infty]{}0
    \]
    En particulier,  si \[
        \sum_{k=n}^{u_n}a_k\xnrightarrow[n\to+\infty]{}0
    \]
    alors la série diverge.
\end{rem}

\begin{ex}[Série harmonique]
    \[
        \sum_{k=n}^{2n}\frac1k\geq \frac{n+1}{2n+1}\xrightarrow[n\to+\infty]{}\frac12\neq 0
    \]
    donc la série harmonique diverge.
\end{ex}

\begin{ex}[Oscillations lentes]
    On va déterminer la nature de \[
        \sum\frac{\cos (\ln n)}n
    \]
    Pour $N\geq 1$, \[
        \sum_{-\frac\pi3+2N\pi\leq \ln n\leq \frac\pi3+2N\pi}\frac{\cos(\ln n)}n \geq \frac12\sum_{e^{-\frac\pi3}e^{2N\pi}\leq n\leq e^{\frac\pi3}e^{2N\pi}}\frac1n\geq \frac12\frac{\floor{e^{\frac\pi3}e^{2N\pi}}-\floor{e^{-\frac\pi3}e^{2N\pi}}}{e^{\frac\pi3}e^{2N\pi}}\xrightarrow[n\to+\infty]{}\frac12(1-e^{-\frac{2\pi}3})\neq 0
    \]
    donc la série diverge.
\end{ex}

\begin{ex}[Série lacunaire]
    On note \[
        u_n=\begin{cases}\frac1{n^\alpha}&\text{ si $n$ n'a pas de $9$ dans son écriture décimale }\\ 0 & \text{ sinon }\end{cases}
    \]
    Les sommes partielles sont croissantes donc la série $\sum u_n$ converge si et seulement si $(S_n)$ a une suite extraite convergente.
    On va considérer des paquets du type \[
        \sum_{k=10^n}^{10^{n+1}-1}u_k\leq \frac1{10^{n\alpha}}\times \underbrace{8\times 9^n}_{\text{nb d'entiers sans }9}
    \]
    On a ainsi \[
        \underbrace{\left( \frac9{10^\alpha} \right)^n}_{\text{tg série CV} \iff \frac9{10^\alpha}<1}\frac8{10^\alpha}\leq \sum_{k=10^n}^{10^{n+1}-1}u_k=S_{10^{n+1}-1}-S_{10^n-1}\leq \left( \frac9{10^\alpha} \right)^n\cdot 8
    \]
    donc $(S_{10^n-1})_n$ converge ssi $\displaystyle \alpha>\frac{\ln 9}{\ln 10}$
\end{ex}

\section{Développement décimal illimité}

Pour $x\in\R$, on pose $a_0=\floor x$ et \[
    a_n=\floor{10^nx}-10\floor{10^{n-1}x}\in\llbracket 0,9\rrbracket
\]
Par récurrence, on montre \[
    \frac{\floor{10^nx}}{10^n}=\sum_{k=0}^n\frac{a_k}{10^k}=\bar{a_0,a_1\cdots a_n}
\]
Puis, $\sum \frac{a_k}{10^k}$ converge et sa somme vaut $x$.

\begin{rem}
    La suite $(a_n)$ n'est pas stationnaire égale à $9$.
\end{rem}

\begin{proof}
    On suppose que la suite est stationnaire égale à $9$: $a_{N-1}\neq 9$ et $(a_n)_{n\geq N}$ constante égale à $9$. Dans ce cas, \[
        x=\sum_{k=0}^N\frac{a_k}{10^k}+\underbrace{\sum_{k=N}^{+\infty}\frac9{10^k}}_{=\frac1{10^{N-1}}}
    \]
    et c'est absurde en calculant explicitement $a_{N-1}$
\end{proof}

On dit que le développement obtenu est le \textbf{développement décimal propre} de $x$ \index{développement décimal propre}

\begin{prop}
    $x$ est rationnel si et seulement si son développement propre est ultimement périodique (périodique à partir d'un certain rang)
\end{prop}

\begin{proof}
    On supose $(a_k)$ $T$-périodique à partir du rang $N$. On a \[
        x=\underbrace{\sum_{k=0}^{N-1}\frac{a_k}{10^k}}_{\in\Q}+\underbrace{\sum_{k=0}^{+\infty}\frac1{10^{kT}}}_{=\frac1{1-\frac1{10^T}}\in\Q} \underbrace{\left(  \frac{a_N}{10^N}+\cdots +\frac{a_{N+T-1}}{10^{N+T-1}}\right)}_{\in \Q}\in\Q
    \]
    On suppose $x=\frac pq\in\Q$. Il y a un nombre fini de restes modulo $p$ donc il existe $n<m$ tels que $10^np\equiv 10^mp\pmod q$.

On a $10^np=q_nq+r_n$ avec $q_n\in\N$, donc $10^nx=q_n+\frac{r_n}q$. On a alors, pour $n, i\in\N^\star$, \begin{align*}
                a_{n+i} &= E(10^{n+i}x)-10E(10^{n+i-1}x) \\
                        &= E\left(10^i \left( q_n+\frac{r_n}q \right)\right)-10E \left( 10^{i-1} \left(  q_n+\frac{r_n}q \right) \right) \\
                        &= E \left( 10^i\frac {r_n}q \right)+10^iq_n-10E \left( 10^{i-1}\frac {r_n}q \right) -10^iq_n \\
                        &= E \left( 10^i\frac {r_n}q \right) - 10E \left( 10^{i-1}\frac {r_n}q \right)
            \end{align*}
         On pose $T=m-n$. On a, pour tout $i\geq 1$, $a_{n+i+T}=a_{m+i}=a_{n+i}$. Donc $(a_k)_k$ est périodique à partir du rang $n$.
\end{proof}

\endchapter


\chapter{Familles sommables}

\ifsolo
    ~

    \vspace{1cm}

    \begin{center}
        \textbf{\LARGE Familles sommables} \\[1em]
    \end{center}
    \tableofcontents
\else
    \chapter{Familles sommables}

    \minitoc
\fi
\thispagestyle{empty}

\ifsolo \newpage \setcounter{page}{1} \fi

\section{Dénombrabilité}

\begin{dfn}
    Un ensemble est dénombrable si et seulement si il est en bijection avec $\mathbb N$
\end{dfn}

\begin{ex}
    $\mathbb N, 2\mathbb N, \mathbb N^\star, \mathbb Z$
\end{ex}

\begin{prop}
    \begin{enumerate}
        \item Toute partie infinie de $\mathbb N$ est dénombrable
        \item Toute partie de $\mathbb N$ est finie ou dénombrable
    \end{enumerate}
\end{prop}

\begin{proof}
    Soit $A$ une partie de $\mathbb N$ infinie. On note $\varphi(0)=\min A$ et \[
        \varphi(n)=\min (A\setminus \llbracket 0, \varphi(n-1)\rrbracket)
    \]
    de sorte que $\varphi$ est croissante donc injective, et surjective car si $a\in A$ n'a pas d'antécédent, $I=\{k, \quad a<\varphi(k)\}$ est non vide et a un minimum $p$, et $\varphi(0)<a$ donc $p\geq 1$ d'où $\varphi(p-1)< a<\varphi(p)$ absurde.
\end{proof}

\begin{prop}
    \begin{enumerate}
        \item Si $f:A\to B$ est injective et $B$ dénombrable alors $A$ est fini ou dénombrable
        \item Si $f:A\to B$ est surjective et $A$ dénombrable alors $B$ est fini ou dénombrable.
    \end{enumerate}
\end{prop}

\begin{proof}~
    \begin{enumerate}
        \item $B$ est en bijection avec $\mathbb N$ donc on peut supposer $B=\mathbb N$. Si $A$ est fini alors ok, sinon $f(A)$ est une partie infinie de $\mathbb N$ en bijection avec $A$.
        \item On note $\psi(b)$ un élément de $f^{-1}(\{b\})$. $f\circ \psi=\id_B$ injective donc $\psi$ est injective et on utilise 1.
    \end{enumerate}
\end{proof}

\begin{prop}
    \begin{enumerate}
        \item Si $I$ est non-vide, fini ou dénombrable, et $(A_i)_{i\in I}$ est une famille d'ensembles finis ou dénombrables, alors \[
                \bigcup_{i\in I}A_i
            \]
            est fini ou dénombrable
        \item $E_1,\cdots, E_n$ des ensembles finis ou dénombrables, non vides. $E_1\times \cdots \times E_n$ est fini ou dénombrable et $E_1\times \cdots \times E_n$ est dénombrable si et seulement si au moins l'un des $E_i$ est dénombrable.
    \end{enumerate}
\end{prop}

\begin{proof}~
    \begin{enumerate}
        \item On note $\varphi:\mathbb N\to I$ surjective et pour chaque $i$ on note $f_i:\mathbb N\to A_i$ bijective de sorte que \[
        \psi: (p, q)\in\mathbb N^2\longmapsto f_{\varphi(q)}(p)\in\bigcup_{i\in I}A_i
            \]
            est surjective, et $\mathbb N^2$ est dénombrable.
        \item Chaque $E_i$ est en bijection avec une partie de $\mathbb N$, on peut supposer $E_i\subseteq \mathbb N$. On note $p_1, \cdots, p_n$ des premiers disjoints. \[
                f: E_1\times \cdots \times E_n\longrightarrow \mathbb N, \quad (a_1, \cdots, a_n)\longmapsto p_1^{a_1}\cdots p_n^{a_n}
            \]
            est injective et $\mathbb N$ dénombrable. Si tous les $E_i$ sont finis alors le produit est fini et si l'un est infini alors le produit est infini.
    \end{enumerate}
\end{proof}

\section{Exemples d'ensembles classiques}

\begin{ex}
    $\mathbb Q$ est dénombrable car $\mathbb Z\times \mathbb N^\star$ est dénombrable et $(p, q)\in\mathbb Z\times \mathbb N^\star\longmapsto \frac pq$ est surjective
\end{ex}

\begin{ex}
    L'ensemble des nombres algébriques est dénombrable: \[
        \mathcal A=\bigcup_{n\in\mathbb N^\star}\bigcup_{P\in\mathbb Z_n[X]}\mathcal Z_{\mathbb C}(P)
    \]
\end{ex}

\begin{ex}[Diagonale de Cantor]
    $\mathbb R$ \index{Cantor!argument diagonal}n'est pas dénombrable. Si $\mathbb R$ est dénombrable alors $[0, 1[$ est dénombrable, et on note $(r_n)_{n\in\mathbb N}$ une énumération des éléments de $[0, 1[$. On construit le réel $x\in [0, 1[$ de la manière suivante: si la $n$-ième décimale de $r_n$ vaut $1$ alors la $n$-ième décimale de $x$ vaut $0$, sinon elle vaut $1$. Le réel ainsi construit ne peut pas être un $r_n$ car sa $n$-ième décimale est différente de celle de $r_n$. On a donc trouvé un réel qui n'est pas dans l'énumération, c'est absurde.
\end{ex}

\section{Discontinuités d'une fonction monotone}

On note $f:[a, b]\in\mathbb R$ croissante. On va montrer que l'ensemble des points de discontinuité de $f$ est fini ou dénombrable.

On note pour $x\in ]a, b[$, \[
    \delta(x)=\lim_{\substack{t\to x\\t>x}}f(t)-\lim_{\substack{t\to x\\t<x}}f(t)
\]

\begin{itemize}
    \item $f$ est croissante donc $\delta$ est positive et $f$ est continue en $a$ ssi $\delta(a)=0$.
    \item Pour tout $n\in\mathbb N^\star$, on note \[
            I_n= \left\{ x\in ]a, b[, \qquad \delta(x)\geq \frac{f(b)-f(a)+1}n \right\}
        \]
    \item Si $I_n$ est de cardinal $\geq n+1$ alors on peut trouver $x_1< \cdots< x_{n+1}$ dans $I_n$, et on se donne $a<t_1<x_1<\cdots <x_{n+1}<t_{n+2}<b$ pour avoir \[
            f(b)-f(a)=\underbrace{f(b)-f(t_{n+2})}_{\geq 0}+\underbrace{f(t_{n+2})-f(b_{n+1})}_{\geq \delta(x_{n+1})} +\cdots +\underbrace{f(t_2)-f(t_1)}_{\geq \delta(x_1)}+\underbrace{f(t_1)-f(b)}_{\geq 0}
        \]
        donc \[
            f(b)-f(a)\geq \delta(x_{n+1})+\cdots +\delta(x_1)\geq n\frac{f(b)-f(a)+1}n
        \]
        absurde donc $I_n$ est fini. $x\in ]a, b[$ est un point de discontinuité si et seulement si $\delta(x)>0$ ssi $\exists n\in\mathbb N, x\in I_n$ ssi \[
            x\in \bigcup_{n\in\mathbb N^\star}I_n
        \]
\end{itemize}

\section{Familles sommables de réels positifs}

Pour un ensemble dénombrable $I$, on note $\mathcal P_f(I)$ l'ensemble des parties finies de $I$

\begin{dfn}
    Soit $(a_i)_{i\in I}\in (\mathbb R_+)^I$. On dira que la famille $(a_i)_{i\in I}$ est \index{famille sommable}\textbf{sommable} si \[
        \left\{ \sum_{j\in J}a_j, \quad J\in\mathcal P_f(I) \right\}
    \]
    est majoré. On note alors \[
        \sum_{i\in I}a_i=\sup\left\{ \sum_{j\in J}a_j, \quad J\in\mathcal P_f(I) \right\}
    \]
    Si la famille n'est pas sommable on peut parfois écrire \[
        \sum_{i\in I}a_i=+\infty
    \]
\end{dfn}

\begin{prop}
    \Hyp $(a_i)_{i\in I}$ famille dénombrable de réels positive et $S\in \mathbb R$
    \Conc $(a_i)_{i\in I}$ est sommable et $S=\sum_{i\in I}a_i$ si et seulement si les deux conditions suivantes sont réunies: \begin{itemize}[left=2cm]
        \item $\forall J\in\mathcal P_f(I), \quad \displaystyle \sum_{j\in J}a_j\leq S$
        \item $\displaystyle \forall \epsilon>0, \exists J\in\mathcal P_f(J), \quad S-\epsilon\leq \sum_{j\in J}a_j$
    \end{itemize}
\end{prop}

\begin{proof}
    C'est une caractérisation de la borne supérieure.
\end{proof}

\begin{rem}
    Dans le cas $I=\mathbb N$ et $(a_i)_{i\in\mathbb N}$ positive, la sommabilité équivaut à la convergence de $\sum a_i$. Si la famille est sommable, $(S_n)$ est croissante majorée donc la série converge, et si la série converge, alors les sommes sur un segment sont majorées par la somme de la série.
\end{rem}

\begin{prop}
    \Hyp $(a_i)_{i\in I}$, $(b_i)_{i\in I}$ des familles de réels positifs
    \begin{concenum}
    \item Si $0\leq a_i\leq b_i$ pour tout $i\in I$ et $(b_i)_i$ est sommable, alors $(a_i)_i$ est sommable et \[
            \sum_{i\in I}a_i\leq \sum_{i\in I}b_i
        \]
    \item Pour $\lambda\in\mathbb R_+$, si $(a_i)_i$ et $(b_i)_i$ sont sommables, $(\lambda a_i+b_i)_i$ aussi et \[
            \lambda\sum_{i\in I}a_i+\sum_{i\in I}b_i=\sum_{i\in I}(\lambda a_i+b_i)
        \]
    \end{concenum}
\end{prop}

\begin{proof}~
    \begin{enumerate}
        \item Ok
        \item Pour tout $J\in \mathcal P_f(I)$, \[
                \sum_{j\in J}(\lambda a_j+b_j)=\lambda\sum_{j\in J}a_j+\sum_{j\in J}b_j\leq \lambda \sum_{i\in I}a_i+\sum_{i\in I}b_i\defeq M
            \]
            Soit $\epsilon>0$. Il existe $J_1, J_2\in\mathcal P_f(I)$ tels que \[
                \sum_{j\in I}a_j-\epsilon'\leq \sum_{j\in J_1}a_j\qquad \text{ et }\qquad \sum_{j\in I}b_j-\epsilon'\leq \sum_{j\in J_2}b_j
            \]
            donc pour $J=J_1\cup J_2$, \[
                M-\lambda\epsilon'-\epsilon'=M-\epsilon\leq \sum_{j\in J}(\lambda a_j+b_j)
            \]
            donc \conc
    \end{enumerate}
\end{proof}

\section{Familles sommables de complexes}

\begin{defprop}
    \Hyp $(a_k)_{k\in I}$ une famille dénombrable de complexes, $\alpha_k=\Re(a_k)$ et $\beta_k=\Im(a_k)$
    \begin{concenum}
    \item $(a_k)_{k\in I}$ est sommable\index{famille sommable} ssi $(|a_k|)_{k\in I}$ est sommable.
    \item Dans ce cas, $(\alpha_k^+)$, $(\alpha_k^-)$, $(\beta_k^+)$ et $(\beta_k^-)$ sont sommables\footnote{$u_n^+=\max(u_n, 0), \quad u_n^-=\max(-u_n, 0)$ pour écrire $u_n=u_n^+-u_n^-$} et on pose \[
            \sum_{i\in I}a_i\defeq\sum_{i\in I}\alpha_i^+-\sum_{i\in I}\alpha_i^-+i\sum_{i\in I}\beta_i^+-i\sum_{i\in I}\beta_i^-
        \]
    \item Si $(a_i)_i$ et $(b_i)_i$ sont sommables et $\lambda\in\mathbb C$ alors $(\lambda a_i+b_i)_{i\in I}$ est sommable et \[
            \lambda\sum_{i\in I}a_i+\sum_{i\in I}b_i=\sum_{i\in I}(\lambda a_i+b_i)
        \]
    \end{concenum}
\end{defprop}

\begin{proof}~
    \begin{enumerate}
        \item Définition
        \item $0\leq \alpha_k^\pm\leq |\alpha_k|\leq |a_k|$ pour la sommabilité
        \item L'inégalité triangulaire donne la sommabilité. La linéarité est admise car la démo est lourde.
    \end{enumerate}
\end{proof}

\begin{thm}[Sommation par paquets\index{sommation par paquets}]
    \Hyp $I$ dénombrable, $(I_n)_{n\in\mathbb N}$ une partition de $I$, $(a_i)_{i\in I}\in \mathbb C^I$
    \Conc Il y a équivalence entre \begin{enumerate}[left=1.1cm,label=\alph{enumi}.]
        \item la famille $(a_i)_i$ est sommable
        \item Pour tout $n\in\mathbb N$, $(a_i)_{i\in I_n}$ est sommable et \[
                \sum_{i\in I_n}a_i
            \]
            est le terme général d'une série absolument convergente. Dans ce cas, \[
                \sum_{i\in I}a_i=\sum_{n\in\mathbb N}\sum_{i\in I_n}a_i
            \]
    \end{enumerate}
\end{thm}

\begin{proof} La démonstration n'est pas au programme. On la fait dans le cas $(a_i)_i$ réelle positive.

    \begin{itemize}
        \item $(b\implies a)$ On prend $J$ une partie finie de $I$, \[
                I=\bigcup_{n\in\mathbb N}I_n
            \]
            donc il existe $n_0$ tel que \[
                J\subset \bigcup_{n\leq n_0}I_n
            \]
            On a \[
                0\leq \sum_{j\in J}a_j\leq \sum_{k=0}^{n_0}\sum{k\in I_n}a_k\leq \sum_{n\in\mathbb N}\sum_{k\in I_n}a_k
            \]
            donc $(a_i)_i$ sommable.
        \item $(a\implies b)$ Pour tout $J$ fini inclus dans $I_n$, $J\subset I$ donc \[
                0\leq \sum_{j\in J}a_j\leq \sum_{j\in I}a_j
            \]
            donc $(a_i)_{i\in I_n}$ est sommable.
            Soit $n_0\in\mathbb N$ et $J_0, \cdots, J_{n_0}$ finis tels que $J_k\subset I_k$.
            $J=J_0\cup \cdots \cup J_{n_0}$ est fini et \[
                \sum_{j\in J}a_j=\sum_{i\in J_0}a_i+\cdots +\sum_{i\in J_{n_0}}a_i\leq \sum_{j\in I}a_j
            \]
            On a alors \[
                \sum_{j\in J_0}a_j\leq \sum_{i\in I}a_i-\sum_{j\in J_1}a_j-\cdots -\sum_{j\in J_{n_0}}a_j
            \]
            et le membre de droite est indépendant de $J_0$. On passe à la borne supérieure par rapport à $J_0$: \[
                \sum_{i\in I_0}a_i=\sup_{\substack{J_0\subset I_0\\J_0\text{ fini}}}\sum_{j\in J_0}a_j\leq \sum_{i\in I}a_i-\sum_{j\in J_1}a_j-\cdots -\sum_{j\in J_{n_0}}a_j
            \]
            puis on recommence pour $J_1, \cdots, J_{n_0}$. On obtient ainsi \[
                \sum_{i\in I_0}a_i+\cdots +\sum_{i\in I_{n_0}}a_i\leq \sum_{j\in I}a_j.
            \]
            Les sommes partielles sont croissantes et majorées donc \[
                \sum \left( \sum_{k\in I_n}a_k \right)
            \]
            est absolument convergente (convergente + termes positifs). Puis, $n_0\to+\infty$ donne \[
                \sum_{n=0}^{+\infty}\sum_{k\in I_n}a_k\leq \sum_{i\in I}a_i
            \]
            Enfin, \[
                \sum_{j\in J}a_j\leq \sum_{k=0}^{n_0}\sum_{j\in I_k}a_j\leq \sum_{k=0}^{+\infty}\sum_{j\in I_k}a_j
            \]
            puis on passe au sup pour $J$ et on trouve la conclusion.
    \end{itemize}
    Pour $(a_i)$ complexe, on utilise la linéarité.
\end{proof}

\begin{cor}[Fubini discret\index{Fubini (théorème discret de -- )}]
    \Hyp $(a_{n, m})_{(n, m)\in\mathbb N^2}\in\mathbb C^{\mathbb N^2}$
    \Conc Il y a équivalence entre: \begin{enumerate}[left=1.2cm,label=\alph{enumi}.]
        \item $(a_{n,m})_{(n,m)\in\mathbb N^2}$ est sommable
        \item $\forall n\in\mathbb N, (a_{n, m})_{m\in\mathbb N}$ est sommable et $\sum_{m\in\mathbb N}|a_{n,m}|$ est le terme général d'une série convergente. Dans ce cas, \[
                \sum_{(n, m)\in\mathbb N^2}a_{n, m}=\sum_{n\in\mathbb N}\sum_{m\in\mathbb N}a_{n,m}=\sum_{m\in\mathbb N}\sum_{n\in\mathbb N}a_{n, m}
            \]
    \end{enumerate}
\end{cor}

\begin{ex}
    \[
        \sum_{n\geq 2}(\zeta(n)-1)
    \]
    On a $\zeta(n)-1=\sum_{k\geq 2}\frac1{k^n}$. La famille $\displaystyle \left( \frac1{k^n} \right)_{n, k\geq 2}$ est-elle sommable ?

    Pour $k\geq 2$, $\left(\frac1{k^n}\right)_{n\geq 2}$ est sommable (géométrique) et \[
        \sum_{n\geq 2}\frac1{k^n}=\frac1{k^2}\frac1{1-\frac1k}=\frac1{k(k+1)}
    \] est le terme général d'une série convergente, donc la famille est sommable et \[
    \sum_{m\geq 2}\sum_{n\geq 2}\frac1{k^n}=\sum_{k\geq 2}\frac1{k-1}-\frac1k=1=\sum_{n\geq 2}\sum_{k\geq 2}\frac1{k^n}=\sum_{n\geq 2}(\zeta(n)-1)
\]
donc \[
    \sum_{n\geq 2}(\zeta(n)-1)=1
\]
\end{ex}

\begin{rem}
    Une suite réelle est sommable si et seulement si elle est le terme général d'une série absolument convergente. Dans ce cas, \[
        \sum_{n\in\mathbb N}a_n=\sum_{n=0}^{+\infty}a_n
    \]
\end{rem}

\begin{prop}[Convergence cumulative]
    \Hyp $(a_n)_{n\in\mathbb N}$ est une suite complexe
    \Conc Il y a équivalence entre \begin{enumerate}[left=1.2cm,label=\alph{enumi}.]
        \item $(a_n)$ est sommable
        \item $\sum a_n$ est absolument convergente
        \item $\forall \sigma\in\mathfrak S(\mathbb N), (a_{\sigma(n)})$ est sommable
    \end{enumerate}
\end{prop}

\begin{proof}~
    \begin{itemize}
        \item $(a\iff b)$ OK
        \item $(b\implies c)$ Soit $\sigma\in\mathfrak S(\mathbb N)$ et $J$ fini. \[
                \sum_{k\in J}|a_{\sigma(k)}|=\sum{k\in\sigma(J)}|a_k|\leq \sum_{k\in\mathbb N}|a_k|
            \]
        \item $(c\implies a)$ $\sigma=\id$
    \end{itemize}
\end{proof}

\begin{thm}[Produit de Cauchy]
    \Hyp $(a_p)_{p\in\mathbb N}$, $(b_q)_{q\in\mathbb N}$ des familles complexes sommables et \[
        c_n=\sum_{p+q=n}a_pb_q=\sum_{k=0}^na_kb_{n-k}=\sum_{k=0}^na_{n-k}b_k
    \]
    \begin{concenum}
    \item $(c_n)_{n\in\mathbb N}$ est sommable
    \item \[
            \sum_{n=0}^{+\infty}c_n= \left( \sum_{p=0}^{+\infty} a_p \right)\cdot \left( \sum_{q=0}^{+\infty}b_q \right)
        \]
    \end{concenum}
\end{thm}

\begin{proof}
    On note \[
        I_n=\{(p, q)\in\mathbb N^2, \quad p+q=n\}
    \]
    de sorte que $(I_n)$ est une partition de $\mathbb N^2$. On note \[
        J_n=\{n\}\times \mathbb N
    \]
    de sorte que $(J_n)$ est aussi une partition de $\mathbb N^2$. La famille $(a_nb_q)_{q\in\mathbb N}$ est sommable et \[
        \sum_{q\in\mathbb N}|a_nb_q|=|a_n|\underbrace{\sum_{q\in\mathbb N}|b_q|}_{\text{fixé}}
    \]
    est le terme général d'une série convergente ($(a_n)$ sommable).
    On a donc $(a_pb_q)_{(p, q)\in\mathbb N^2}$ sommable et on somme par paquets sur les $I_n$:\[
        \sum_{(p, q)\in\mathbb N^2}a_pb_q=\sum_{n=0}^{+\infty}\sum_{(p, q)\in\mathbb N^2}a_pb_q=\sum_{n\in\mathbb N}c_n
    \]
\end{proof}

\section{Exemples}

\subsection{Théorème de réarrangement de Steinitz}

\index{Steinitz (théorème de réarrangement de -- )}

On note $(a_n)_n$ une suite réelle telle que $\sum a_n$ est semi-convergente. On va montrer que pour tout $\ell\in\bar{\mathbb R}$, il existe $\sigma\in\mathfrak S(\mathbb N)$ tel que \[
    \sum_{n=0}^{+\infty}a_{\sigma(n)}=\ell
\]
On note $a_n=a_n^+-a_n^-$ donc $\sum a_n^+$ et $\sum a_n^-$ divergent car $\sum a_n$ semi-convergente. On fixe $\alpha\in\mathbb R$, \[I^+=\{n\in\mathbb N, a_n\geq 0\}\qquad I^-=\{n\in\mathbb N, a_n<0\}\]
On note $(i_n)_n$ (resp. $(j_n)_n$) la suite croissante de tous les éléments de $I^+$ (resp. $I^-$).
On note $k_1$ le plus petit indice $k\geq 0$ tel que \[
    -\alpha+a_{i_1}+\cdots +a_{i_{k_1}}>0
\]
et $\rho_1$ le plus petit indice tel que \[
    -\alpha+a_{i_1}+\cdots +a_{i_{k_1}}+a_{j_1}+\cdots +a_{j_{\rho_1}}<0
\]
On itère le processus pour construire $\sigma=(i_1, \cdots, i_{k_1}, j_1, \cdots j_{\rho_1}, i_{k_1+1}, \cdots, i_{k_2}, \cdots)\in\mathfrak S(n)$ (c'est bijectif car $I^+\sqcup I^-=\mathbb N$).

Puis, \[
    \left| -\alpha+\sum_{k=0}^na_{\sigma(k)} \right|\leq |a_{p_n}|+|a_{q_n}|\xrightarrow[n\to+\infty]{}0
\]
avec $p_n$ le plus grand entier tel que $\rho_{p_n}\leq \sigma(n)$ et $q_n$ le plus petit entier tel que $\sigma(n)\leq \rho_{q_n}$
Pour $\ell=\pm\infty$ c'est le même principe.

\subsection{Inégalité de Carleman}

\index{Carleman (inégalité de -- )}

On suppose que $\sum a_n$

\section{Produits infinis}

Soit $(a_n)$ une suite numérique\index{produit infini} et \[
    p_n=\prod_{k=0}^n(1+a_k)
\]
On dira que le produit $\prod(1+a_k)$ converge si et seulement si $p_n$ a une limite réelle \textbf{non nulle}

\begin{res}
    Si $\sum a_k$ est absolument convergente et $(a_n)$ ne vaut jamais $-1$, alors $\prod(1+a_k)$ converge
\end{res}

\begin{proof}
    On le montre dans le cas réel. On a $a_n\longrightarrow 0$ donc à partir d'un certain rang $n_0$, $(1+a_n)$ est strictement positive. Pour $N>n_0$, \[
        p_n=\underbrace{\prod_{k=0}^{n_0}(1+a_k)}_{\text{fixe } \neq 0}\times \prod_{k=n+1}^{N}\underbrace{(1+a_k)}_{>0}
    \]
    Or pour $k\geq n_0$, $|\ln(1+a_k)|\sim |a_k|$ donc $\sum \ln(1+a_k)$ ACV et sa somme vaut $\ell\in\mathbb R$. On a alors \[
        p_n=e^{\ln p_n}\xrightarrow[n\to+\infty]{}e^\ell \prod_{k=0}^{n_0}(1+a_k) \neq 0
    \]
\end{proof}

\begin{rem}
    Le resultat est vrai avec une suite complexe.
\end{rem}

\begin{ex}
    On note $f:\mathbb N^\star\longrightarrow \mathbb R_+$ complètement multiplicative\footnote{$\forall x,y, f(xy)=f(x)f(y)$} (le résultat que l'on va montrer est toujours vrai pour $f$ multiplicative\footnote{$\forall x,y, x\land y=1\implies f(xy)=f(x)f(y)$})

    On suppose que $\sum f(n)$ est absolument convergente. On va montrer que \[
        \prod_{p\text{ premier }} \frac1{1-f(p)}=\sum_{n=1}^{+\infty}f(n)
    \]
    On note $(p_i)_i$ la suite croissante des nombres premiers. Si il existe un premier $p$ tel que $f(p)=1$ alors $f(p^n)=1$ et $f(n)\xnrightarrow{}0$ ce qui est absurde, donc pour tout $p$ premier, $f(p)\neq 1$.
    On a alors $\sum f(p_i)$ ACV donc et $f(p_i)\neq 1$ donc $\prod (1-f(p_i))$ converge donc $\prod \frac1{1-f(p_i)}$ aussi. Puis, \[
        P_N=\prod_{i=1}^N\frac1{1-f(p_i)}=\prod_{i=1}^N \sum_{k=0}^{+\infty}f(p_i)^k\underset{\Pi\text{ de Cauchy}}=\sum_{k_1, \cdots, k_N}f(p_1^{k_1}\cdots p_N^{k_N})
    \]
    d'où \[
        \sum_{n=0}^{M_N}f(n)\leq P_N\leq\sum_{n=0}^{+\infty}f(n)
    \]
    avec $M_N=\max \{k, \quad \llbracket 1, k\rrbracket\text{ se décompose avec } p_1, \cdots p_N\}$. Puis, $N\to+\infty$ donne le résultat.
\end{ex}

Avec $s>1$ et $f:n\longmapsto \frac1{n^s}$, on trouve l'identité d'Euler\index{Euler!identité}: \[
    \prod_{p\in\mathbb P}\frac1{1-\frac1{p^s}}=\zeta(s)
\]

\endchapter


\chapter{Intégration sur un intervalle quelconque}

\ifsolo
    ~

    \vspace{1cm}

    \begin{center}
        \textbf{\LARGE Intégration sur un intervalle quelconque} \\[1em]
    \end{center}
    \tableofcontents
\else
    \minitoc
\fi
\thispagestyle{empty}

\ifsolo \newpage \setcounter{page}{1} \fi

\section{Rappels}

L'opération d'intégration sur $\mathcal C^0_{PM}([a, b], \mathbb R)$ est linéaire positive croissante. On a \[
    \left| \int_a^bf \right|\leq \int_a^b|f|\leq (b-a)\sup_{[a,b]} |f|
\]
et similairement \[
    \left| \int_a^bfg \right|\leq \sup_{[a,b]}|f|\int_a^b|g|
\]
Si $f$ est continue alors \[
    \left| \int_a^bf \right|=\int_a^b|f|\iff f\text{ de signe constant }
\]
et si $g$ est continue positive, \[
    \inf_{[a,b]}f\int_a^bg\leq \int_a^bfg\leq \sup_{[a,b]}f\int_a^bg
\]
et si $g$ n'est pas identiquement nulle, \[
    \int_a^bg>0\implies \dfrac{\int_a^bfg}{\int_a^bg}\in[\inf f,\sup f]=f([a,b])\implies \exists c, \quad \int_a^bfg=f(c)\int_a^bg.
\]
Le résultat reste vrai si $g$ est identiquement nulle. Si $f$ est continue positive sur $[a, b]$ et $\int_a^bf=0$ alors $f=0$.


% \chapter{Suites et séries de fonctions}

\chapter{Algèbre générale}


\ifsolo
    ~

    \vspace{1cm}

    \begin{center}
        \textbf{\LARGE Algèbre générale} \\[1em]
    \end{center}
    \tableofcontents
\else
    \chapter{Algèbre générale}

    \minitoc
\fi
\thispagestyle{empty}

\ifsolo \newpage \setcounter{page}{1} \fi
\section{Groupes}

\textbf{Rappel.} $(G, \star)$ est un groupe\index{groupe} ssi \begin{itemize}
    \item $\star$ est une loi de composition interne associative
    \item $\star$ a un neutre dans $G$
    \item Tous les éléments de $G$ ont un inverse dans $G$ pour $\star$
\end{itemize}

\begin{notation}
    Si $\star$ est commutative, on la notera en général $+$ et l'inverse sera alors noté $-x$. Sinon, on notera souvent la loi $\cdot$ et l'inverse $x^{-1}$. Le neutre de $G$ se note $e_G$ (parfois $0_G$ ou $0$ dans le cas d'une loi notée additivement)
\end{notation}

Si $G$ est de cardinal fini, alors on note $|G|$ ou $\#G$ son cardinal.

\section{Constructeurs de groupes}

\begin{prop}~
    \begin{enumerate}
        \item Si $(G, \cdot)$ est un groupe et $X$ un ensemble non vide alors on muni l'ensemble de fonction de $X$ dans $G$ de la loi $\star$ en posant \[
                \forall f, g: X\longrightarrow G, \quad f\star g:x\longmapsto f(x)\cdot g(x)
            \]
            $(\mathcal F(X, G), \star)$ est alors un groupe.
        \item Soit $(G, \land)$, $(G', \star)$ deux groupes. On munit $G\times G'$ de la loi $\cdot$ en posant \[
                (x, x')\cdot (y, y')=(x\land x', y\star y')
            \]
            $(G\times G', \cdot)$ est alors un groupe.
    \end{enumerate}
\end{prop}

\begin{proof}
    Facile mais pénible.
\end{proof}

\section{Sous-groupes}

\begin{thmdef}
    \Hyp $(G, \cdot)$ un groupe, $H\subset G$
    \Conc On dira que $H$ est un sous-groupe\index{groupe!sous-groupe} de $G$ si $H\neq \emptyset$, $\cdot$ est une l.c.i. sur $H$ et $(H, \cdot)$ est un groupe. Dans ce cas, $e_H=e_G$. Il y a équivalence entre: \begin{enumerate}[label=(\alph*)]
        \item $H$ est un sous-groupe de $G$
        \item $H\subset G$, $H\neq \emptyset$ et $\forall x, y\in H, x\cdot y^{-1}\in H$
    \end{enumerate}
\end{thmdef}

\begin{notation}
    $(G, \cdot)$ groupe, $H\subset G$, $a\in G$. \[
        aH=\{a\cdot h, \quad h\in H\}
    \]
    \[
        H^{-1}={h^{-1}, \quad h\in H}
    \]
    et si $A\subset G$,
    \[
        AH = \{a\cdot h, \quad a\in A, h\in H\}
    \]
\end{notation}

Ainsi $H\subseteq G$ est un sous-groupe de $G$ si et seulement si $H\neq \emptyset$ et $HH^{-1}\subseteq H$

\begin{exo}
    Soient $H, K$ des sous-groupes de $G$. Montrer que $HK$ est un sous groupe si et seulement si $HK=KH$.
\end{exo}

\begin{proof}[Résolution]
    Si $HK$ sous-groupe de $G$ alors $HK=(HK)^{-1}=K^{-1}H^{-1}=KH$. Si $HK=KH$ alors $HK(HK)^{-1}=HKK^{-1}H=HKKH=HHKK\subseteq HK$
\end{proof}

\subsection{Centre d'un groupe}

\begin{dfn}
    Le centre d'un groupe\index{centre (d'un groupe)}\index{groupe!sous-groupe!centre} $G$ est le groupe \[
        Z(G)=\{x\in G, \quad \forall y\in G, xy=yx\}
    \]
\end{dfn}

\begin{rem}
    C'est un sous-groupe de $G$, on peut le vérifier manuellement.
\end{rem}

\subsection{Normalisateur}

On note $H$ un sous-groupe de $G$. On appelle \textbf{normalisateur}\index{normalisateur} de $H$ l'ensemble \[
    N(H)=\{x\in G, \quad xHx^{-1}=H\}
\]
On peut vérifier que c'est un sous-groupe de $G$. On dira que $H$ est un sous-groupe distingué \index{groupe!sous-groupe!distingué} de $G$ si $N(H)=G$.

\begin{rem}
    On note $\mathcal R$ la relation \[
        x\mathcal Ry\iff y^{-1}x\in H
    \]
    C'est bien une relation d'équivalence (réfléxive symétrique transitive) et si $a\in H$, $x\in \bar a \iff x\in aH$. De même, $x\mathcal R'y\iff xy^{-1}\in H$ définit une autre relation d'équivalence dont les classes sont les $Ha$. Lorsque $H$ est distingué (on notera $H\triangleleft G$), $aH=Ha$
\end{rem}

\subsection{Sous-groupe de torsion}

Soit $G$ un groupe \emph{abélien}. On note \[
    \tau(G)=\{g\in G, \quad \exists n\in\mathbb N^\star, g^n=e_G\}
\]
C'est un sous-groupe\footnote{on a besoin de l'hypothèse de commutativité de $G$, ça n'est pas vrai sinon} de $G$ appelé sous-groupe de torsion\index{groupe!sous-groupe!de torsion}.

\subsection{Sous-groupes de $\mathbb Z$ et $\mathbb R$}

On va montrer que les sous-groupe de $\mathbb Z$ sont de la forme $d\mathbb Z$, et $d$ est unique au signe près.
On va pour cela montrer la propriété plus générale suivante (hors-programme): si $A$ est un anneau euclidien, alors c'est un anneau principal.

\begin{dfn}[Hors-Programme]
    Un anneau $A$ commutatif est dit euclidien\index{anneau!euclidien} lorsqu'il existe une application $v:A\setminus\{0\} \to \mathbb N$ qui satisfait les propriétés suivantes:
    \begin{enumerate}
        \item \[
                \forall (a, b)\in A\times A\setminus\{0\}, \exists q, r\in A, \qquad a=qb+r\quad \text{ et } \quad r=0\text{ ou } v(r)<v(b)
            \]
        \item \[
                \forall a, b\in A\setminus\{0\}, \qquad v(b)\leq v(ab)
            \]
    \end{enumerate}
    Une telle application sera appelée \textbf{stathme euclidien}\index{anneau!euclidien!stathme, préstathme} sur $A$. Une application qui ne satisfait que la première propriété sera appelée \textbf{préstathme euclidien} sur $A$.

    Un anneau $A$ commutatif sera dit principal\index{anneau!principal} lorsque tous ses idéaux\footnote{notion abordée plus en détail dans la section \textbf{Idéal d'un anneau}} (sous-groupes $I$ de $(A, +)$ tels que pour tout $a\in A$, $aI\subset I$) sont principaux (i.e. de la forme $aA$)
\end{dfn}

L'anneau $\mathbb Z$ est euclidien (la valeur absolue est un stathme euclidien sur $\mathbb Z$). Si $H$ est un sous-groupe de $\mathbb Z$ alors pour tout $k\in\mathbb Z$, $kH=H+H+\cdots +H\subset H$ donc $H$ est un idéal de l'anneau $\mathbb Z$. Il sera donc suffisant de montrer que tout anneau euclidien est principal.

\begin{prop}[Hors-Programme]
    Un anneau euclidien est principal
\end{prop}

\begin{proof}
    Soit $A$ un anneau commutatif euclidien (de stathme $v$), et $I$ un idéal de $A$. On va montrer que $I$ s'écrit $aA$ pour un $a\in A$. Si $I=\{0\}$ alors $a=0$ convient.

    On suppose $I\neq \{0\}$. L'ensemble $v(I\setminus\{0\})$ admet un minimum, on note $a\in I\setminus\{0\}$ un élément minimal pour $v$. On a clairement $aA\subset I$ car $a\in I$ et $I$ idéal. On prend $b\in I$ et \[
        \exists q, r, \qquad b=qa+r\quad \text{ et }\quad r=0\text{ ou }v(r)<v(a).
    \]
    On a $r=b-qa\in I$ donc $v(r)<v(a)$ est impossible par définition de $a$, donc $r=0$ et $b=qa$. On a donc $I\subset aA$ d'où $aA=I$
\end{proof}

\begin{rem}
    Si $I$ est un idéal principal, il s'écrit $I=aA$. S'il s'écrit également $a'A$ alors $a\;|\;a'$ et $a'\;|\;a$ donc $a=\varepsilon a'$ pour $\varepsilon\in\mathcal U(A)$ (le groupe des inversibles de $A$). Dans $\mathbb Z$, les inversibles sont $-1$ et $1$, d'où l'unicité du générateur positif.
\end{rem}

\begin{rem}
    On aurait pu écrire la démonstration directement pour $A=\mathbb Z$, $v:x\mapsto |x|$
\end{rem}

On va maintenant montrer qu'un sous groupe $H$ non trivial de $\mathbb R$ est soit dense dans $\mathbb R$ soit de la forme $a\mathbb Z$. On note $\ell=\inf H\cap \mathbb R_+^\star$. Il y a deux cas: \begin{itemize}
    \item $\ell>0$. Si $l\not\in H$ alors il existe $a\in\mathbb R_+^\star$ tel que $a\in ]\ell, 2\ell[$ et $a\in H$, il existe $b\in H$ tel que $b\in ]\ell, a[$ donc $a-b\in H\cap \mathbb R_+^\star$ ce qui est absurde car $a-b>\ell$ par définition de $\ell$. Donc $\ell \in H$ et $\ell \mathbb Z\subset H$. Pour l'autre inclusion, fait comme pour $\mathbb Z$: on écrit $x=\lfloor x/\ell \rfloor \ell + r$ et $r=0$.
    \item $\ell=0$. Soient $0<x<y$. Il existe $\alpha\in H\cap \mathbb R_+^\star$ tel que $0<\alpha<y-x$ et \[
            x<\underbrace{\floor{\frac x\alpha}+\alpha}_{\in H}<x+y-x=y
        \]
        donc $H$ est dense dans $\mathbb R$ (les autres configurations de $x$ et $y$ se traitent similairement)
\end{itemize}

Par exemple, $a\mathbb Z+b\mathbb Z$ est dense dans $\mathbb R$ si et seulement si $\frac ab\not\in \mathbb Q$, de sorte que $ \cos(\mathbb N) $ est dense dans $\mathbb R$ car $\mathbb Z+2\pi\mathbb Z$ est dense dans $\mathbb R$

\section{Opérations sur les sous-groupes}

\begin{prop}
    \Hyp $(H_i)_{i\in I}$ famille de sous-groupes de $G$
    \Conc L'ensemble \[
        \bigcap_{i\in I}H_i
    \]
    est un sous-groupe de $G$
\end{prop}

\begin{rem}
    En général, l'union de deux sous-groupes n'est pas un sous-groupe
\end{rem}

\section{Sous-groupes engendrés}

\begin{thmdef}
    \Hyp $(G, \cdot)$ un groupe, $A\subset G$, $S_A=\{H\in\mathcal P(G), A\subset H, H\text{ sg de }G\}$
    \begin{concenum}
    \item On appelle sous-groupe engendré\index{groupe!engendré} par $A$ l'ensemble \[
            \langle A\rangle = \bigcap_{H\in S_A}H
        \]
    \item Si $H$ est un sous-groupe de $G$ contenant $A$ alors $\langle A\rangle \subseteq H$
    \item \[
            \langle A\rangle =\{e_G\}\cup \{x_1\cdots x_n, \quad n\in\mathbb N, x_1, \cdots, x_n\in A\cup A^{-1}\}
        \]
    \end{concenum}
\end{thmdef}

\begin{proof}~
    \begin{enumerate}
        \setcounter{enumi}{2}
    \item On note $E=\{e_G\}\cup \{x_1\cdots x_n, \quad n\in\mathbb N, x_1, \cdots, x_n\in A\cup A^{-1}\}$. On a \[
            E\subset \langle A\rangle\qquad \text{ car }A,A^{-1}\subset \langle A\rangle
        \]
        et $A\subset E$ donc $\langle A\rangle \subset E$ par 2. donc $E=\langle A\rangle$
\end{enumerate}
\end{proof}

\begin{ex}
    \begin{itemize}
        \item $\langle \exp \left( \frac{2i\pi}n \right)\rangle=\mathbb U_n$
        \item $\mathbb Z=\langle 1\rangle = \langle 2, 3\rangle$
        \item $\mathfrak S_n=\langle \tau_{1,2}, (1\;\cdots\; n)\rangle$
    \end{itemize}
\end{ex}

\begin{dfn}
    Un groupe $G$ est dit monogène\index{groupe!monogène} s'il existe $a\in G$ tel que $G=\langle a\rangle$. Si de plus $G$ est fini, on dira qu'il est cyclique\index{groupe!cyclique}
\end{dfn}


\section{Morphismes}

\begin{dfn}
    Soient $(G,\cdot), (G',\star)$ deux groupes. Une application $f:G\to G'$ est un morphisme ssi \[
        \forall x, y\in G, f(x\cdot y)=f(x)\star f(y)
    \]
\end{dfn}

\begin{notation}
    $f:G\to G'$ est un endomorphisme ssi c'est un morphisme. C'est un isomorphisme si c'est une bijection. Si $G=G'$, on dira qu'un morphisme de $G$ dans $G'$ est un endomorphisme de $G$. Un endomorphisme bijectif est appelé automorphisme. On note $\Hom (G, G')$ l'ensemble des morphisme de $G$ dans $G'$ et $\Aut(G)$ l'ensemble des automorphismes de $G$.\index{groupe!isomorphisme} \index{groupe!endomorphisme}
\end{notation}

\begin{prop}
    \Hyp $f\in\Hom(G, G')$
    \begin{concenum}
    \item $f(e_G)=e_G$
    \item $\forall x\in G, \quad f(x^{-1})=f(x)^{-1}$
    \item On note $\Ker f=f^{-1}(\{e_G\})$. \begin{enumerate}
            \item $\Ker f$ est un sous groupe de $G$
            \item $f$ est injective si et seulement si $\Ker f=\{e_G\}$
        \end{enumerate}
    \item $f^{-1}$ est un isomorphisme si $f$ est un isomorphisme.
    \item $g\circ f$ est un morphisme si $f:G\to G'$ et $g:G'\to G''$ le sont
    \end{concenum}
\end{prop}

\begin{csq}
    $(\Aut(G), \circ)$ est un groupe
\end{csq}

\begin{csq}
    Les $i_a:x\in G\longmapsto axa^{-1}$ pour $a\in G$ sont appelés automorphismes intérieurs\index{groupe!automorphismes intérieurs}. $\varphi:a\longmapsto i_a$ est un morphisme.
\end{csq}

Les morphismes ont la propriété intéressante suivante:

\begin{prop}
    L'image et la préimage d'un groupe par un morphisme est un groupe.
\end{prop}

\begin{proof}
    Facile mais pénible
\end{proof}

\section{Groupes finis}

\subsection{Théorème de Lagrange}

\index{Lagrange!théorème}

On note $G$ un groupe fini et $H$ un sous groupe de $G$. On va montrer que $|H|$ divise $|G|$. On a vu que \[
    x\mathcal Ry\iff x^{-1}y\in H
\]
définit une relation d'équivalence sur $G$, les classes d'équivalence forment donc une partition de $G$. On a ainsi ($\sqcup$ désigne l'union disjointe) \[
    G=\bigsqcup_{i=1}^p a_iH \implies |G|=p|H|
\]

\subsection{\texorpdfstring{$|G|=|\Ker f|\cdot |\Img f|$}{\#G=\#Ker f × \#Im f}}

On note $G$ un groupe fini et $f:G\to G'$ un morphisme. On note $\Img f=\{y_1, \cdots, y_p\}$ de sorte que \[
    G=\bigsqcup_{i=1}^p f^{-1}(\{y_i\}) \implies |G|=\sum_{i=1}^p \left| f^{-1}(\{y_i\}) \right|
\]
On se donne $y\in \Img f$ fixé et $x\in f^{-1}(\{y\})$. On pose \[
    \varphi: t\in \Ker f\longmapsto xt\in f^{-1}(\{y\}).
\]
Cette application est injective et si $u\in f^{-1}(\{y\})$ alors $u=\varphi(x^{-1}u)$ donc elle est aussi surjective. Elle est donc bijective, d'où \[
    |f^{-1}(\{y\})|=|\Ker f|\quad \implies \quad |G|=|\Ker f|\cdot |\Img f|
\]

\subsection{Ordre d'un élément}

\begin{defprop}
    \Hyp $G$ un groupe fini, $a\in G$
    \begin{concenum}
    \item $\varphi:\mathbb Z\longrightarrow G, n\longmapsto a^n$ est un morphisme. Il existe un unique entier naturel non nul tel que $\Ker \varphi=d\mathbb Z$, qu'on appelle \textbf{ordre} de $a$ dans $G$\index{groupe!ordre}, noté $\ord_G (a)$ ou $\ord (a)$
    \item $\ord (a)=|\langle a\rangle|=\min\{n\in\mathbb N^\star, a^n=e_G\}$
    \item $a^n=e_G \iff \ord(a)\;|\; n$ et \[
            \ord(a^p)=\frac{\ord (a)}{p\land \ord (a)}
        \]
    \end{concenum}
\end{defprop}

\begin{proof}~
    \begin{enumerate}
        \item $\varphi$ n'est pas injective car $G$ fini donc $\Ker \varphi$ est un sous-groupe non trivial de $\mathbb Z$, ce qui conclut
        \item $d=\min \Ker\varphi\cap \mathbb N^\star=\min\{n\in\mathbb N^\star, a^n=e_G\}$ puis \[
                a^k=a^{dq+r}=a^r
            \] avec $r\in\llbracket 0, d-1\rrbracket$ donc $\langle a\rangle = \{e_G, a, \cdots, a^{d-1}\}$ et pour $0\leq r_1<r_2<d$, $a^{r_1}=a^{r_2} \implies r_2-r_1\in d\mathbb Z\cap \rrbracket 0, d-1 \llbracket $ absurde donc $|\langle a \rangle |=d$
        \item $a^n=e_G\iff n\in\Ker \varphi=d\mathbb Z\iff d\;|\;n$ puis \[
                \left( a^p \right)^{\frac{\ord (a)}{p\land \ord(a)}}=a^{k\ord(a)}=e_G\quad \text{ donc }\quad \ord(a^p)\;\Big|\;\frac{\ord(a)}{p\land\ord(a)}
            \]
            et \[
                (a^p)^{\ord (a^p)}=e_G \implies \ord(a)\;\Big|\;p\ord (a^p) \implies \frac{\ord(a)}{p\land \ord(a)}\;\Big|\; \frac{p}{p\land \ord(a)}\times \ord(a^p)
            \]
            donc par le lemme de Gauss, \[
                \frac{\ord(a)}{p\land \ord(a)}\;\Big|\; \ord(a^p)
            \]
            d'où l'égalité
    \end{enumerate}
\end{proof}

\begin{rem}
    Si $G$ est fini abélien et $a, b\in G$, \[
        (ab)^{\ord(a)\lor \ord(b)}=e_G\implies \ord (ab)\;|\; \ord(a)\lor \ord(b)
    \]
    Si $\ord (a)$ et $\ord (b)$ sont premiers entre eux alors \[
        (ab)^{\ord(ab)\ord (b)}=e_G=a^{\ord(ab)\ord(b)}\implies \ord (a)\;|\;\ord (b)\ord(ab)\underset{\text{Gauss}}\implies \ord(a)\;|\;\ord(ab)
    \]
    et par symétrie, $\ord(b)\;|\;\ord (ab)$ don $\ord(ab)=\ord(a)\ord(b)$
\end{rem}

\begin{prop}
    \Hyp $(G, \cdot)$ groupe fini, $a\in G$
    \Conc $\ord(a)\;|\;\# G$
\end{prop}

\begin{proof}
    C'est évident avec le théorème de Lagrange. Sinon, \[
        \prod_{x\in G}x=\prod_{x\in G}ax=a^{|G|}\prod_{x\in G}x \implies a^{|G|}=e_G
    \]
\end{proof}

\subsection{Groupe de cardinal $p^2$}

On note $p$ un nombre premier et on suppose que $G$ est un groupe d'ordre $p^2$. Les seuls ordres possibles sont $1, p$ et $p^2$. S'il y a un élément d'ordre $p^2$, alors $G=\langle a\rangle$ est abélien et $a^p$ est d'ordre $p$. Sinon, tous les éléments sont d'ordres $p$. Dans les deux cas, il y a un élément d'ordre $p$.

\subsection{Groupe de cardinal divisible par $p$}

On note $p$ premier et $G$ un groupe d'ordre $kp$. On va montrer qu'il existe un élément d'ordre $p$. Il suffit de montrer qu'un élément est d'ordre $p$. On note \[
    E=\{(x_1, \cdots, x_p)\in G^p, \quad x_1 \cdots x_p=e_G\}\qquad \qquad \#E=\#G^{p-1}.
\]
Cet ensemble est stable par permutation circulaire (car $x_1\cdots x_p=e_G \implies x_2\cdots x_px_1=e_G$). On note $\sigma$ la permutation circulaire $(1\;\; 2 \;\; \cdots \;\; p)$. On va faire agir $\langle \sigma \rangle$ sur $E$. Pour $C\in \mathfrak S_n$, $x=(x_1, \cdots, x_p)\in E$, on note $C.x=(x_{\sigma(1)}, \cdots, x_{\sigma(p)})$.

Pour $x\in E$, on note $O_x=\{C.x, C\in\langle \sigma\rangle\}$ l'orbite de $x$. La relation $x\mathcal Ry\iff y\in O_x$ est une relation d'équivalence sur $E$ de sorte que les $O_x$ forment une partition de $E$. On note \[
    \varphi_x:c\in \langle \sigma \rangle \longmapsto c.x\in O_x.
\]
Si $\varphi_x$ est injective alors $O_x=\Img \varphi_x$ et $|O_x|=|\langle \sigma \rangle|=p$. Sinon, il existe $1\leq k_1 < k_2\leq p$ tel que $\sigma^{k_1}.x=\sigma^{k_2}.x$ donc $\sigma^{k_2-k_1}.x=x$. On a alors $\forall n\geq 0, \sigma^n.x=\sigma^{n+k_2-k_1}.x$ donc $\forall i\in\llbracket 1, p\rrbracket$, $x_i=x_{i+n(k_2-k_1)\pmod p}$ or $k_2-k_1$ engendre $(\mathbb Z_p, +)$ (car $p$ premier) donc $(n(k_2-k_1))_n$ parcourt tous les entiers modulo $p$ et $x=(x_1, \cdots, x_1)$, d'où finalement $\#O_x=1$

Ainsi, les orbites sont de cardinal $p$ ou $1$. On note $O_{x_1}, \cdots, O_{x_r}$ les orbites de cardinal $p$ et $O_{x_1'}, \cdots, O_{x_s'}$ celles de cardinal $1$. On a \[
    \#E=rp+s=\#G^{p-1}\equiv 0\pmod p \implies r\equiv 0\pmod p
\]
donc $r\geq p\geq 2$ et il existe $x=(a, \cdots, a)\neq (e_G, \cdots, e_G)$ dans $E$, donc $a^p=e_G$ et $a\neq e_G$ donc $\ord (a)=p$ (car $p$ premier n'est divisible que par $1$ et $p$).

\subsection{Moyenne dans un groupe}

On note $E$ un $K$-e.v., $G$ un sous-groupe de $\mathrm{GL}(E)$ fini d'ordre $n$. On note \[
    p= \frac{1}{|G|} \sum_{g\in G} g
\]
Si $g_0\in G$, alors $g\longmapsto g_0\circ g$ est bijective et \[
    g_0 p=\frac1{|G|}\sum_{g\in G}g_0 g=p=pg_0
\]
donc \[
    p\circ p=p\circ \left( \frac1{|G|}\sum_{g\in G} g\right)=\frac1{|G|}\sum_{g\in G}\underbrace{pg}_{=p}=\frac1{|G|}|G|p=p
\]
donc $p$ est un projecteur


\begin{exo}
    \begin{itemize}
        \item ~[3/2] On note $f\in \mathcal L(E)$, \[
                f^\star = \frac1{|G|}\sum_{g\in G}g^{-1}\circ f\circ g
            \]
            Montrer que $f^\star$ commute avec les éléments de $G$ et $f^{\star\star}=f^\star$
        \item ~[5/2] Montrer que si $F$ est un s.e.v stable par tous les éléments de $G$ alors il a un supplémentaire stable par tous les éléments de $G$.
    \end{itemize}
\end{exo}

\section{Groupes cycliques}

\begin{dfn}
    Un groupe cyclique est un groupe fini monogène\index{groupe!cyclique}
\end{dfn}

\subsection{Sous-groupes d'un groupe cyclique}

On note $G=\langle a\rangle$ cyclique d'ordre $n$. On note $H$ un sous-groupe de $G$, et $r=\min\{p\in\mathbb N^\star, a^p\in H\}$. On a \begin{itemize}
    \item $\langle a^r\rangle\subseteq H$
    \item Soit $x=a^m\in H$. On écrit $m=rq+s$, $s\in\llbracket 0, r-1\rrbracket$ de sorte que $a^s\in H$ donc $s=0$ et donc $x\in \langle a^p\rangle=H$
\end{itemize}

Ainsi, $H$ est cyclique.

\subsection{Générateurs d'un groupe cyclique}

Avec les mêmes notations, on considère\index{Euler!indicatrice} \[
    \varphi: n\longmapsto \{k\in\llbracket 1, n\rrbracket, \quad k\land n=1\}.
\]
On appelle cette application l'indicatrice d'Euler.
$g=a^p$ est un générateur de $G$ ssi $\ord(a^p)=n$ ssi $n=\frac{n}{p\land n}$ ssi $p\land n=1$. Il y a donc $\varphi(n)$ générateurs de $G$.

Plus généralement, il y a $\varphi(d)$ éléments d'ordre $d$ (on peut le voir en appliquant ce raisonnement à un sous-groupe d'ordre $d\;|\; n$). On a ainsi \[
    \sum_{d\;|\; n}\varphi(d)=n
\]

\subsection{Sous-groupes d'ordre $d$}

On note $H$ un sous-groupe de cardinal $d$ de $G$ (donc $d$ divise $n$).
\begin{itemize}
    \item $H$ a $\varphi(d)$ générateurs (donc d'ordre $d$)
    \item $G$ a $\varphi(d)$ éléments d'ordre $d$
\end{itemize}
donc $H$ contient tous les éléments d'ordre $d$, il est donc unique.

\section{Le groupe \texorpdfstring{$\mathfrak S_n$}{symétrique d'ordre n}}

Le $n$-ième groupe symétrique\index{groupe!symétrique} $(\mathfrak S_n, \circ)$ est un groupe d'ordre $n!$. On note $\tau_{i, j}$ la transposition de support $\{i, j\}$. Si $\sigma \in \mathfrak S_n$, alors $\sigma$ s'identifie au $n$-uplet \[
    (\sigma(1), \cdots, \sigma(n))
\]

On note $c=(i_1\;\;i_2\;\;\cdots \;\;i_p)$ la permutation \[ 
    c: n \in\llbracket 1, n\rrbracket \longmapsto \begin{cases}
        i_{k+1}& \text{ si }n=i_k, k<p\\
        i_1 &\text{ si }n=i_p\\
        n & \text{sinon}
    \end{cases}
\] 
On dira que $c$ est un cycle de support $\{i_1, \cdots, i_p\}$ et d'ordre\footnote{il s'agit bien de l'ordre au sens de $\ord_{\mathfrak S_n}(c)$} $p$ (on prend les $i_k$ deux à deux distincts)

\subsection{Décomposition en cycles}

On note $\sigma \in\mathfrak S_n$ et $O_1, \cdots O_p$ les orbites pour l'action de $\langle \sigma \rangle$ sur $\llbracket1, n\rrbracket$ (c'est à dire les classes d'équivalence pour $x\mathcal Ry\iff \exists k, y=\sigma^k(x)$)

Les orbites s'écrivent $O_k=\{i_k, \sigma(i_k), \cdots, \sigma^{p_k}(i_k)\}$. On note $c_k=(i_k\;\;\sigma(i_k)\;\;\cdots\;\;\sigma^{p_k}(i_k))$ de sorte que \[
    \sigma=c_1\circ \cdots \circ c_p
\]
On a alors \[
    \sigma^k=\id \iff c_1^k=\cdots =c_p^k=\id \iff \ord(c_1)\lor \cdots \lor \ord(c_p)\;|\;k
\]

\subsection{Signature}

Une inversion pour $\sigma \in \mathfrak S_n$ est un couple $(i, j)\in\llbracket 1, n\rrbracket^2$ tel que $i<j$ et $\sigma(j)<\sigma(i)$. On note $N_\sigma$ le nombre de ces inversions et \[
    \varepsilon: \sigma\longmapsto (-1)^{N_\sigma}
\]
est appelée \textbf{signature}\index{signature (d'une permutation)} de $\sigma$

\begin{prop}
    \begin{enumerate}
        \item \[
                \forall \sigma\in\mathfrak S_n, \qquad \varepsilon(\sigma)=\prod_{i<j}\frac{\sigma(i)-\sigma(j)}{i-j}=\prod_{\{i, j\}\in \mathcal P_2(\llbracket 1, n\rrbracket)}\frac{\sigma(i)-\sigma(j)}{i-j}
            \]
        \item $\varepsilon:\mathfrak S_n\to \mathbb U_2$ est un morphisme
        \item La signature d'une transposition vaut $-1$
    \end{enumerate}
\end{prop}

\begin{proof}~
    \begin{enumerate}
        \item La seconde égalité est claire. Puis, \[
                \left| \prod_{i<j}\frac{\sigma(i)-\sigma(j)}{i-j} \right|=\frac{\prod_{i<j}|\sigma(i)-\sigma(j)|}{\prod_{i<j}|i-j|}=1
            \]
            et le signe est bon.
        \item $\sigma, \sigma'\in\mathfrak S_n$. \[
                \varepsilon(\sigma \circ \sigma')=\prod_{i<j}\frac{\sigma\circ \sigma'(i)-\sigma\circ \sigma'(j)}{\sigma'(i)-\sigma'(j)}\times \frac{\sigma'(i)-\sigma'(j)}{i-j}=\varepsilon(\sigma)\varepsilon(\sigma')
            \]
        \item On compte les inversions de $\tau_{1, 2}$ et $\sigma\circ \tau_{i, j}\circ \sigma^{-1}=\tau_{1, 2}$.
    \end{enumerate}
\end{proof}

\subsection{Théorème de Cayley}

On note $G=\{g_1, \cdots, g_n\}$ un groupe fini. On a \[
    \varphi: G\to \mathfrak S(G), x\longmapsto (t\longmapsto xt)
\]
Cette application est injective car \[
    \varphi(x)=\varphi(x')\implies \varphi(x)(e_G)=\varphi(x')(e_G)\implies x=x'
\]
C'est un morphisme car $\varphi(xx')(t)=xx't=x\varphi(x')(t)=\varphi(x)\circ \varphi(x')(t)$.
Ainsi, $G$ est isomorphe à $\varphi(G)$ sous-groupe de $\mathfrak S(G)$ qui est isomorphme à $\mathfrak S_n$. Donc, $G$ est isomorphe à un sous-groupe de $\mathfrak S_n$.

\section{Anneaux et corps}

\begin{dfn}
    \index{anneau} Un ensemble $A$ muni des lois $+$ et $\times$ est un anneau si \begin{itemize}
        \item $(A, +)$ est un groupe abélien de neutre $0_A$
        \item $(A\setminus\{0\}, \times)$ est un monoïde\footnote{mêmes hypothèses que la structure de groupe sans l'inversibilité} de neutre $1_A$
        \item $\times$ est distributive par rapport à $+$
    \end{itemize}
\end{dfn}

\begin{dfn}
    Un anneau commutatif est intègre\index{anneau!intègre} si tous ses éléments sont réguliers pour $\times$
\end{dfn}

\begin{dfn}
    $(A, +, \times)$ est un corps\index{corps} si $(A\setminus\{0\}, \times)$ est un groupe abélien.
\end{dfn}

\begin{rem}
    Un anneau intègre fini est un corps, car $\forall a\in A\setminus\{0\}, x\longmapsto ax$ est injective donc bijective donc tous les éléments sont inversibles.
\end{rem}

\subsection{Calcul dans un anneau}

\begin{prop}[Binôme de Newton]
    \Hyp $a, b\in A$ un anneau, $ab=ba$.
    \Conc \index{Newton!binôme}\[
        (a+b)^n=\sum_{k=0}^n\binom nk a^kb^{n-k}
    \]
\end{prop}

\begin{proof}
    Vu en sup, facile mais long.\footnote{«Moralement, c'est vrai»}
\end{proof}

\begin{rem}
    On prend la convention \[
        \forall k\not\in\llbracket 0, n\rrbracket, \qquad \binom nk=0
    \]
\end{rem}

La formule du binôme se généralise ainsi: pour $a_1, \cdots, a_p\in A$ qui commutent $2$ à $2$ et pour $n\in\mathbb N$, \[
    (a_1+\cdots +a_p)^n=\sum_{a_1+\cdots +a_p=n}\binom{n}{n_1, \cdots, n_p}a_1^{n_1}\cdots a_p^{n_p}
\]
avec \[
    \binom{n}{n_1, \cdots, n_p}=\frac{n!}{n_1!\cdots n_p!}
\]
Pour le voir, on va développer $(a_1+\cdots +a_p)\times\cdots\times(a_1+\cdots +a_p)$. On obtient une somme de termes en $a_1^{n_1}\cdots a_p^{n_p}$ avec $n_1+\cdots +n_p=n$. \begin{itemize}
    \item Il y a $\binom n{n_1}$ blocs pour $a_1$
    \item Il y a $\binom{n-n_1}{n_2}$ pour $a_2$
    \item $\cdots$
    \item Il y a $\binom{n-n_1-\cdots -n_{p-1}}{n_p}$ blocs pour $a_p$
\end{itemize}

Donc les coefficients valent \[
    \prod_{i=1}^p\binom{n-n_1-\cdots -n_{i-1}}{n_i}=\frac{n!}{n_1!\cdots n_p!}=\binom n{n_1,\cdots, n_p}
\]

\textbf{Techniques pour les sommes de coefficients binomiaux}
\begin{itemize}
    \item $(1+X)^p(1+X)^q=(1+X)^{p+q}$ donc \[
    \binom{p+q}\ell=\sum_{i+j=\ell}\binom pi\binom qj
\]
En particulier, \[
    \binom {2n}n=\sum_{i=0}^p\binom pi^2
\]
\item Si on note $S:x\longmapsto (1+x)^n$ \[
        S'(x)=n(1+x)^{n-1}=\sum_{k=0}^n\binom nk kx^{k-1}
    \]
    donc \[
        \sum_{k=0}^n\binom nkk=n2^{n-1}
    \]
\item \textbf{Sommes lacunaires}: On note $p\in\mathbb N^\star$ fixé, on va calculer \[
        S_n=\sum_{pk\leq n}\binom n{pk}
    \]
    On note $\omega=\exp \left(  \frac{2i\pi}p\right)$ de sorte que \[
        1+\omega^k+\cdots +\omega^{(p-1)k}= \begin{cases}
            p&\text{ si } w^k=1 \text{ i.e. } p\;|\; k \\ 0 &\text{ sinon }
        \end{cases}
    \]
    On pose $f(x)=(1+x)^n$ et \[
        f(1)+f(\omega)+\cdots +f(\omega^{p-1})=\sum_{k=1}\binom nk(1+\cdots +\omega^{k(p-1)})=pS_n
    \]
    d'où \[
        S_n=\frac1p \left( f(1)+\cdots +f(\omega^{p-1}) \right)
    \]
\end{itemize}

\subsection{Élément nilpotent}

\begin{dfn}
    Un élément $a\in A$ ($A$ est un anneau) est \textbf{nilpotent}\index{nilpotence} si il existe un entier naturel $n$ tel que $a^n=0_A$. On appelle \textbf{indice de nilpotence}\index{nilpotence!indice de nilpotence} l'entier naturel minimal tel que $a^n=0_A$
\end{dfn}

Si $a, b$ sont nilpotents et commutent alors \[
    (a+b)^{n+m}=\sum_{k=0}^{n+m}\binom{n+m}k a^kb^{n+m-k}=0_A
\]
avec $n$ et $m$ les indices de nilpotence respectifs de $a$ et $b$. Ainsi, $a+b$ est nilpotent. Si $A$ est un anneau commutatif alors l'ensemble $\mathcal N$ des éléments nilpotents est un idéal de $A$. Pour $a$ nilpotent, $a^n=0_A$ et \[
    1=1-a^n=(1-a)(1+a+\cdots +a^{n-1})
\]
donc $a-1$ est inversible.

\section{Constructeurs d'anneaux, sous-anneaux, sous-corps}

Comme pour les groupes, si $A, A'$ sont des anneaux et $X$ un ensemble non vide, $A\times A'$ a une structure d'anneau et $\mathfrak F(X, A)$ aussi

\begin{defprop}
    \begin{enumerate}
        \item On dira que $A'$ est un sous-anneau de $A$ si \begin{itemize}
            \item $A'\subset A$
            \item $1_A\in A$
            \item $\forall x, y\in A', x-y\in A'$ et $xy\in A'$
        \end{itemize}
    \item On dira que $K'$ est un sous-corps de $K$ si \begin{itemize}
        \item $K'\subset K$
        \item $1_K\in K$
        \item $\forall x, y\in K', x-y\in K'$ et si $y\neq 0$, $xy^{-1}\in K'$
    \end{itemize}
    \end{enumerate}
\end{defprop}

\begin{ex}~
    \begin{itemize}
        \item $\mathbb Z[\sqrt 2]$ est un sous-anneau de $\mathbb R$
        \item $\mathbb Q[\sqrt 2]$ est un sous-corps de $\mathbb R$ qu'on note $\mathbb Q(\sqrt 2)$
        \item $\mathbb Q[\omega]$ pour $\omega$ algébrique est un sous-corps de $\mathbb R$ (c'est un espace vectoriel de dimension finie et $\varphi: t\longmapsto tx$ est injective donc bijective car linéaire ce qui montre l'inversibilité)
    \end{itemize}
\end{ex}

\begin{rem}
    Si on fixe $n$ et on note \[
        S_n=\left\{\sum_{i=1}^nx_i^2\;, \qquad x_1, \cdots, x_n\in K\right\}
    \]
    alors on peut montrer que $S_n$ est stable par $\times$ si $n=2^k$
\end{rem}

\section{Idéal d'un anneau}

\begin{dfn}
    Un idéal\index{idéal} à gauche (resp. à droite) $I\subset A$ d'un anneau $A$ est un sous-groupe de $(A, +)$ tel que \[
        \forall a\in A, x\in I,\quad  ax\in I\quad \text{(resp. $xa\in I$)}
    \]
    Un idéal bilatère (ou simplement idéal) est un idéal à gauche et à droite.
\end{dfn}


\begin{rem}
    Pour un anneau commutatif $A$ et $a\in A$, $I=aA$ est un idéal.
\end{rem}

L'ensemble des idéaux d'un anneau commutatif est stable par somme et par intersection. De plus, si $I$ et $J$ sont des idéaux, alors \[
    IJ \defeq \left\{ \sum_{k\in K}i_kj_k, \quad K\subset \mathbb N \text{ fini }, i_k\in I, j_k\in J \right\}
\]
est aussi un idéal. Le radical $\sqrt I$ de $I$, définit par \[
    \sqrt I \defeq \left\{ x\in A, \quad \exists n\in\mathbb N, x^n\in I \right\},
\]
est également un idéal (c'est un sous-groupe par binôme de Newton).

\begin{dfn}[Hors-Programme]
    Un anneau $A$ est principal si tous ses idéaux sont principaux\index{idéal!principal}, c'est à dire de la forme $I=aA$. Un anneau est noéthérien\index{anneau!noéthérien} si toute suite croissante d'idéaux est stationnaire
\end{dfn}

\begin{exo}
    Décrire les idéaux de $\mathbb K^2$ pour un corps $\mathbb K$
\end{exo}

\section{Groupe des inversibles d'un anneau}

Pour un anneau $A$, on note $A^\star$ le groupe des inversibles\index{anneau!inversibles} de $A$.

\begin{ex}
    On note $A=\mathbb Z[\sqrt 2]$, on va décrire $\mathbb Z[\sqrt 2]^\star$.

    On note $G=\mathbb Z[\sqrt 2]^\star$ de sorte que $G_+=G\cap \mathbb R_+^\star$ est un groupe pour $\times$. Le logarithme est un morphisme injectif donc $\ln(G_+)$ est un sous-groupe de $\mathbb R$, donc dense si et seulement si $a=\inf (\ln(G_+)\cap \mathbb R_+^\star)=0$.

    On a $e^a=\inf(G_+\cap ]1; +\infty[)$ et \[
        1\leq e^a< 1+\sqrt 2 \implies \exists p,q \in \mathbb Z, 1\leq e^a\leq \underbrace{p+q\sqrt 2}_{\in G_+}< 1+\sqrt 2
    \]
    Si $p>0, q<0$ alors $1\leq p+q\sqrt 2\leq p-q\sqrt 2<p^2-2q^2=1$ (la dernière égalité vient de l'inversibilité de $p+\sqrt2 q$, cf. remarque). C'est absurde. On arrive à la même conclusion si $p<0, q>0$. On a donc $e^a=1+\sqrt 2$ et $G=\langle 1+\sqrt 2 \rangle$
\end{ex}

\begin{rem}
    Dans un anneau euclidien $A$ de stathme $v$, les éléments inversibles sont exactement les éléments qui minimisent le stathme. Si $u$ minimise le stathme, $1=qu+r$ avec $r=0$ ou $v(r)<v(u)$ donc $r=0$ et $1=qu$ donc $u$ inversible. Sinon, $v(u)>v(1)$ donc $v(u)\geq 2$ donc $\forall q, 2\leq v(uq)$ donc $u$ n'est pas inversible.

    Par exemple, dans $\mathbb Z$ euclidien pour le stathme $|\,\cdot\,|$, les inversibles sont $-1$ et $1$. Dans $\mathbb Q[X]$ euclidien pour le stathme $\deg$, les inversibles sont tous les polynômes de degré $0$ (les polynômes constants).
\end{rem}

\section{Morphismes d'anneaux}

\begin{dfn}
    Pour $A, A'$ des anneaux, $\varphi:A\to A'$ est un morphisme\index{anneau!morphisme} d'anneau si \begin{itemize}
        \item $f(1_A)=1_{A'}$
        \item $\forall x, y\in A, f(x+y)=f(x)+f(y)$
        \item $\forall x, y\in A, f(xy)=f(x)f(y)$
    \end{itemize}
\end{dfn}

Comme dans les groupes ou les espaces vectoriels (dans les e.v. les applications linéaires jouent le rôle de morphismes), les morphismes d'anneaux préservent la structure (l'image et la préimage d'un anneau par un morphisme est un anneau).

\section{Structure d'algèbre}

\begin{dfn}
    On note $(A, +, \times, \cdot)$ muni de deux lois internes $+$ et $\times$ et d'une loi externe $\cdot$ sur le corps $K$. On dira que $A$ est une $K$-algèbre\index{algèbre (structure d' -- )} si $(A, +,\times)$ est un anneau, $(A, +, \cdot)$ est un $K$-e.v, et \[
        \forall \lambda\in K, \forall x, y\in A, \lambda\cdot (x\times y)=(\lambda\cdot x)\times y=x\times (\lambda\cdot y)
    \]
\end{dfn}

\begin{rem}
    Un corps $K$ est toujours une $K$-algèbre, la multiplication servant de loi interne et externe.
\end{rem}
\endchapter


\chapter{Arithmétique}

\ifsolo
    ~

    \vspace{1cm}

    \begin{center}
        \textbf{\LARGE Arithmétique} \\[1em]
    \end{center}
    \tableofcontents
\else
    \minitoc
\fi
\thispagestyle{empty}

\ifsolo \newpage \setcounter{page}{1} \fi
\section{Arithmétique classique}

\subsection{Rappels}

On se place dans l'anneau $(\mathbb Z, +, \times)$, et on note $\mathbb P$ l'ensemble des nombres premiers.

\begin{thm}[Division Euclidienne]
    \index{division euclidienne} Pour $a, b\in\mathbb Z$, $b$ non nul, il existe des uniques entiers $q, r$ tels que \[
        a=qb+r\qquad \qquad \text{ et } 0\leq r<|b|
    \]
\end{thm}

\begin{rem}
    L'existence d'un tel $r$ montre que $\mathbb Z$ est euclidien pour le stathme $|\cdot|$
\end{rem}

\begin{thm}[Théorème fondamental de l'arithmétique]
    Tout entier relatif $n$ non nul s'écrit de manière\index{théorème fondamental de l'arithmétique} unique (à l'ordre près des facteurs et à une constante inversible près) comme produit de puissances de premiers: \[
        \forall n\in\mathbb Z^\star, \exists! p_1, \cdots, p_r\in\mathbb P, \exists! \alpha_1, \cdots, \alpha_r\in\mathbb N^\star, \exists \epsilon\in\mathcal U_{\mathbb Z}=\{-1, 1\}, \qquad n=\epsilon p_1^{\alpha_1}\cdots p_r^{\alpha_r}
    \]
\end{thm}

\begin{thmdef}[Bézout, PGCD]
    Pour $a, b\in\mathbb Z$, il existe un unique entier positif $d$ tel que $a\mathbb Z+b\mathbb Z=d\mathbb Z$. On appelle cet entier le PGCD de $a$ et $b$, et on le note $d=a\land b=(a, b)=\mathrm{PGCD}(a, b)$. Si $d=a\land b$ alors il existe $u, v\in \mathbb Z$ tels que $au+bv=d$ (c'est la  propriété de Bézout)\index{Bézout (théorème de -- )}
\end{thmdef}

\begin{rem}
    Cette définition du PGCD coïncide avec la définition usuelle.
\end{rem}

\begin{thmdef}[Identité de Bézout, Coprimalité]
    Deux entiers $a$ et $b$ sont dits premiers entre eux si les ensembles de leurs diviseurs stricts sont disjoints. C'est équivalent à \[
        \exists u, v\in\mathbb Z, \qquad au+bv=1
    \]
\end{thmdef}

\begin{thm}[Lemme de Gauss]
    \index{Gauss!lemme} Si $a, b, c\in \mathbb Z$ et $a\land c=1$ alors \[
        a\;|\;bc\implies a\;|\; b
    \]
\end{thm}

\subsection{Utilisation d'un Vandermonde}

On note $a_1, \cdots, a_n\in\mathbb Z$. On veut montrer que \[
    \prod_{1\leq i<j\leq n}(j-i)\;\Big|\; \prod_{1\leq i<j\leq n}(a_j-a_i)
\]

On introduit \[
    \Delta(X_1, \cdots, X_n)= \begin{vmatrix}
        \binom{X_1}0 & \cdots & \binom{X_1}{n-1} \\
        \vdots & \ddots & \vdots \\
        \binom{X_n}0 & \cdots & \binom{X_n}{n-1}
    \end{vmatrix}\qquad \text{ avec }\binom Xk=\frac{X(X-1)\cdots (X-k+1)}{k!}
\]
On a \[
    \binom X{n-1}=\frac{X^{n-1}}{(n-1)!}+\underbrace{c_{n-2}X^{n-2}+\cdots +c_0}_{\in \Vect \left( \binom X0, \cdots, \binom X{n-2} \right)}
\]
de sorte que dans le déterminant on peut remplacer la dernière colonne par \[
    \begin{pmatrix}
        \frac{X_1^{n-1}}{(n-1)!} \\
        \vdots \\
        \frac{X_n^{n-1}}{(n-1)!}
    \end{pmatrix}
\]
et en itérant sur les colonnes de droite à gauche, \[
    \Delta(X_1, \cdots, X_n)= \begin{vmatrix}
        \;\;\;1\;\; & \dfrac{X_1}1 & \cdots & \dfrac{X_1^{n-1}}{(n-1)!} \\
        \;\vdots & \ddots & \ddots & \vdots \\
        \;\vdots & \ddots & \ddots & \vdots \\
        \;\;\;1\;\; & \cdots & \cdots & \dfrac{X_n^{n-1}}{(n-1)!}
    \end{vmatrix}
\]
et \[
    \Delta(a_1, \cdots, a_n)=\frac1{1!2!\cdots (n-1)!}V(a_1, \cdots, a_n)=\frac1{\prod_{1\leq i<j\leq n}(j-i)}\prod_{1\leq i<j\leq n}(a_j-a_i)
\]
or $\Delta(a_1, \cdots, a_n)$ est bien un entier car les $\binom{a_i}k$ sont des entiers, ce qui conclut.

\subsection{Méthode de la descente infinie}

On note $p$ premier et pour $n\geq 3$, on considère l'équation $x^n+py^n=p^2z^n$.
\begin{itemize}
    \item $(0, 0, 0)$ est solution. On suppose qu'il y en a d'autres et on note $(x, y, z)\neq (0, 0, 0)$ une solution minimale pour $|x|+|y|+|z|$.
    \item On a $p\;|\;p^2z^n-py^n=x^n$ donc $p$ divise $x$: on écrit $x=px'$. On a \[
            p^{n-1}x'^n+y^n=pz^n
        \]
        donc $p$ divise $y$ qu'on écrit $y=py'$ et donc $p$ divise $z$ qui s'écrit $z=pz'$. On a ainsi trouvé une nouvelle solution $(x', y', z')$ strictement inférieure à la précédente, ce qui est absurde. Ainsi, la seule solution est la solution nulle.
\end{itemize}

\subsection{Formule de Legendre}

\begin{dfn}
    \index{valuation $p$-adique} Pour tout premier $p$, on définit la \textbf{valuation $\bm p$-adique} $v_p$ par : \[
        v_p: n\in\mathbb Z^\star \longmapsto \max\{k\in\mathbb N, \quad p^k\;|\; n\}
    \]
\end{dfn}

On va montrer la formule de Legendre\index{Legendre!formule}: \[
    v_p(n!)=\sum_{i=1}^{+\infty}\floor{\frac n{p^i}}
\]
Pour cela, on remarque \[
    v_p(n!)=\sum_{m=1}^nv_p(m)=\sum_{m=1}^n\sum_{k=1}^{v_p(m)}1=\sum_{k=1}^{+\infty}\sum_{\substack{k\leq v_p(m)\\ 1\leq m\leq n}}1=\sum_{k=1}^{+\infty}\sum_{d\leq m=p^ku\leq n}1=\sum_{k=1}^{+\infty}\floor{\frac n{p^k}}
\]

\begin{exo}[X]
     Pour $m, n\in\mathbb N$, montrer \[
        \frac{(2m)!(2n)!}{(m+n)!n!m!}\in\mathbb N
     \]
\end{exo}

\subsection{Coefficients du binôme}

\todo{Retrouver ça}

\section{L'anneau $\mathbb Z_n$}

Soit $A$ un anneau commutatif et $I$ un idéal de $A$. On définit la relation d'équivalence suivante: $x\mathcal Ry\iff x-y\in I$. On note $A/I$ l'ensemble des classes d'équivalences: c'est un anneau commutatif avec les lois $\bar{x+y}=\bar x+\bar y$ et $\bar{x\times y}=\bar x\times \bar y$.

Pour $A=\mathbb Z$, $I=n\mathbb Z$, il y a $n$ classes d'équivalence dans $\mathbb Z_n\defeq\mathbb Z/n\mathbb Z$, ce sont les $\bar 0, \cdots , \bar{n-1}$.

\section{Groupe $\mathbb Z_n^\star$ des inversibles de $\mathbb Z_n$}

\begin{prop}
    \begin{itemize}
        \item La classe de $a\in\mathbb Z$ est inversible dans $\mathbb Z_n$ si et seulement si $a\land n=1$
        \item $\#\mathbb Z_n^\star=\varphi(n)$
        \item (Euler\index{Euler!théorème (arithmétique modulaire)}) Pour $\bar a$ inversible, $\bar a^{\varphi(n)}\equiv \bar 1\pmod n$
    \end{itemize}
\end{prop}

\begin{proof}~
    \begin{enumerate}
        \item $\bar a\in\mathbb Z_n^\star\iff \exists \bar u\in\mathbb Z_n, \bar a\times \bar u\equiv 1\pmod n\iff \exists u, v\in\mathbb Z, au=1+nv \iff a\land n=1$
        \item C'est évident vu 1.
        \item $\ord(\bar a)\;|\;\#\mathbb Z_n^\star=\varphi(n)$
    \end{enumerate}
\end{proof}

\begin{cor}[Petit théorème de Fermat]
    \index{Fermat!petit théorème} Pour tout $a\in\mathbb Z, \quad \bar a^p\equiv \bar a\pmod p$ pour $p$ premier
\end{cor}

\section{Caractérisations des corps}

\begin{prop}
    $\mathbb Z_n$ est un corps si et seulement si $n$ est premier
\end{prop}

\begin{proof}
    $\varphi(n)=n-1\iff n\in \mathbb P$
\end{proof}

\subsection{Théorème de Wilson}

\begin{res}[Théorème de Wilson]
    $p\geq 3$ est premier si et seulement si $(p-1)!\equiv -1\pmod p$\index{Wilson (théorème de -- )}
\end{res}

\begin{proof}~
    \begin{itemize}
        \item $(\implies)$ Dans $\mathbb Z_p$, $x^2\equiv 1$ n'a que deux solutions: $1$ et $-1$ (car on est dans un corps donc un anneau intègre). Ce sont donc les deux seuls éléments qui sont leurs propres inverses, les autres sont couplés deux à deux dans l'écriture de $(p-1)!$. On a donc \[
                (p-1)!\equiv 1\times (2\times \cdots \times (p-2))\times (p-1)\equiv 1\times (p-1)\equiv -1 \pmod p
            \]
        \item $(\impliedby)$ On note $n\geq 3$ non premier, il s'écrit donc $n=ab$ avec $a\land b=1$ et $a, b\geq 2$. Si $a\neq b$ alors $ab\;|\;(n-1)!$. Sinon, $n=a^2$ et $a\geq 2$ donc $a^2\;|\;(n-1)!$ car $a$ et $2a$ apparaissent dans l'écriture de la factorielle. Dans les deux cas, $(n-1)!\equiv 0\pmod n$
    \end{itemize}
\end{proof}

\subsection{RSA}

On va voir (rapidement) le principe du chiffrement RSA. On choisit $p, q$ des nombres premiers distincts, et on note $n=pq$. On note $a, b$ tels que $ab\equiv 1\pmod{\varphi(n)}$. Les nombres $n$ et $a$ sont publics, et $b$ (qu'on appelle clé privée) n'est connu que de Bob, qui veut recevoir un message chiffré d'alice

Alice veut transmettre $x\in\mathbb Z_n^\star$. Elle calcule $x^n\equiv y\pmod n$, et $y$ est le message chiffré.

Bob calcule $y^b\equiv x^{ab}\equiv x\times x^{\varphi(n)k}\equiv x\pmod n$

\subsection{$\mathbb Z_{p^\alpha}^\star$ pour $p$ premier impair}

\subsubsection{$\mathbb Z_p^\star$ est cyclique}

On va d'abord montrer les résultats suivants, vrais dans tous les groupes abéliens finis: \begin{itemize}
    \item L'ensemble des ordres est stable par ppcm
    \item L'exposant d'un groupe est l'ordre d'un élément
\end{itemize}
On utilisera ensuite la structure de corps de $\mathbb Z/p\mathbb Z$ pour conclure.

Soit $(G, \cdot)$ un groupe abélien fini. On va montrer que l'ensemble des ordres est stable par ppcm. Soient $a, b$ deux éléments d'ordres respectifs $m, n$. On cherche un élément d'ordre $m\lor n$. On note $d=m\land n$ de telle sorte que $m=dm'$, $n=dn'$ et $m\land n'=1$. De cette manière, on a $m\lor n=mn'$
\[
    \ord \left(a\right)=m\qquad \qquad \ord\left(b^d\right)=n'
\]
et donc $c=ab^d$ est un élément d'ordre $mn'=m\lor n$, ce qui conclut.

On définit l'exposant\index{groupe!exposant} $\lambda$ du groupe $G$ comme l'ordre maximal des éléments de $G$. L'ensemble des ordres étant stable par ppcm, et le ppcm de tous les ordres majorant l'ensemble, le maximum des ordres $\lambda$ est le ppcm des ordres des éléments. Puis, la stabilité garantit qu'il existe un élément d'ordre $\lambda$.

On va maintenant montrer que $\mathbb Z_p^\star$ est cyclique. $\mathbb Z_p$ est un corps (déjà vu). Ainsi, un polynôme de $\mathbb Z_p[X]$ admet au plus autant de racines que son degré. On pose $P=X^\lambda-1$. Ce polynôme admet au moins $p-1$ racines distinctes (les éléments non nuls du corps) par définition de $\lambda$. On a donc $\lambda\geq p-1$. Puis, $\lambda$ est un ordre donc divise $\#\mathbb Z_p^\star=p-1$ (Lagrange). Ainsi, $\lambda\mid p-1$ et $\lambda=p-1$.

L'exposant du groupe vaut le cardinal du groupe, ce groupe est donc cyclique (un élément d'ordre $\lambda$ est automatiquement un générateur du groupe).

\subsection{$\mathbb Z_{p^\alpha}^\star$ est cyclique}

On va d'abord montrer par récurrence \[
    \forall n\in\mathbb N, \quad (1+p)^{p^n}\equiv 1+p^{n+1}\pmod {p^{n+2}}.
\]

\begin{itemize}
    \item Le cas $n=0$ est direct.
    \item On suppose le résultat vrai au rang $n$: \[
        (1+p)^{p^n}=1+p^{n+1}+p^{n+2}q.
    \]
    On a alors \begin{align*}
        (1+p)^{p^{n+1}}&\equiv(1+p^{n+1}+p^{n+2}q)^p\\
                       &\equiv 1+\sum_{k=1}^{p}\binom pk p^{(n+1)k}(1+pq)^k\\ 
                       &\equiv 1 + p^{n+1}(1+pq)p \\
                       &\equiv 1+p^{n+2}&\pmod{p^{n+3}}.
    \end{align*}
    Donc l'égalité est vraie pour tout entier naturel $n$.
\end{itemize}
Il est suffisant de trouver un élément de $(\mathbb Z/p^\alpha\mathbb Z)^\star$ d'ordre $\varphi(p^\alpha)=(p-1)p^{\alpha-1}$ (cette égalité s'obtient via la définition de l'indicatrice d'Euler), et vu la stabilité des ordres par ppcm vue précédemment, il est suffisant de trouver un élément d'ordre $p-1$ et un élément d'ordre $p^{\alpha-1}$.

On va d'abord exhiber un élément d'ordre $p^{\alpha-1}$. On a 
\[
(1+p)^{p^{\alpha-2}} 
\equiv 1+p^{\alpha-1}\pmod {p^\alpha}    \]
et \[
    (1+p)^{p^{\alpha-1}}\equiv 1\pmod {p^\alpha}.
\]
Donc $\ord(1+p)\mid p^{\alpha-1}$ mais $\ord(1+p)\nmid p^{\alpha-2}$ donc $\ord(1+p)=p^{\alpha-1}$.

Il reste à trouver un élément d'ordre $p-1$. On sait que $\mathbb F_p^\star$ est cyclique, on note $g$ un générateur. On a:
\[
    \forall k<p-1, \quad p\;\nmid\;g^k-1\quad \text{ie}\quad p^\alpha\;\nmid\; g^k-1
\]
et $p\mid g^{p-1}-1$ donc $\ord g=k(p-1)$ est $g^k$ et d'ordre $p-1$.

\section{Théorème chinois}

\begin{thm}
    \Hyp $m, n\in\mathbb N^\star$, $m\land n=1$
    \begin{concenum}
    \item La classe de $\bar x\in\mathbb Z_{mn}$ dans $\mathbb Z_n$ ne dépend pas du choix de représentant dans $\mathbb Z_{mn}$
    \item L'application \[
            \begin{matrix}
                \varphi: & \mathbb Z_{m, n} & \longrightarrow & \mathbb Z_n\times \mathbb Z_m \\
                         & \bar x &\longmapsto & (\bar x, \bar x)
            \end{matrix}
        \]
        est un isomorphisme d'anneaux
    \item $\varphi$ est une bijection de $\mathbb Z_{m, n}^\star$ dans $\mathbb Z_n^\star \times \mathbb Z_m^\star$
    \item Pour $a, b\in\mathbb Z$ fixés, le système \[
            \begin{cases}
                x=a\pmod n\\
                x=b\pmod m
            \end{cases}
        \]
        admet une unique solution modulo $mn$
    \end{concenum}
\end{thm}

\begin{proof}~
    \begin{enumerate}
        \item $a, b$ de représentants de $\bar x$ dans $\mathbb Z_{m, n}$. $a\equiv b\pmod mn\implies a\equiv b\pmod n$ donc $\bar a=\bar b$ dans $\mathbb Z_n$
        \item C'est bien un morphisme, il est injectif car $\Ker\varphi=\{\bar 0\}$ donc c'est un isomorphisme car les ensembles d'antécédents et d'images sont équipotents finis.
        \item C'est clair ($a\land n\neq 1\implies a\land mn\neq 1$ d'où on déduit l'injectivité).
        \item C'est 2.
    \end{enumerate}
\end{proof}

\begin{rem}
    Le point 3. donne la multiplicativité de l'indicatrice d'euler: pour $m, n$ premiers entre eux, $\varphi(mn)=\varphi(m)\varphi(n)$
\end{rem}

\begin{rem}
    On peut généraliser le résultat à $n_1, \cdots, n_r$ deux à deux premiers entre eux
\end{rem}

\begin{rem}
    Si $un+vm=1$ alors $x=avm+bun$ est une solution
\end{rem}

\subsection{Puissances parfaites}

On va montrer qu'il existe $n$ entiers consécutifs qui ne sont pas des puissances parfaites.

On note $p_1, \cdots, p_n$ des premiers deux à deux distincts. Le système \[
    \begin{cases}
        x\equiv p_1-1\pmod {p_1^2}\\
        \hspace{1cm}\vdots\\
        x\equiv p_n-n\pmod{p_n^2}
    \end{cases}
\]
possède une solution $x\in\mathbb N$ et $v_{p_i}(x+i)=1$ donc $x+i$ n'est pas une puissance parfaite (car $x+i\neq p_i$ quitte à translater $x$ de $p_1^2\cdots p_n^2$)

\subsection{Points visibles}

\begin{center}
    \includegraphics{src/figures/arithmetique-points-visibles.pdf}
\end{center}

On considère l'ensemble des points (qu'on appelle \emph{points visibles}) de $\mathbb Z^2$ tels que pour un point $(x, y)$, l'ensemble $\{(1-t)(0, 0)+t(x, y), \quad t\in ]0, 1[\}$ ne contient pas de point de $\mathbb Z^2$. On veut montrer qu'il existe des rectangles pleins arbitrairement grands dans cet ensemble, c'est à dire que pour tout $n\in\mathbb N$, il existe $(a, b)$ tel que $\forall i, j\leq n, (a+i, a+j)$ est invisible.

On vérifie facilement que cet ensemble est l'ensemble $\{(a, b)\in\mathbb Z^2, \quad a\land b=1\}$. On note $A$ une matrice $k\times k$ dont les coefficients sont des nombres premiers deux à deux distincts. On note $m_i$ le produit de la $i$-ième ligne, $M_i$ le produit de la $i$-ième colonne. Le théorème chinois donne l'existence de $a$, $b$ tels que \[
    \begin{cases}
        a\equiv -1\pmod {m_1} \\
        \hspace{.4cm}\vdots \\
        a\equiv -k\pmod{m_k}
    \end{cases}
    \qquad \begin{cases}
        b\equiv -1\pmod {M_1}\\
        \hspace{.4cm}\vdots \\
        b\equiv -k\pmod{M_k}
    \end{cases}
\]
de sorte que $(a+i)\land (b+j)$ est divisible par $m_i\land M_j\neq 1$.

\subsection{Les congruences de Lucas}

On note $p\in\mathbb P$ et \[
    n\defeq\sum_{k=0}^Na_kp^k\geq \sum_{k=0}^Nb_kp^k\defeq m
\]
On va montrer que \[
    \binom nm\equiv \binom{a_0}{b_0}\cdots \binom{a_N}{b_N}\pmod p
\]

On a $(1+X)^p=1+X^p$ dans $\mathbb Z_p[X]$. Puis, \[
    (1+X)^n=\prod_{i=0}^N((1+X)^p)^{a_i}=\prod_{i=0}^n\sum_{k=0}^{a_i}\binom{a_i}{k}X^{kp^i}
\]
et il suffit d'identifier le coefficient de $X^m$

\section{Principe d'inclusion-exclusion}

Pour $E$ fini, $A_1, \cdots, A_N\in\mathcal P(E)$, \[
    \#A_1\cup\cdots \cup A_N=\sum_{i_1=1}^N\#A_{i_1}-\sum_{i_1<i_2}\#A_{i_1}\cap A_{i_2}+\cdots +(-1)^{N-1}\sum_{i_1<\cdots <i_N}\#A_1\cap \cdots \cap A_N
\]

Pour le voir, on remarque: \[
    \sum_{x\in E}\mathbbm 1_{A}(x)=\#A \qquad\quad \mathbbm 1_{A\cap B}=\mathbbm 1_A \cdot \mathbbm 1_B \qquad \quad \mathbbm 1_{A^c}=1-\mathbbm 1_A
\]
et \[
    (A_1\cup \cdots \cup A_N)^c=A_1^c\cap \cdots \cap A_N^c
\]
donc \[
    \mathbbm1_{(A_1\cup \cdots \cup A_N)^c}=\prod_{i=1}^N(1-\mathbbm 1_{A_i^c})
\]
puis on développe le produit et on somme la fonction prise sur les éléments de $E$


\begin{ex}[Application aux permutations sans points fixes]
    On note $d_n$ le nombre de permutations de $\mathfrak S_n$ sans points fixes.

    Il y a $\displaystyle \binom ni d_{n-i}$ permutations qui ont $i$ points fixes donc \[
        n!=\sum_{i=0}^n\binom nid_{n-i}=\sum_{i=0}^n\binom nid_i.
    \]
    Matriciellement, \[
        \begin{pmatrix}
            n!\\[1em]
            \vdots\\[1em]
            \vdots\\[1em]
            0!
            \end{pmatrix}= \underbrace{\begin{pmatrix}
        \displaystyle\binom nn & \cdots & \cdots & \displaystyle\binom n0 \\[1em]
        0 & \displaystyle\binom{n-1}{n-1} & \ddots & \vdots \\[1em]
        \vdots & \ddots & \ddots & \vdots \\[1em]
        0 & \cdots & 0 & \displaystyle\binom 00
    \end{pmatrix}}_{\defeq A_n\in \mathrm{GL}_n(\mathbb R)} \begin{pmatrix}
        d_n\\[1em]
        \vdots \\[1em]
        d_1\\[1em]
        d_0
        \end{pmatrix}
    \]
    On a $A_n^\top=\mathcal {M}_{\text{base canonique}}(f)$ avec $f:P\in\mathbb R_n[X]\longmapsto P(X+1)$. Puis, $(A_n^\top)^{-1}=(A_n^{-1})^\top=\mathcal M(f^{-1})$ or $f^{-1}:P\longmapsto P(X-1)$ d'où \[
        (A_n^{-1})^\top = \begin{pmatrix}
            \displaystyle\binom nn & 0 & \cdots & 0 \\[1em]
            \displaystyle -\binom n{n-1} & \ddots & \ddots & \vdots \\[1em]
            \vdots &  & \ddots & 0 \\[1em]
            (-1)^n\displaystyle\binom n0 & \;\;\;\cdots\;\;\; & \;\;\;\cdots\;\;\; & \displaystyle\binom 00
        \end{pmatrix}
    \]
    d'où \[
        d_n=\sum_{i=0}^n\binom n{n-i}(n-i)!(-1)^i=\sum_{i=0}^n\frac{(-1)^in!}{i!}\sim \frac{n!}e
    \]
\end{ex}

\section{Principe des tiroirs}


\chapter{Réduction des endomorphismes}

\ifsolo
    ~

    \vspace{1cm}

    \begin{center}
        \textbf{\LARGE Réduction des endomorphismes} \\[1em]
    \end{center}
    \tableofcontents
\else
    \minitoc
\fi
\thispagestyle{empty}

\ifsolo \newpage \setcounter{page}{1} \fi

Dans tout le chapitre, on note $\mathbb K$ un corps.

\section{Rappels sur les polynômes}

\begin{thm}
    \Hyp $A, B\in \mathbb K[X]$, $B\neq 0$
    \Conc Il existe un unique couple de polynômes $(Q, R)$ dans $\mathbb K[X]$ tels que \[
        A=QB+R
    \]
    et $\deg R<\deg B$
\end{thm}

\begin{thmdef}
    $a\in\mathbb K$ est une racine de $P$ ssi \[
    P(a)=0\iff (X-a)\;|\;P
\]
et $a$ est racine de multiplicité $\alpha \geq 1$ ssi \[
    \begin{cases}
        P(a)=P'(a)=\cdots=P^{(\alpha - 1)}(a)=0 \\
        P^{(\alpha)}(a)\neq 0
    \end{cases}
    \iff \begin{cases}
        (X-a)^\alpha \; |\; P\\
        (X-a)^{\alpha+1}\;\not|\;P
    \end{cases}
\]
Si $a$ n'est pas racine de $P$, on convient que $a$ est racine de multiplicité $0$.
\end{thmdef}

\begin{thm}[Interpolation de Lagrange]
    \index{Lagrange!interpolation} Pour $a_1, \cdots, a_n\in\mathbb R$ des abscisses distinctes et $b_1, \cdots, b_n\in\mathbb R$ des ordonnées distinctes, il existe un unique polynôme de degré au plus $n-1$ tel que \[
        \forall i\in\llbracket 1, n\rrbracket, \qquad P(a_i)=b_i
    \]
\end{thm}
\begin{proof}
    \[
        \varphi:P\in\mathbb K_{n-1}[X]\longmapsto (P(a_1), \cdots, P(a_n))\in\mathbb R^n
    \]
    est un isomorphisme.
\end{proof}

\section{Idéaux de {$\mathbb K[X]$} et arithmétique}

\begin{thm}
    $\mathbb K[X]$ est principal
\end{thm}
\begin{proof}
    C'est un anneau euclidien donc principal (cf. chapitre \textbf{Algèbre générale})
\end{proof}

\begin{thmdef}
    \begin{enumerate}
        \item Il existe un unique polynôme unitaire $Q$ tel que $P\mathbb K[X]+Q\mathbb K[X]=Q\mathbb K[X]$. C'est le PGCD de $A$ et $B$. Deux polynômes sont premiers entre eux si leur PGCD vaut $1$.
        \item Il existe un unique polynôme unitaire $Q$ tel que $P\mathbb K[X]\cap Q\mathbb K[X]=Q\mathbb K[X]$. C'est le PPCM de $A$ et $B$.
        \item (Gauss) \index{Gauss!lemme (polynômial)} Si $A$ et $B$ sont premiers entre eux et $A\;|\;BC$ alors $A\;|\;C$
        \item (Bézout) \index{Bézout (théorème de -- )} $A$ et $B$ sont premiers entre eux si et seulement si il existe $U, V$ tels que $AU+BV=1$.
    \end{enumerate}
\end{thmdef}

\begin{ex}
    $(X^n-1)\land (X^m-1)=X^{n\land m}-1$
\end{ex}

\section{Décomposition en facteurs irréductibles}

\begin{dfn}
    Un anneau $A$ est factoriel si \begin{enumerate}
        \item Pour tout $a\in A$, il existe des éléments irréductibles $p_1, \cdots, p_n$ tels que $a=p_1\cdots p_n$
        \item Cette décomposition est unique à l'ordre des facteurs près et à association près (deux éléments $a$ et $b$ sont associés si il existe un inversible $u$ tel que $a=ub$)
    \end{enumerate}
\end{dfn}

\begin{thm}
    $\mathbb K[X]$ est un anneau factoriel
\end{thm}

\begin{ex}
    Dans $\mathbb C[X]$, \begin{itemize}
        \item \[
                X^n-1=\prod_{\omega\in\mathbb U_n}(X-\omega)
            \]
        \item \[
                X^{2m}-1=(X+1)(X-1)\prod_{k=1}^{n-1} \left( X-\exp \left( \frac{ik\pi}n \right) \right) \left( X-\exp \left( -\frac{ik\pi}n \right) \right)
            \]
    \end{itemize}
\end{ex}

\begin{exo} Montrer que
    \begin{itemize}
        \item $P(\mathbb R)\subset \mathbb R_+ \iff \exists A, B\in\mathbb R[X], P=A^2+B^2$
        \item $P(\mathbb R_+)\subset \mathbb R_+\iff \exists A, B\in\mathbb R[X], P=A^2+XB^2$
    \end{itemize}
\end{exo}

\begin{proof}[Résolution]
    Éléments de résolution pour le premier point: l'ensemble des $A^2+B^2$ est stable par produit, et en décomposant $P\geq 0$ en facteurs irréductibles, chacun des facteurs est somme de deux carrés.
\end{proof}

\begin{exo}
    $P$ unitaire de degré $n$ dans $\mathbb R[X]$. Montrer \[
        P\text{ scindé dans }\mathbb R[X] \iff \forall z\in\mathbb C, \quad |P(z)|\geq |\mathfrak{Im}(z)|^n
    \]
\end{exo}

\begin{proof}~
    \begin{itemize}
        \item $(\impliedby)$ $P$ scindé dans $\mathbb C$ de racines $(a_i)$. $0=|P(a_i)|\geq |\mathfrak{Im}(z)|^n$ donc $a_i\in\mathbb R$
        \item $(\implies)$ \[
                \left| P(x+iy) \right|= \left| \prod_{j=1}^n(x-a_j+iy) \right|=\prod_{j=1}^n|x-a_j+iy|\geq \|\mathfrak{Im}(z)|^n
            \]
    \end{itemize}
\end{proof}


\section{Polynômes d'endomorphismes}

\begin{defprop}
    \Hyp On note $E$ un $\mathbb K$-e.v., $u\in\mathcal L(E)$ et \[
        P=\sum_{k=0}^{+\infty}a_kX^k\in\mathbb K[X]
    \]
    \begin{concenum}
    \item ~ \[
            P(u)\defeq\sum_{k=0}^{+\infty} a_ku^k
        \]
    \item $\varphi: P\in\mathbb K[X]\longmapsto P(u)\in\mathcal L(E)$ est un morphisme d'algèbres.
    \end{concenum}
\end{defprop}

\begin{proof}
    On justifie $(PQ)(u)=P(u)\circ Q(u)$. On considère \[
        \psi: (P, Q)\in\mathbb K[X]^2\longmapsto (PQ)(u)-P(u)\circ Q(u)\in\mathbb L(u)
    \]
    \begin{itemize}
        \item $\psi$ est bilinéaire
        \item À $Q$ fixé, $P\longmapsto \psi(P, Q)$ AL nulle ssi $\forall j\in\mathbb N, \psi(X^j, Q)=0$
        \item Pour chaque $j$, $Q\longmapsto \psi(X^j, Q)$ nulle ssi $\forall k\in\mathbb N, \psi(X^j, X^k)=0$ qui est vrai.
    \end{itemize}
\end{proof}

\begin{dfn}
    Un polynôme $P$ est dit \textbf{polynôme annulateur\index{polynôme!annulateur}} de $u\in\mathcal L(E)$ si $P(u)$ est l'application nulle
\end{dfn}

\begin{ex}
    $p$ est un projecteur si et seulement si $X^2-X$ est annulateur. $s$ est une symétrie si et seulement si $X^2-1$ est annulateur
\end{ex}

\begin{defprop}
    L'ensemble des polynômes annulateurs d'un endomorphisme $u$ est un idéal. Il possède un unique générateur unitaire, qu'on appelle \textbf{polynôme minimal\index{polynôme!minimal}}, qu'on notera $\mu_u$ ou $\Pi_u$
\end{defprop}

\begin{proof}
    $(\id, u, \cdots, u^{n^2})$ famille liée de $\mathcal L(E)$ donc l'idéal n'est pas trivial
\end{proof}

\begin{prop}
    $\mathbb K[u]\defeq \{P(u), P\in\mathbb K[X]\}$ est une $\mathbb K$-algèbre de dimension $\deg \mu_u$ de base $(1, u, \cdots, u^{\deg \mu_u-1})$
\end{prop}

\begin{ex}
    $P=X^2-3X+2$ polynôme minimal \[
        A= \begin{pmatrix}
            -2 & -6 \\ 2 & 5
        \end{pmatrix}
    \]
    Puis $X^n=\mu_A\times Q+\alpha_n X+\beta_n$. En injectant $1$ et $2$ (les racines de $\mu_A$) on peut exprimer $\alpha_n$ et $\beta_n$ et $A^n=\alpha_nA+\beta_n$
\end{ex}

\section{Théorème de décomposition des noyaux}

\begin{thm}
    \Hyp $E$ un $\mathbb K$-e.v., $u\in\mathcal L(E)$
    \begin{concenum}
    \item Si $P, Q\in\mathbb K[X]$ sont premiers entre eux, alors $\Ker (PQ)(u)=\Ker P(u)\oplus \Ker Q(u)$. En particulier, si $PQ$ est annulateur alors $\Ker(PQ)(u)=E$
    \item Si $P_1, \cdots, P_q\in\mathbb K[X]$ deux à deux premiers entre eux et si $P=P_1\cdots P_q$ alors \[
            \Ker P(u)=\Ker P_1(u)\oplus\cdots\oplus \Ker P_q(u)
        \]
    \end{concenum}
\end{thm}


% \chapter{Analyse réelle et fonctions convexes}

% \chapter{Espaces vectoriels normés}

% \chapter{Limite -- Continuité -- Fonctions vectorielles}

% \chapter{Compacité -- Connexité}

% \chapter{Séries entières}

% \chapter{Variables aléatoires discrètes}

% \chapter{Espérance -- Variance}

% \chapter{Les intégrales à paramètres}

% \chapter{Les équations différentielles}

% \chapter{Espaces préhilbertiens réels}

% \chapter{Fonctions de plusieurs variables}

\chapter{Calcul différentiel}

\ifsolo
~

\vspace{1cm}

\begin{center}
\textbf{\LARGE Calcul différentiel} \\[1em]
\end{center}
\tableofcontents
\else
\minitoc
\fi
\thispagestyle{empty}

\ifsolo \newpage \setcounter{page}{1} \fi
\section{Notions, Définitions}

$E$ est un $\mathbb R$-e.v de dimension $n$ (donc isomorphe à $\mathbb R^n$), $F$ un sev de $E$ et $\Omega$ un ouvert de $E$.

\begin{rem}
Pour $a\in\Omega, v\in E, \exists \epsilon>0, \forall t\in ]-\epsilon, \epsilon[, a+tv\in\Omega$ (par définition d'un ouvert)
\end{rem}

\begin{dfn}
On note $f:\Omega \longrightarrow F$. On dira que $f$ possède une dérivée en $a\in\Omega$ selon $v\in E$ si \[
    \lim_{\substack{t\to 0\\ t\neq 0}}\frac{f(a+tv)-f(a)}t
\] existe. On note cette limite $D_vf(a)$.

Si $\mathcal B=(e_1,\cdots, e_n)$ base de $E$ et $f$ a une dérivée en $a\in\Omega$ selon $e_{i_0}$ alors on appelle $i_0$-ième dérivée partielle $D_{e_{i_0}}f(a)$, notée \[
    \frac{\partial f}{\partial x_{i_0}}(a)\qquad \text{ ou }\qquad \partial_{i_0}f(a)
\]
\end{dfn}

\begin{rem}
Pour $\mathcal B=(e_1, \cdots, e_n)$ on note $x\underset{\mathcal B}{\longleftrightarrow}(x_1, \cdot, x_n)$ et on a $f(x)=f(x_1e_1+\cdots + x_ne_n)$ qu'on note abusivement $f(x_1, \cdots, x_n)$. \[
    \frac{f(x+te_1)-f(x)}{t}=\frac{f(x_1+t, x_2, \cdots, x_n)-f(x_1, \cdots, x_n)}{t}
\]
Pour calculer $\partial_if(a)$, on fait comme si les autres variables étaient constantes et on dérive par rapport à $x_i$.
\end{rem}

\begin{rem}
$D_v$ est linéaire
\end{rem}

\begin{rem}
Avoir une dérivée en $a$ n'implique pas la continuité en $a$. Contre-exemple: \[
    f(x, y)=\begin{cases}
    \frac{y^2}x &\text{ si }x\neq 0\\ y & \text{ sinon}
    \end{cases}
\]
\end{rem}

\section{Différentielle}

\begin{defprop}
\Hyp $f:\Omega\longrightarrow F$, $a\in \Omega$. \begin{concenum}
    \item On dira que $f$ est différentiable en $a$ s'il existe $\varphi\in\mathcal L(E, F), \epsilon>0$ tels que \[
        \forall h\in\mathcal B_o(0, \epsilon), f(a+h)=f(a)+\varphi(h)+o_0(h)
    \]
    Dans ce cas $\varphi$ est unique et se note $\diff f_a$
    \item Si $f$ est différentiable en $a$ alors $f$ est continue en $a$
    \item Si $f$ est différentiable en $a$ alors $f$ a une dérivée en $a$ selon tout $v$ et \[
        D_vf(a)=df_a(v)
    \]
\end{concenum}
\end{defprop}

\begin{proof} ~
\begin{enumerate}
    \item Si $\varphi, \psi$ conviennent, $\varphi(tv)-\psi(tv)=o_0(tv)$ ie \[
        \varphi(v)-\psi(v)=o_0(v)\xrightarrow[t\to 0]{}0
    \]
    d'où $\varphi\equiv \psi$.
    \item $f(a+h)=f(a)+\varphi(h)+o_0(h)\xrightarrow[h\to 0]{}f(a)$
    \item \[
        \frac{f(a+tv)-f(a)}{t}=\diff f_a(v)+o_0(v)\xrightarrow[\substack{t\to 0\\t\neq 0}]{}\diff f_a(v)
    \]
\end{enumerate}
\end{proof}

\begin{thmdef}
~ \begin{enumerate}
    \item Soit $f\in\mathcal L(E, F)$ et $a\in E$. La fonction $f$ est différentiable en $a$ et $df_a=f$.
    \item Si $\mathcal B=(e_1, \cdots, e_n)$ est une base de $E$ et si $\pi_i$ est l'application $i$-ième coordonnée alors $\forall a\in E, \diff \pi_i=\pi_i$ se dépend pas de $a$. On note cette application $\diff x_i$.
    \item Si $\mathcal B$ est une base de $E$ et $f$ est différentiable en $a$ alors \[
        \diff f_a=\sum_{i=1}^n\frac{\partial f}{\partial x_i}(a)\diff x_i
    \]
\end{enumerate}
\end{thmdef}

\begin{proof} ~
\begin{enumerate}
    \item $f(a+h)=f(a)+f(h)$
    \item Ok
    \item $h=h_1e_1+\cdots +h_ne_n$, \[
        \diff f_a(h)=h_1\diff f_a(e_1)+\cdots +h_n\diff f_a(e_n)=\sum_{i=1}^n\frac{\partial f}{\partial x_i}(a)\diff x_i
    \]
\end{enumerate}
\end{proof}

\begin{rem}
Pour une fonction $f:\mathbb R\longrightarrow F$, $f'(a)=\diff f_a(1)$
\end{rem}

\section{Calculs de différentielles}

\subsection{Quelques exemples}
\begin{ex} Étude en $0$ de
\[
    f(x, y)=\begin{cases}
    \frac{y^2}x &\text{ si }x\neq 0\\ y & \text{ sinon}
    \end{cases}
\]
\end{ex}

\[
\frac{\partial f}{\partial x}(0, 0)=0 \qquad\qquad \frac{\partial f}{\partial y}(0, 0)=f(0, 1)=1
\]
donc \[
\frac{\partial f}{\partial x}(0, 0) \diff x+ \frac{\partial f}{\partial y}(0, 0)\diff y
\]
existe mais $f$ est non différentiable car non $\mathcal C^0$

\begin{ex} Calcul de la différentielle de
{$f:A\in\mathcal M_n(\mathbb R)\longmapsto A^2$}
\end{ex}

\[
    f(A+H)=f(A)+AH+HA+H^2
\]
et pour une norme sous-multiplicative, $\|H^2\|\leq \|H\|^2=o_0(H)$ donc \[
    \diff f_A(H)=AH+HA
\]

\begin{ex} Calcul de la différentielle de
{$f:M\in GL_n(\mathbb R)\longmapsto M^{-1}$}
\end{ex}

On fixe $M\in GL_n(\mathbb R)$. Alors $\exists \epsilon >0 / \mathcal B_o(M, \epsilon)\subset GL_n(\mathbb R)$. Pour $H$ dans $\mathcal B_o(0, \epsilon)$, \[
    M+H\in\mathcal B_o(M, \epsilon)
\]
et \[
    \left((M+H)^{-1}-M^{-1}\right)(M+H)=I_n-M^{-1}(M+H)=-M^{-1}H
    \] donc
    \begin{align*}
    (M+H)^{-1}-M^{-1}&=-M^{-1}H(M+H)^{-1}\\&=-M^{-1}HM^{-1}+M^{-1}Ho_0(1)\\&=\underbrace{-M^{-1}HM^{-1}}_{\diff f_M(H)}+o_0(H)
\end{align*}

\begin{exo} Calcul de la différentielle de
{$f:A\in\mathcal M_n(\mathbb R)\longmapsto A^\intercal A$}
\end{exo}

\needspace{5\baselineskip}
\subsection{Interprétation matricielle}

\begin{thmdef}
On note $f:\Omega\to \mathbb R^m$, $\Omega\subset \mathbb R^n$, on écrit $f=(f_1, \cdots, f_m)$. On suppose que $f$ est différentiable en $a\in\Omega$ et on note $\mathcal C_n$ (resp. $\mathcal C_m$) la base canonique de $\mathbb R^n$ (resp. $\mathbb R^m$). Alors: \begin{enumerate}
    \item On appelle matrice jacobienne de $f$ en $a$ la matrice \[
        J_f(a)=\left(\frac{\partial f_i}{\partial x_j}(a)\right)_{\substack{1\leq i\leq n\\ 1\leq j\leq n}}
    \]
    \item Si $x\underset{\mathcal C_n}{\longleftrightarrow}X$ alors \[
        \diff f_a(x)\underset{\mathcal C_m}\longleftrightarrow J_f(a)X
    \]
\end{enumerate}
\end{thmdef}

\begin{proof}
\[
    \diff f_a(x)=\sum_{i=1}^n\underbrace{\frac{\partial f}{\partial x_i}(x)}_{\displaystyle \left(\frac{\partial f_1}{\partial x_i}(x), \cdots, \frac{\partial f_m}{\partial x_i}(x)\right)}\overbrace{\diff x_i(x)}^{=x_i} \qquad \underset{\mathcal C_n}{\longleftrightarrow}\qquad \begin{pmatrix}
    \sum_{i=1}^n\limits \frac{\partial f_1}{\partial x_i}(a)x_i \\
    \vdots \\
    \sum_{i=1}^n\limits \frac{\partial f_m}{\partial x_i}(a)x_i 
    \end{pmatrix}=J_f(a)X
\]
\end{proof}

\section{Opérations sur les différentielles}

\begin{prop}
Si $f,g:\omega\subset E\to F$ sont différentiables en $a$, alors pour $\lambda\in\mathbb R$, $\lambda f+g$ est différentiable en $a$ et $\diff (\lambda f+g)_a=\lambda df_a+\diff g_a$
\end{prop}

\begin{prop}
    \Hyp $B:E_1\times E_2\to F$ bilinéaire, $f_1:\Omega_1\to E_1$ différentiable en $a_1$, $f_2:\Omega_2\to E_2$ différentiable en $a_2$.
    \Conc $f=B(f_1, f_2)$ est différentiable en $a=(a_1, a_2)$ et \[\diff f_a(h_1, h_2)=B(\diff {f_1}_{a_1}(h_1), f_2(a_2))+B(f_1(a_1), \diff {f_2}_{a_2}(h_2)) \]
\end{prop}

\begin{proof}
On développe $f(a+h)=B(f_1(a_1+h_1), f_2(a_2+h_2))$
\end{proof}

\begin{ex} ~
\begin{itemize}
    \item $(E, (\; |\; ))$ euclidien, $\varphi:(x, y)\longmapsto (x|y)$ \[
        \diff \varphi_{(x, y)}(h, k)=(h|y)+(x|k)
    \]
    \item $\varphi = \det $. On note \[
        A = (C_1 \cdots C_n) \quad H=(H_1 \cdots H_n)
    \]
    alors \[
        \det(A+H)=\det(A)+\det(H_1, C_2, \cdots, C_n)+\cdots + \det(C_1, \cdots, C_{n-1}, H_n)+o_0(H)
    \]
    On en déduit \begin{align*}
        \diff \varphi_A(H) &= \sum_{j=1}^n \det(C_1, \cdots, H_i, \cdots, C_n)\\
        &= \sum_{j=1}^n\sum_{i=1}^n(-1)^{i+j}h_{i, j}\Delta_{i, j}(A) \\
        &= \Tr(\Com(A)^\intercal H)
    \end{align*}
\end{itemize}
\end{ex}

\section{Composition des fonctions différentielles}

\begin{prop}
    \Hyp $f : \Omega \subset E \to F$ différentiable en $a$ telle que $f(\Omega)\subset \Omega'$, $g : \Omega'\subset F \to G$ différentiable en $f(a)$
    \Conc $g\circ f$ est différentiable en $a$ et $\diff (g\circ f)_a=\diff g_{f(a)}\circ \diff f_a$
\end{prop}

\begin{proof}
    On a \[
        g(f(a)+k)=g(f(a))+\diff g_{f(a)}(k)+o(k)
    \]
    et \[
        f(a+h)=f(a)+\underbrace{\diff f_a(h)+o(h)}_{k}
    \]
    donc \[
        g\circ f(a+h)=g(f(a))+\diff g_{f(a)}(\diff f_a(h)+o_0(h)) + \underbrace{o(\underbrace{\diff f_a(h)+o_0(h)}_{o(h)})}_{o(h)}
    \]
    d'où \conc
\end{proof}

\begin{rem}
    $E=\mathbb R^n$, $F=\mathbb R^m$, $G=\mathbb R^p$, $f=(f_1, \cdots, f_m):E\to F$, $g=(g_1, \cdots, g_p):F\to G$, $\mathcal C_k$ base canonique de $\mathbb R^k$.

    Si $x \underset{\mathcal C_n}\longleftrightarrow X$ alors $\diff f_a(x)\underset{\mathcal C_m}\longleftrightarrow J_f(a)X$ et $\diff (g\circ f)_a(x)\underset{\mathcal C_p}\longleftrightarrow J_{g\circ f}(a)X\underset{\mathcal C_p}\longleftrightarrow J_g(f(a))J_f(a)X$. C'est vrai pour tout $X$ donc \[
        \underbrace{J_{g\circ f}(a)}_{\in\mathcal M_{p,n}} = \underbrace{J_g(f(a))}_{\in\mathcal M_{p,m}}\underbrace{J_f(a)}_{\in \mathcal M_{m,n}}
    \]
    On va identifier le coefficient $i, j$. \begin{align*}
        & \left[ J_{g\circ f}(a) \right]_{i, j}=\sum_{k=1}^m \left[ J_g(f(a)) \right]_{i, k} \left[ J_f(a) \right]_{k, j} \\ 
        \iff & \frac{\partial (g_i\circ f)}{\partial x_j} (a)=\sum_{k=1}^m \frac{\partial g_i}{\partial y_k}(f(a)) \frac{\partial f_k}{\partial x_j} (a) \\
        \iff & \colorboxed{red}{\partial_j (g_i\circ f)(a)=\sum_{k=1}^m\partial_kg_i(f(a))\partial_jf_k(a)} \quad \text{ \emph{(règle de la chaîne)} }
    \end{align*}

    En particulier, si $p = n = 1$, $\varphi:t\in \mathbb R\to g(f_1(t), \cdots, f_m(t))$ alors \[
        (g\circ f)'(a)=\sum_{k=1}^m \frac{\partial g}{\partial y_k}(f(a))f'_k(a)
    \]
\end{rem}

\begin{rem}
    Si $f:\Omega\subset E\to \mathbb R$ est différentiable en $a$ et $f(a)\neq 0$ alors $\frac 1f$ est différentiable en $a$ car $x\longmapsto \frac 1x$ est différentiable en $f(a)$
\end{rem}

\begin{rem}
    Si $f:\Omega\subset E\to F$ est $a\in\Omega$, $v\in E$ alors $\varphi:t \longmapsto f(a+tv)$ est définie au $\mathcal V(0)$ et $\varphi'(t)=\diff f_{a+tv}(v)$
\end{rem}

\needspace{5cm}
\section{Applications de classe {$\mathcal C^1$}}

\begin{thmdef}
    \Hyp $f:\Omega\subset E\longrightarrow F$
    \begin{concenum}
    \item On dira que $f$ est $\mathcal C^1$ sur $\Omega$ si $f$ est différentiable sur $\Omega$ et si $a\in\Omega \longmapsto \diff f_a\in\mathcal L(E, F)$ est $\mathcal C^0$
        \item Une fonction $\mathcal C^1$ est continue
        \item On note $\mathcal B=(e_1, \cdots, e_n)$ une base de $E$, $\mathcal C=(u_1, \cdots, u_m)$ une base de $F$ et $f=(f_1, \cdots, f_m)$. Il y a équivalence entre: \begin{enumerate}
            \item $f$ est $\mathcal C^1$ sur $\Omega$
            \item Les $f_i$ ont toutes leurs dérivées partielles sur $\Omega$ et les $\partial_jf_i$ sont $\mathcal C^0$ sur $\Omega$
        \end{enumerate}
    \item On note $\mathcal C^1(\Omega, F)$ l'ensemble de ces fonctions
    \end{concenum}
\end{thmdef}

\begin{proof}~
    \begin{enumerate}
        \setcounter{enumi}{1}
        \item Les fonctions $\mathcal C^1$ sont différentiables donc continues.
        \item $(a) \implies (b)$: \[
                J_f(a)=\mathcal M_{\mathcal B, \mathcal C}(\diff f_a)= \left( \frac{\partial f_i}{\partial x_j} (a) \right)_{\substack{1\leq i\leq m\\1\leq j\leq n}}
            \]
            est continue par rapport à $a$ par hypothèse.

            $(b) \implies (a)$: Il suffit de montrer que $f$ est différentiable sur $\Omega$. On commence par le cas $F=\mathbb R$. \begin{align*}
                f(a+h)-f(a) = \phantom{+} & f(a_1+h_1, \cdots, a_n+h_n)-f(a_1, a_2+h_2, \cdots, a_n+h_n) \\
                            + & f(a_1, a_2+h_2, \cdots, a_n+h_n)-f(a_1, \cdots, a_i, a_{i+1}+h_{i+1}, \cdots, a_n+h_n) \\
                            + & f(a_1, \cdots, a_i, a_{i+1}+h_{i+1}, \cdots, a_n+h_n) + \cdots \\
                            + & f(a_1, \cdots, a_{n-1}, a_n+h_n)-f(a_1, \cdots, a_n)
            \end{align*}
            et chaque ligne est fonction d'une seule variable, donc le théorème des accroissements finis donne 
            \begin{align*}
                \exists c_{i,h}\in]a_i,a_i+h_i[, \quad f(a+h)-f(a)=\phantom+ &
                    h_1\frac{\partial f}{\partial x_1}(c_{1,h}, a_2+h_2, \cdots, a_n+h_n)\\
                + & h_2 \frac{\partial f}{\partial x_2}(a_1, c_{2,h},a_3+h_3,\cdots, a_n+h_n)\\
                + & \hspace{2cm}\vdots \\
                + & \underbrace{h_n \frac{\partial f}{\partial x_n} (a_1, \cdots, a_{n-1}, c_{n,h})}_{\longrightarrow h_n \frac{\partial f}{\partial x_n}(a) \text{ donc } =h_n \frac{\partial f}{\partial x_n}(a)+o(1) } \\
                =\phantom + & \underbrace{h_1 \frac{\partial f}{\partial x_1}(a)+\cdots + h_n \frac{\partial f}{\partial x_n} (a) }_{\text{AL}}  + \underbrace{h_1o(1)+\cdots h_no(1)}_{o(h)}
            \end{align*}
            donc $f$ est différentiable sur $\Omega$. Si $F\neq R$ on raisonne coordonnée par coordonnée.
    \end{enumerate}
\end{proof}

\begin{prop} ~
    \begin{enumerate}
        \item $\mathcal C^1(\Omega, F)$ est un $\mathbb R$-e.v.
        \item Si $\mathcal C$ base de $F$ est $f:\Omega\longrightarrow F$ donnée par $f=(f_1, \cdots, f_n)$ (dans $\mathcal C$) alors il y a équivalence entre \begin{enumerate}
            \item $f$ est $\mathcal C^1$
            \item Les $f_i$ sont toutes $\mathcal C^1$
        \end{enumerate}
    \item Si $f:\Omega\longrightarrow F$ et $g: \Omega'\supset f(\Omega)\longrightarrow G$ sont $\mathcal C^1$ alors $g\circ f$ aussi
    \end{enumerate}
\end{prop}

\begin{proof}
    1. et 2. sont faciles, 3. est déjà vu.
\end{proof}

\begin{thm}
    \Hyp $f:\Omega\subset E\longrightarrow \mathbb R$, $a,b\in\Omega$ tels que $[a, b]\subset \Omega$
    \begin{concenum}
    \item \[ f(b)-f(a)=\int_0^1\mathrm df_{(1-t)a+tb}(b-a)\diff t\]
    \item Si $\nop \diff f_x\nop\leq M$ pour $x\in[a, b]$ alors $|f(b)-f(a)|\leq M\|b-a\|$
    \end{concenum}
\end{thm}

\begin{proof}
    \begin{enumerate}
        \item $\varphi:t\longmapsto f(a+t(b-a))$ est telle que $\varphi'(t)=df_{(1-t)a+tb}(b-a)$ et \[
                f(b)-f(a)=\varphi(1)-\varphi(0)=\int_0^1\varphi'
            \]
        \item \[
                |f(b)-f(a)|\leq \int_0^1\|\diff f_{(1-t)a+tb}(b-a)\|\diff t\leq M\int_0^1\|b-a\|\diff t
            \]
    \end{enumerate}
\end{proof}

\section{Étude de régularité}

\subsection{$f:(x, y)\longmapsto \min(x^2, y^2)$}

D'abord, \[
    f(x, y)=\frac{x^2+y^2-|x^2-y^2|}2
\]
donc $f$ est $\mathcal C^0$. On note $\Delta$ le fermé $\{(x, y)\in\mathbb R^2, x^2=y^2\}$, $\mathbb R^2\setminus \Delta$ est ouvert de $f\in \mathcal C^1(\mathbb R^2\setminus \Delta, \mathbb R)$ comme composée de fonctions $\mathcal C^1$.

Il reste à savoir si $f$ est différentiable sur $\Delta$ et si éventuellement la différentielle est continue.


\chapter{Compléments sur les polynômes}

\ifsolo
~

\vspace{1cm}

\begin{center}
\textbf{\LARGE Compléments sur les polynômes} \\[1em]
\end{center}
\fi

Tous les théorèmes ou propositions présents dans ce chapitre sont soit déjà connus (donc au programme) soit nouveaux et sont alors des compléments de cours (hors programme). Le chapitre dans sa totalité n'est pas au programme, seuls quelques rappels du programme y sont présents car pertinents et/ou utiles au développement d'autres compléments.

\ifsolo
\tableofcontents
\else
\minitoc
\fi
\thispagestyle{empty}

\ifsolo \newpage \setcounter{page}{1} \fi
\section{Rappels}

$\mathbb K$ désigne un corps (donc commutatif). Un polynôme de $\mathbb K[X]$ s'écrit $P$ ou $P(X)$, l'évaluation en $a\in\mathbb K$ de $P$ se note $P(a)$. L'application \[
    P\in\mathbb K[X] \longmapsto (a\longmapsto P(a)\in\mathbb K)
\]
est un morphisme d'algèbres qui est injectif si $\mathbb K$ est infini. Pour $P\in\mathbb K[X]\setminus \{0\}$, $a\in\mathbb K$ est une racine de $P$ ssi \[
    P(a)=0\iff (X-a)\;|\;P
\]
et $a$ est racine de multiplicité $\alpha \geq 1$ ssi \[
    \begin{cases}
        P(a)=P'(a)=\cdots=P^{(\alpha - 1)}(a)=0 \\
        P^{(\alpha)}(a)\neq 0
    \end{cases}
    \iff \begin{cases}
        (X-a)^\alpha \; |\; P\\
        (X-a)^{\alpha+1}\;\not|\;P
    \end{cases}
\]
Si $a$ n'est pas racine de $P$, on convient que $a$ est racine de multiplicité $0$.

\begin{notation} ~
    \begin{itemize}
        \item
            On note $Z_\mathbb K(P)$ l'ensemble des racines de $P$ dans $\mathbb K$
        \item On note $\mathrm{mult}(P, a)$ la multiplicité de la racine $a$ dans $P$
    \end{itemize}
\end{notation}

\begin{thm}
    Soit $P\in \mathbb K[X]\setminus \{0\}$. Le polynôme $P$ a au plus $\deg P$ racines dans $\mathbb K$.
    En particulier, le polynôme nul est le seul polynôme qui a une infinité de racines.
\end{thm}

\begin{thm}
    Soit $P\in\mathbb K[X]$.
    \begin{enumerate}
        \item \[
                P(X)=\sum_{k=0}^{+\infty} \frac{P^{(k)}(0)}{k!}X^k
            \]
        \item \[
                \forall a\in\mathbb K, \quad P(X)=\sum_{k=0}^{+\infty}\frac{P^{(k)}(a)}{k!}(X-a)^k
            \]
    \end{enumerate}
\end{thm}
\begin{proof} ~
    \begin{enumerate}
        \item L'application \[
                \psi: P\longmapsto P(X)-\sum_{k=0}^{+\infty}\frac{P^{(k)}(a)}{k!}(X-a)^k
            \]
            est linéaire nulle sur les vecteurs de la base canonique.
    \end{enumerate}
\end{proof}

\begin{rem}
    On montre de même que \[
        P(a)=\sum_{k=0}^{+\infty}\frac{P^{(k)}(X)}{k!}(a-X)^k
    \]
\end{rem}

\section{Factorisation des polynômes}

\subsection{Rappels}

\begin{thm}
    Soit $P\in\mathbb K[X]$, $a_1, \cdots, a_p$ deux à deux distincts et $n_1, \cdots, n_p\in\mathbb N^\star$. Si $a_1, \cdots, a_p$ sont des racines de multiplicités $n_1, \cdots, n_p$ de $P$, alors \[
        (X-a)^{n_1}\cdots (X-a_p)^{n_p}\;|\;P
    \]
    En particulier, si $\deg P=\sum_i n_i$ alors il existe $\lambda\in\mathbb K^\star$ tel que \[
        P(X)=\lambda(X-a_1)^{n_1}\cdots (X-a_p)^{n_p}
    \]
\end{thm}

\subsection{Théorème de d'Alembert-Gauss}

\begin{thm}[D'Alembert-Gauss] \index{D'Alembert-Gauss (théorème de -- )}
    Tout polynôme non constant de $\mathbb C[X]$ admet au moins une racine
\end{thm}

\begin{proof}[Preuve de Gauss]
    \begin{itemize}
        \item $P\in \mathbb C[X]$ non constant a même racine que $\frac1{\deg P}P$ donc on suppose $P$ unitaire.
        \item Si $P(X)=a_0+\cdots+ a_{n-1}X^{n-1}+X^n$, et $M=\max |a_i|$, alors pour $|z|\geq M+2$, \begin{align*}
                |P(z)|&\geq |z|^n-M\frac{|z|^n-1}{|z|-1}=\frac{|z|^{n+1}-|z|^n-M|z|^n+M}{|z|-1}\\
                      &\geq \frac{|z|^n}{|z|-1} \left( |z|-1-M \right)\geq |z|^n\left(1-\frac M{M+1}\right)\geq \frac{(M+2)^n}{M+1}>M
            \end{align*}
        \item \begin{align*}
                \inf_{\mathbb C}|P|&=\min\left(\underbrace{\inf_{|z|>M+2}|P(z)|}_{>M}, \underbrace{\inf_{|z|\leq M+2}|P(z)|}_{\leq |P(0)|\leq M}\right) \\&=\inf _{|z|\leq M+2}|P(z)|
            \end{align*}
            Par compacité il existe $z_0\in\mathcal B_f(0, M+2)$ tel que $\inf_{\mathbb C}|P|=|P(z_0)|$.
        \item Si $P(z_0)\neq 0$ alors il existe $k\geq 1$ tel que \begin{align*}
                P(z_0+h)&=P(z_0)+\frac{P^{(k)}(z_0)}{k!}h^k+o(h^k), \quad P(z_0)\neq 0 \\
                        &=P(z_0)\left(1+h^k\frac{P^{(k)}(z_0)}{P(z_0)k!}+o(h^k) \right)
            \end{align*}
        \item En écrivant \[
                \frac{P^{(k)}(z_0)}{P(z_0)k!}=\rho e^{i\theta}, \quad h=\epsilon e^{-i\frac{\theta+\pi}k}
            \]
            on obtient \[
                P(z_0+\epsilon e^{-i\frac\theta k})=\underbrace{P(z_0)\underbrace{\left(1 - \epsilon^k \rho + o(\epsilon^k)\right)}_{|\;\;|<1}}_{|\;\;|<|P(z_0)|}
            \]
            absurde.
    \end{itemize}
\end{proof}

\section{Polynômes sous contraintes}

\subsection{Contraintes sur les zéros}

\begin{exo}
    Si $Z_{\mathbb C}(P)=Z_{\mathbb C}(Q)$ et $Z_{\mathbb C}(P+1)=Z_{\mathbb C}(Q+1)$ pour $P,Q$ non constants de $\mathbb C[X]$, alors $P=Q$
\end{exo}

On suppose par l'absurde que $P\neq Q$ et par exemple $n=\deg P\geq \deg Q=m$. On note $u_1, \cdots, u_r$ les racines de $P$ (et donc de $Q$), $v_1, \cdots, v_s$ les racines de $P+1$ (donc de $Q+1$). Les $u_i$ et les $v_j$ sont tous distincts, donc \[
    \deg P=n\geq \deg(P-Q)\geq r+s \quad\text{car }P-Q\text{ s'annule en $u_i$ et $v_j$ et }P-Q\neq 0
\]
On a $P'=(P+1)'$ donc $P'$ a $n-r$ racines dans les $u_i$ avec multiplicité, et $n-s$ dans les $n_j$ donc \[
    n-r+n-s\leq \deg P'=n-1\iff n+1\leq r + s
\]
Or $r+s\leq n$ absurde et donc $P = Q$

\subsection{Coefficients unitaires}

\begin{exo}
    Déterminer les $P\in\mathbb R[X]$ non constants avec les coefficients $\pm1$ et dont toutes les racines sont réelles.
\end{exo}

\begin{proof}[Résolution]
    On note $P$ un polynôme de degré $n\geq 2$ qui convient et quitte à changer $P$ en $-P$, $P$ est unitaire.
    \begin{itemize}
        \item $P(X)=X^n+p_{n-1}X^{n-1}+\cdots +p_0$
        \item $r_1, \cdots, r_n$ les racines (non nulles car $p_0\neq 0$) donne \[
                r_1^2+\cdots +r_n^2=p_{n-1}^2-2p_{n-2}>0 \implies p_{n-2}=-1 \text{ et }r_1^2+\cdots+r_n^2=3
            \]
        \item \[
                p_0 X^n P \left( \frac1X \right)=X^n+p_0p_1X^{n-1}+p_0p_2X^{n-2}+\cdots
            \]
            vérifie les mêmes propriétés donc \[
                \frac 1{r_1^2} +\cdots +\frac1{r_n^2}=3
            \]

        \item \[
                \sum_{i=1}^n\underbrace{\left(r_i^2+\frac1{r_i^2}\right)}_{\geq 2}=6 \implies n\leq 3
            \]
    \end{itemize}
    On doit étudier $12$ polynômes
    \begin{center}
        \begin{tabular}{lll}
            \sout{$X^2+X+1$} & \sout{$X^3+X^2+X+1$} & \sout{$X^3-X^2+X+1$} \\
            $X^2+X-1$ & \sout{$X^3+X^2+X-1$} & \sout{$X^3-X^2+X-1$} \\
            \sout{$X^2-X+1$} & \sout{$X^3+X^2-X+1$} & $X^3-X^2-X+1$ \\
            $X^2-X-1$ & $X^3+X^2-X-1$ & \sout{$X^3-X^2-X-1$}
        \end{tabular}
    \end{center}
\end{proof}

\subsection{Équation fonctionnelle}

\begin{exo}
    Déterminer les polynômes $P\in\mathbb C[X]$ tels que $P(X)P(X+2)+P(X^2)=0$
\end{exo}

\begin{proof}[Résolution] $-1$ et $0$ sont les deux polynômes constants solutions.

    Soit $P$ non constant solution et $a$ une racine. $P(a)=0\implies P(a^2)=0$ donc $P$ admet une infinité de racines sauf si $|a|=1$ ou $a=0$.

    Si $a=0$ est racine alors $P(-2)P(0)+P(4)=0$ donc $4$ est racine ce qui est absurde. Donc $|a|=1$.

    Si $a$ est racine alors $(a-2)^2$ racine donc $|a-2|=1=|a|$ donc $a=1$ (intersection de deux cercles), donc $1$ seule racine de $P$ et donc pour un $k\geq 1$, \[
        P=\lambda(X-1)^k
    \]
    En calculant, on trouve $\lambda^2+\lambda=0$ et $\lambda \neq 0$ donc $\lambda = -1$
\end{proof}

\begin{exo}
    Trouver les polynômes $P\in\mathbb R[X]$ tels que \[
        (X+1)P(X-1)-(X-1)P(X)
    \]
    est un polynôme constant
\end{exo}

\begin{proof}[Résolution]
    On suppose que $P$ convient et $(X+1)P(X-1)-(X-1)P(X)=\lambda$. \begin{itemize}
        \item $2P(-1)=\lambda$ et $2P(0)=\lambda$ donc $P(0)=P(-1)$
        \item $Q(X)=P(X)-\frac\lambda 2$ s'annule en $-1$ et $0$ donc $Q(X)=X(X+1)A(X)$ i.e.  $P(X)=X(X+1)A(X)+\frac\lambda 2$.
        \item \[
                (X+1)((X-1)X A(X-1)+\frac\lambda2)-(X-1)X(X+1)A(X)-(X-1)\frac\lambda2=\lambda
            \]
            donc \[
                X(X-1)(X+1)(A(X-1)-A(X))=0 \implies A(X)=A(X-1)
            \]
            et donc $A$ est un polynôme constant. Donc $P$ est de la forme $P(X)=aX(X+1)+\mu, \quad a, \mu\in\mathbb R$
    \end{itemize}
\end{proof}

\section{Liens coefficients et racines}

\subsection{Relations de Viète}

\index{Viète!relations}

Pour $x_1, \cdots, x_n\in\mathbb C$ on note: \[
    \sigma_i(x_1, \cdots, x_n)=\sum_{\substack{I\subset \llbracket 1, n\rrbracket\\|I|=i}}\prod_{l\in I}x_l
\]
Soit $P(X)=a_nX^n+\cdots+a_0$ un polynôme de degré $n$ et $x_1, \cdots, x_n$ ses racines dans $\mathbb C$. \begin{itemize}
    \item \begin{align*}
            P(X)=a_n(X-x_1)\cdots (X-x_n)=a_n(X^n-\sigma_1 X^{n-1}+ \sigma_2 X^{n-2}+\cdots (-1)^n\sigma_nX^{n-n})
        \end{align*}
    \item Ainsi, \[
            \forall i\in\llbracket 1, n\rrbracket, \quad (-1)^i\sigma_i=\frac{a_{n-i}}{a_n}
        \]
\end{itemize}
En particulier, on en déduit \[
    x_1+\cdots+x_n = -\frac{a_{n-1}}{a_n}
\]
et \[
    x_1\times \cdots\times x_n=(-1)^n\frac{a_0}{a_n}
\]

\subsection{Somme de Newton}

On note $P(X)=a_nX^n+\cdots+a_0$ un polynôme de $\mathbb C[X]$, $x_1, \cdots, x_n$ ses racines. On note $S_k=x_1^k+\cdots + x_n^k$ appelée $k-$ième somme de Newton \index{Newton!sommes}.

Pour $k\geq n$, \begin{align*}
    a_nS_k+a_{n-1}S_{k-1}+\cdots + a_0S_{k-n}&=a_nx_1^k+\cdots +a_0x_1^{k-n}+\cdots +a_nx_n^k+\cdots +a_0x_n^{k-n} \\ &=x_1^{k-n}P(x_1)+\cdots + x_n^{k-n}P(x_n)=0
\end{align*}

Avec les relations de Viète, pour $k\geq n$, \[
    \underbrace{\sigma_0}_{=\; 0 \text{ convention}} S_k-\sigma_1S_{k-1}+\cdots + (-1)^nS_{k-n}=0
\]

Il est possible de calculer $S_1, \cdots, S_{n-1}$ par des formules de récurrence. \begin{itemize}
    \item $S_0=n$ clair. On va montrer que \[
            k\sigma_k=\sum_{i=1}^k(-1)^{i-1}\sigma_{k-i}S_i
        \]
    \item On note $P$ le polynôme $P/\dom P$\[
            P(x)=\prod_{i=1}^n(x-x_i)=\frac1{a_n}\sum_{k=0}^na_kx^k=\sum_{k=0}^n\sigma_{n-k}(-1)^{n-k}x^k
        \]
        d'où \[
            x^nP\left(\frac1x\right)=\underbrace{\prod_{k=1}^n(1-xx_k)}_{Q(x)}=\sum_{k=0}^n\sigma_{n-k}(-1)^{n-k}x^{n-k}=\sum_{k=0}^n\sigma_k(-1)^kx^k
        \]
    \item \begin{align*}
            \sum_{k=0}^n k\sigma_k(-1)^kx^k&= x\sum_{k=1}^n \frac{Q(x)}{1-xxk}(-x_k) \\
                                           &=-\sum_{k=1}^n\frac{xx_k}{1-xx_k}Q(x) \\
                                           &=-\left(\sum_{k=1}^n\sum_{i\geq 0}(xx_k)^{i+1}\right)Q(x) \\
                                           &=-\sum_{i\geq 0}x^{i+1}S_{i+1}\sum_{l=0}^n(-1)^k\sigma_kx^k
        \end{align*}

        En identifiant le coefficient de $x^k$, \begin{align*}
            (-1)^kk\sigma_k=-\sum_{l=0}^{k-1}(-1)^l\sigma_lS_{k-l} \iff k\sigma_k &=-\sum_{l=0}^{k-1}(-1)^{k-l}\sigma_lS_{k-l} \\
                                                                                  & =\sum_{l=1}^k(-1)^{l-1}\sigma_{k-l}S_l
        \end{align*}
\end{itemize}

\section{Polynômes scindés sur $\mathbb R$}

\subsection{Une CNS}

\begin{exo}
    Montrer que $P\in\mathbb R[X]$ de degré $n$ est scindé sur $\mathbb R$ ssi \[
        \forall z\in\mathbb C, \quad |P(z)|\geq |\dom(P)|\times |\Im(z)|^n
    \]
\end{exo}

\begin{proof}[Résolution]
    \emph{($\Longrightarrow$)} clair, les racines complexes de $P$ ont une partie imaginaire nulle.

     \emph{($\implies$)} On écrit $P(X)=\dom(P)(X-a_1)\cdots (X-a_n), \quad a_i\in\mathbb R$. En injectant $z=a+ib$, \[
        |P(z)|=|\dom P|\times |a-a_1+ib| \cdots |a-a_n+ib|\geq |\dom P| |b|^n
    \]
\end{proof}

\subsection{Opérations stabilisant $\mathcal S$}

On note $\mathcal S$ l'ensemble des polynômes scindés de $\mathbb R[X]$. On note \[
    P(X)=\lambda (X-a_i)^{\alpha_1}\cdots (X-a_p)^{\alpha_p}, \quad a_1<\cdots <a_p
\]
et on a les propriétés de stabilité suivantes.

\vspace{.2cm}

 Rolle donne que $P'$ s'annule au moins une fois sur $]a_i, a_{i+1}[$ pour $1\geq i < p$. Il y a alors au total $(\alpha_1-1)+\cdots+(\alpha_p-1)+p-1=n-1$ racines réelles (avec multiplicité) et $\deg P'=n-1$ donc $P'$ est scindé et $\mathcal S$ est stable par dérivation.

\vspace{.2cm}

 Pour $a\in\mathbb R$ fixé, on note $g(x)=e^{ax}P(x)$, $g'(x)=(P'(x)+aP(x))e^{ax}$. $P'+aP$ s'annule en $a_i$ avec multiplicité $\geq \alpha_i-1$. Rolle donne que $g'$ s'annule au moins une fois sur $]a_i, a_{i+1}[$ donc $P'+aP$ aussi. Ce polynôme a $n-1$ racines réelles et est de degré $n$, donc par division la dernière racine est réelle. Donc $P'+aP\in\mathcal S$

\vspace{.2cm}

 Plus généralement, si $D$ est l'opérateur de dérivation polynomiale, pour $Q\in\mathcal S$ alors
\begin{itemize}
    \item
\[
    Q(X)=\lambda (X-b_1)\cdots (X-b_p)
\]
\item \[
        Q(D)=\lambda(D-b_1\mathrm{id})\circ \cdots \circ (D-b_p \mathrm{id})
    \]
\end{itemize}
et pour $P\in\mathcal S$, $Q(D)(P)\in\mathcal S$.

\subsection{Faisceau de polynômes à racines entrelacées (ENS)}

\begin{exo}
    On note $P, Q\in\mathbb R[X]$ deux polynômes scindés à racines simples distinctes tels qu'entre deux racines de l'un, il y en aie au moins une de l'autre. Montrer que $\lambda P+\mu Q$ est scindé sur $\mathbb R$
\end{exo}

\begin{proof}[Résolution]
    On suppose que $P$ a la plus petite racine, on note $x_1 <\cdots<x_n$ les racines de $P$. Sur $]x_1, x_2[$, il y a au moins une racine de $Q$ et au plus une (sinon $x_1$ et $x_2$ ne sont pas consécutives) donc les $n-1$ premières racines de $Q$ satisfont \[
        x_1<y_1<\cdots <x_{n-1}<y_{n-1}<x_n
    \]
    Il y a lors deux cas: $Q$ possède une racine $y_n>x_n$, ou $Q$ a $n-1$ racines réelles. On se place dans le premier cas (l'autre se traite similairement).

    Comme $\lambda, \mu\in\mathbb R$ sont quelconques, on peut supposer $P$ et $Q$ unitaires donc \[
        P(X)=(X-x_1)\cdots (X-x_n) \qquad Q(X)=(X-y_1)\cdots (X-y_n)
    \]
    Pour $\lambda, \mu\neq 0$, les racines réelles de $\lambda P+\mu Q$ sont différentes des $x_i, y_i$, donc sont les mêmes que celles de $\frac{\lambda P+\mu Q}{PQ}$ sur $\mathbb R\setminus \{x_i, y_i\}$. \[
        \varphi(x)=\frac{\lambda P(x)+\mu Q(x)}{P(x)Q(x)}=\frac\lambda{Q(x)}+\frac\mu{P(x)}
    \]
    Si $\lambda, \mu>0$, alors $\varphi$ s'annule $n$ fois car il y a changement de signe entre $]x_i, y_i[$ (les polynômes sont unitaires donc on connaît exactement les signes de $P$ et $Q$). Si $\lambda>0>\mu$, on trouve $n-1$ racines réelles de la même manière, et par division la dernière est aussi réelle. Les autres cas se traitent de la même manière.
\end{proof}

\section{Valeurs des polynômes}

\subsection{Polynôme à valeurs entières}

On va caractériser les polynômes de $\mathbb C[X]$ tels que $P(\mathbb Z)\subset \mathbb Z$

On définit \[
    \binom X0=1\qquad \text{ et }\qquad \forall n\geq 1,\quad \binom Xn=\frac1{n!}X(X-1)\cdots(X-n+1)
\]

\begin{lmm}
    Pour tout $n\in\mathbb N$ et pour tout $x\in\mathbb Z$, \[
        \binom xn\in\mathbb Z
    \]
\end{lmm}

\begin{proof}
    Pour $0\leq x\leq n-1$, $\binom xn=0$. Pour $x\geq n$, $\binom xn$ est un coef du binôme. Pour $x<0$, $(-1)^n\binom xn$ est un coef du binôme.
\end{proof}

 On introduit \[
    \Delta : P\in\mathbb K[X]\longmapsto P(X+1)-P(X)
\]
Pour $n\geq 1$, \[
    \Delta\binom Xn=\binom X{n-1}\qquad\text{ et }\qquad \Delta\binom X0=0
\]
et $\displaystyle \left(\binom Xn\right)_{n\geq 0}$ est une base de $\mathbb C[X]$ (degré échelonné). On suppose que $P\in\mathbb C[X]$ vérifie $P(\mathbb Z)\subset \mathbb Z$. On note $a_0, \cdots, a_n$ les coefficients de la décomposition de $P$ dans cette base. Alors \[
    P(0)=a_0\in\mathbb Z, \quad a_1=\Delta P(0)\in\mathbb Z, \cdots, \quad a_n=\Delta^nP(0)\in\mathbb Z
\]

Ainsi, $a_0, \cdots, a_n\in\mathbb Z$ et la réciproque est évidente. On a ainsi établi \[
    P(\mathbb Z)\subset \mathbb Z\iff \exists n\in\mathbb N, \exists, a_0, \cdots, a_n\in\mathbb Z, P(X)=\sum_{k=0}^na_k\binom Xk
\]

\begin{rem}
    On a aussi établi la formule de Gregory\index{Gregory (formule de -- )} \[
        P(X)=\sum_{k\geq 0}\Delta^kP(0)\binom Xk
    \]
\end{rem}

\begin{exo}[ENS]
    On note $P\in\mathbb R[X]$ un polynôme de degré $n$ tel que $P(0), P(1), P(4), \cdots, P(n^2)\in\mathbb Z$. Montrer que $\forall a\in\mathbb Z, P(a^2)\in \mathbb Z$.
\end{exo}

\begin{proof}[Résolution]
    $Q(X)=P((X-n)^2)$ de degré $2n$ envoie $2n+1$ valeurs de $\mathbb Z$ dans $\mathbb Z$ donc a des coefficients entiers dans la base explicitée précédemment donc est à valeurs dans $\mathbb Z$ ce qui conclut.
\end{proof}

\subsection{Polynôme stabilisant le cercle unité}

\begin{exo}
    Déterminer les polynômes $P\in\mathbb C[X]$ tels que $\mathbb U$ est stable par $P$.
\end{exo}

\begin{proof}[Résolution]
    Le cas des polynômes constants est immédiat. On note $P(z)=a_0+\cdots +a_nz^n$. On a \[
        P(e^{i\theta})\overline{P(e^{i\theta})}=1\implies P(e^{i\theta})\overline P(e^{-i\theta})=1 \implies P(e^{i\theta})\overline P(e^{-i\theta})e^{in\theta}=e^{in\theta}
    \]
    Si on note $Q(X)=\overline P\left(\frac1X\right)X^n=\bar a_0X^n+\cdots +\bar a_n$ alors $PQ-X^n$ s'annule sur $\mathbb U$ donc est nul donc $P(X)\;|\;X^n$, $\deg P=n$ et $P=\lambda X^n$, où $\lambda=P(1)\in\mathbb U$. La réciproque est claire.
\end{proof}

\begin{exo}
    Variantes \begin{itemize}
        \item $P(\mathbb C)\subset R$. Solution: Formule de Gregory, interpolation de Lagrange, polynôme conjugué, ...
        \item $P(\mathbb Q)\subset \mathbb Q$. Solution: Formule de Gregory
        \item $P(\mathbb Q)\subset \mathbb R\setminus \mathbb Q$ et $P(\mathbb R\setminus \mathbb Q)\subset \mathbb Q$
        \item $P(\mathbb Q)=\mathbb Q$
    \end{itemize}
\end{exo}

\begin{proof}[Résolution]
    Résolution pour $P(\mathbb Q)=\mathbb Q$. Déjà $P\in\mathbb Q[X]$ (par la $2$-ième variante), donc il existe $d$ entier positif minimal tel que $Q=dP\in\mathbb Z[X]$. On se donne $v$ premier avec les coefs $a_i$ de $Q$ et $d$ et $v$ nombre premier. Il existe $\frac ab$ tel que $P(\frac ab)=\frac1v$ donc \[
        v(a_na^n+\cdots +a_0b^n)=db^n \implies v\;|\; b^n\implies v\;|\; b
    \]
    et si $n\geq 2$ alors $v^2\;|\;db^n$ et donc $v|a_na^n+\cdots$ donc $v\;|\;a$ absurde. Par suite, $n=1$, et la réciproque est évidente.
\end{proof}

\subsection{Polynômes positifs}

On va caractériser les polynômes tq $P(\mathbb R_+)\subset \mathbb R_+$.
On pose \[
    \mathcal E=\{P\in\mathbb R[X] \;/\; \exists A, B\in\mathbb R[X], P=A^2+XB^2\}
\]
et on note $\mathcal S$ l'ensemble des polynômes qui conviennent. On a clairement $\mathcal E\subset \mathcal S$, puis $\mathcal E$ est stable par produit car \[
    (A^2+XB^2)(A_1^2+XB_1^2)=(AA_1+XBB_1)^2+X(BA_1+AB_1)^2
\]
Soit $P\in\mathcal S$. Dans la décomposition de $P$ en irréductibles de $\mathbb R[X]$, si on note $\alpha_1, \cdots, \alpha_p$ les puissances des facteurs de degré $1$, on a $(X-a_i)$ change de signe si $\alpha_i$ impair et dans ce cas $a_i<0$ et $X-a_i=\sqrt{-a_i}^2+X1^2\in\mathcal E$. Les autres termes sont dans $\mathcal E$ et donc $\mathcal S=\mathcal E$.

\section{Localisation des racines}

\subsection{Théorème de Gauss-Lucas}

\index{Gauss-Lucas (théorème de -- )}

On note $P\in\mathbb C[X]$ avec $n=\deg P\geq 2$. On va montrer que si $z_1, \cdots, z_n$ sont les racines de $P$ (répétées avec multiplicités) alors les racines de $P'$ sont dans l'enveloppe convexe de $z_1, \cdots, z_n$. On a:
\[
    \frac{P'}P(z)=\frac{\alpha_1}{z-z_1}+\cdots +\frac{\alpha_p}{z-z_p}
\]
donc si $z$ est une racine de $P'$ différente des $z_1, \cdots, z_n$, alors \[
    \alpha_1\frac{z-z_1}{|z-z_1|^2}+\cdots+\alpha_p\frac{z-z_p}{|z-z_p|^2}=0 \iff z \left( \frac{\alpha_1}{|z-z_1|^2}+\cdots+\frac{\alpha_p}{|z-z_p|^2} \right)=\frac{\alpha_1}{|z-z_1|^2}z_1+\cdots+\frac{\alpha_p}{|z-z_p|^2}z_p
\]
Donc $z$ est bien un barycentre de $z_1, \cdots, z_n$.

\subsection{Borne de Cauchy}

\index{Cauchy!borne}

On note $P(z)=z^n+a_{n-1}z^{n-1}+\cdots+a_0\in\mathbb C[X]$ et on pose $M=\max_{i\leq n-1}\limits |a_i|$. On note $z\in\mathbb C$ tel que $|z|>M+1$. Par l'absurde si $P(z)=0$ alors \begin{align*}
    |z^n|=|a_{n-1}z^{n-1}+\cdots +a_0|\leq M\frac{|z|^n-1}{|z|-1} &\implies |z|^{n+1}-|z|^n\leq M|z|^n-M \\ &\implies |z| |z|^n\leq (M+1)|z|^n \\ &\implies |z|\leq M+1\quad\text{ absurde }
\end{align*}
Bilan les racines de $P$ sont dans le disque $D(0, M+1)$

\subsection{Enestrom-Kakeya}

\begin{res}[Cauchy]
    \index{Cauchy!localisation de racines}
    Si $P(x)=x^n-b_{n-1}x^{n-1}-\cdots -b_0$ est un polynôme tel que les $b_i$ sont positifs non tous nuls alors $P$ a une seule racine $a$ dans $\mathbb R_+^\star$, $a$ est racine simple et $\mathbb Z_{\mathbb C}(P)\subset \mathcal B_f(0, a)$.
\end{res}

\begin{proof}
    Pour $x>0$, $f(x)=-\frac{P(x)}{x^n}=\frac{b_0}{x^n}+\cdots + \frac{b_{n-1}}{x}-1$ est décroissante.

    \begin{center}
        \begin{tikzpicture}
            \tkzTabInit{$x$ / 1 , $f(x)$ / 1}{$0$, $a$, $+\infty$}
            \tkzTabVar{+/ $+\infty$, R /, -/ $-1$}
        \end{tikzpicture}
    \end{center}
    Donc $f$ a une unique racine $a>0$, puis \[
        f'(a)=-n\frac{b_0}{x^{n+1}}-\cdots-\frac{b_{n-1}}{x^2}<0 \implies P'(a)\neq 0
    \]
    donc $a$ est racine simple de $P$. Si $z_0$ est une racine de $P$ dans $\mathbb C$ alors $P(z_0)=0$ donc \[
        |z_0^n|=|b_{n-1}z_0^{n-1}\cdots b_0|\leq b_{n-1}|z_0^n|+b_0
    \]
    donc $P(|z_0|)\leq 0$ et $f(|z_0|)\geq 0$ donc $|z_0|\leq a$
\end{proof}

\begin{res}[Enestrom-Kakeya]
    \index{Enestrom-Kakeya (théorème de -- )}
    Si $P(x)=a_{n-1}x^{n-1}+\cdots +a_0$ est un polynôme à coefficients strictement positifs, alors si $e$ est une racine complexe de $P$, \[
        \delta=\min_{i\leq n-1}\frac{a_{i-1}}{a_i} \leq |z|\leq \max_{i\leq n-1}\frac{a_{i-1}}{a_i}=\gamma
    \]
\end{res}

\begin{proof}
    On considère le polynôme \begin{align*}
        Q(x)=(x-\gamma)P(x)&=a_{n-1}x^{n}+(a_{n-2}-\gamma a_{n-1})x^{n-1}+\cdots +(a_0-\gamma a_1)x -\gamma a_0\\ &= a_{n-1}\underbrace{\left(x^n-\frac{\gamma a_{n-1}-a_{n-2}}{a_{n-1}}x^{n-1} -\cdots -\frac{\gamma a_1-a_0}{a_{n-1}}x-\frac{\gamma a_0}{a_{n-1}}\right)}_{\text{vérifie les hypothèses du résultat précédent donc a une unique racine réelle $>0$, $\gamma$}}
    \end{align*}
    d'où la seconde inégalité. La première s'obtient en remarquant que $z$ racine de $P$ entraîne $z$ non nul et $\frac 1z$ racine de $X^{n-1}P\left(\frac1X\right)$ auquel on applique le même raisonnement.
\end{proof}

\subsection{Disques de Gershgorin}

\index{Gershgorin (disques -- )}

On se donne $A\in\mathcal M_n(\mathbb C)$ et $\lambda\neq a_{1, 1}, \cdots, a_{n, n}$. On note $D=\mathrm{diag} (a_{1, 1}, \cdots, a_{n, n})$ et $E=A-D$. Si $\lambda$ est valeur propre, alors \[
    A_\lambda=\underbrace{A-\lambda I_n}_{\text{non inversible}}=\underbrace{D-\lambda I_n}_{\text{inversible}}+E
\]
Puis $\exists X\in\mathcal M_{n, 1}(\mathbb C), X\neq 0 / A_\lambda X=0$ donc \begin{align*}
    (D-\lambda I_n)X+EX=0&\iff X=-(D-\lambda I_n)^{-1}EX \\&\iff \forall i\in\llbracket 1, n\rrbracket, \quad x_i=-\frac1{a_{i, i}-\lambda}\sum_{k\neq i}a_{i, k}x_k
\end{align*}
On considère un indice $i$ tel que $|x_i|$ est maximal et donc \[
    |a_{i, i}-\lambda|=|\sum_{k\neq i}a_{i, k}\frac{x_k}{x_i}\leq \sum_{k\neq i}|a_{i, k}|=R_i
\]
et donc \[
    \lambda\in\mathcal D(a_{i, i}, R_i) \qquad \text{ reste vrai si } \lambda=a_{i, i}
\]

On a donc \[
    \mathrm{Sp}_{\mathbb C}(A)\subset \bigcup_{i=1}^n \underbrace{\mathcal D(a_{i, i}, R_i)}_{\text{disques de Gershgorin}}
\]

\section{Arithmétique des polynômes}

\subsection{Division euclidienne}

\index{division euclidienne}

\begin{exo}
    Effectuer la division euclidienne de $X^n-1$ par $X^m-1$, pour $n, m\geq 1$
\end{exo}

\begin{proof}[Résolution]
    On note $(u_k)$ la suite définie par \[
        \begin{cases}
            u_0=n, u_1=m \\ u_{k+1} \quad \text{ reste de la division euclidienne de $u_{k-1}$ par $u_k$}
        \end{cases}
    \]
    tant que $u_k\neq 0$ et on note \[
        \begin{cases}
            U_0=X^n-1, U_1=X^m-1\\ U_{k+1}\quad \text{ reste de la division euclidienne de $U_{k-1}$ par $U_k$}
        \end{cases}
    \]
    tant que $U_k\neq 0$

    On vérifie par récurrence que $U_k=X^{u_k}-1$. Les cas $k=0, 1$ sont immédiats. On a $u_{k+1}+au_k=u_{k-1}$ donc par HR si $U_k\neq 0$ (ie $u_k\neq 0$), alors \[
        U_{k-1}=U^{u_{k+1}}X^{au_k}-1=\underbrace{X^{u_{k+1}}-1}_{\deg < \deg U_k}+\underbrace{X^{u_{k+1}}\left(X^{au_k}-1\right)}_{\text{divisible par } U_n}
    \]
    d'où le résultat. On en déduit que les deux suites s'arrêtent en même temps et donc \[
        U_0\land U_1=U_{k-1}=X^{m\land n}-1
    \]
\end{proof}

\subsection{Théorème chinois}

\index{chinois (théorème)}

\begin{exo}
On se donne $P_1, \cdots, P_n$ deux à deux premiers entre eux et $R_1, \cdots, R_n\in\mathbb C[X]$. Montrer qu'il existe $Q_1, \cdots, Q_n$ tels que \[
    P_1Q_1+R_1=\cdots=P_nQ_n+R_n
\]
\end{exo}

\begin{proof}[Résolution]
    C'est une application directe du théorème chinois (dans un anneau principal): si $d_i=\deg P_i-1$ et $N=-1+\sum d_i$, alors \[
        \varphi: P\in \mathbb C_N[X]\longmapsto (P\mod P_1, \cdots, P\mod P_n)\in\mathbb C_{d_1}[X]\times \cdots \times \mathbb C_{d_n}[X]
    \]
    est un isomorphisme d'algèbres.
\end{proof}

\subsection{Théorèmes de Mason, Snyder, et Fermat}
On va montrer le théorème de Mason (ou Mason-Stothers) avec une démonstration de
Snyder\footnote{\url{https://doi.org/10.1007/s000170050074}}

\begin{lmm}
    % \index{Snyder (théorème de -- )}
    Si $P\in\mathbb C[X]$ non nul, alors \[
        \deg P=\deg \left(P\land P'\right)+n_0(P)
    \]
    où $n_0(P)=\#Z_{\mathbb C}(P)$.
\end{lmm}

\begin{proof}
    Le cas $P$ constant est immédiat. On suppose maintenant $P$ non constant: \[
        P(X)=c(X-a_1)^{\alpha_1}\cdots (X-a_p)^{\alpha_p}
    \]
    donc \[
        P\land P'=(X-a_1)^{\alpha_1-1}\cdots (X-a_p)^{\alpha_p-1}
    \]
\end{proof}

\begin{thm}[Mason]
    \index{Mason (théorème de -- )}
    Si $A, B, C$ sont des polynômes complexes deux à deux premiers entre eux avec $A+B=C$ et $(A', B', C')\neq (0, 0, 0)$, alors \[
        \deg C\leq n_0(ABC)-1
    \]
\end{thm}

\begin{proof}
Comme $A=C-B$ et $B=C-A$ les rôles de $A, B, C$ sont symétriques. On peut supposer $A$ non constant. On a alors \[
    \left|\begin{matrix}A&C\\A'&C'\end{matrix}\right|=\left|\begin{matrix}A&B\\A'&B'\end{matrix}\right|
\]
donc $AC'-CA'=AB'-BA'$ donc \[
    A\land A', \quad B\land B', \quad C\land C'\quad\Big|\quad AC'-A'C=AB'-BA'
\]
Si $AC'-A'C=0$ alors $AC'=A'C$ et $C\;|\;C'$ donc $C=0$, et de même $B=0$ donc $A=0$ absurde. Les trois polynômes $A, B, C$ sont premiers entre eux donc les trois PGCD le sont aussi et \[
    A\land A' \quad \times \quad B\land B'\quad \times \quad C\land C'\quad \;\Big|\; AC'-A'C\neq 0
\]
donc \[
    \deg A\land A' + \deg B\land B'+\deg C\land C'\leq \deg(AC'-A'C)
\]
et par le résultat précédent \[
    \deg A+\deg B+\deg C \leq \underbrace{AB'-A'B}_{\leq \deg A+\deg B-1}+ \underbrace{n_0(A)+n_0(B)+n_0(C)}_{n_0(ABC)}
\]
\end{proof}

\begin{thm}[Fermat]
    \index{Fermat!grand théorème (version polynômiale)}
    Si $A, B, C\in\mathbb C[X]$ non constants, $A$ et $B$ premiers entre eux et $A^n+B^n=C^n$ alors $n\leq 2$
\end{thm}

\begin{proof}
On note $p=\max (\deg A, \deg B, \deg C)$ et \[
    np=\max(\deg A^n, \deg B^n, \deg C^n)< n_0(A^nB^nC^n)=n_0(ABC)\leq 3p \implies n<3
\]
\end{proof}

\subsection{Résultant}

\index{résultant}

On note $P=a_pX^p+\cdots + a_0$, $Q=b_qX^q+\cdots + b_0$ de degrés respectifs $p, q$.

On considère l'application linéaire \[
    \varphi: (U, V)\in\mathbb K_{q-1}[X]\times \mathbb K_{p-1}[X]\longmapsto UP+VQ
\]
Si $P\land Q=1$, alors pour $U, V\in\mathrm{Ker}\varphi$, on a $Q\;|\;U$ et $P\;|\;V$ donc $U=V=0$ en comparant les degrés. Sinon, si $P\land Q=\Delta$ non constant, alors $\varphi(Q/\Delta, -P/\Delta)=0$. Ainsi, $\varphi$ est un isomorphisme si et seulement si $P\land Q=1$

La matrice de $\varphi$ dans la base $(X^{q-1}, 0), \cdots, (X, 0), (1, 0), (0, X^{p-1}), \cdots, (0, 1)$ est \[
% Attention pas la même base que les notes du prof ! Comme ça le copier/collé de wikipedia marche bien :-)
\begin{pmatrix}a_{p}&0&\cdots &0&b_{q}&0&\cdots &0\\a_{p-1}&a_{p}&\ddots &\vdots &\vdots &b_{q}&\ddots &\vdots \\\vdots &a_{p-1}&\ddots &0&\vdots &&\ddots &0\\\vdots &\vdots &\ddots &a_{p}&b_{1}&&&b_{q}\\a_{0}&&&a_{p-1}&b_{0}&\ddots &\vdots &\vdots \\0&\ddots &&\vdots &0&\ddots &b_{1}&\vdots \\\vdots &\ddots &a_{0}&\vdots &\vdots &\ddots &b_{0}&b_{1}\\0&\cdots &0&a_{0}&0&\cdots &0&b_{0}\\\end{pmatrix}\defeq R(P, Q)
\]
On appelle résultant de $P$ et $Q$ le déterminant de cette matrice. Ainsi, \[
    P\land Q\neq 1 \iff \det R(P, Q)=0
\]

\section{Polynômes irréductibles}

\index{polynôme!irréductible}

\subsection{Contenu de Gauss}

\index{Gauss!contenu}

On note $P(X)=a_nX^n+\cdots + a_0\in\mathbb Z[X]$ de degré $n\geq 1$. On appelle contenu de $P$ le nombre \[
    a_n\land\cdots\land a_0
\]

\begin{dfn}
    $P\in\mathbb Z[X]$ est dit primitif\index{polynôme!primitif} si $C(P)=1$
\end{dfn}

On va montrer que si $P, Q\in\mathbb Z[X]$ non constants alors \[
    C(PQ)=C(P)C(Q)
\]

On commence par supposer que $P$ et $Q$ sont primitifs et on suppose par l'absurde $C(PQ)\neq 1$. On note $p$ premier diviseur de $C(PQ)$. On note $i_0$ (resp $j_0$) le plus petit indice tel que $p$ ne divise pas $a_{i_0}$ (resp. $b_0$). \[
    P(X)=a_nX^n+\cdots +a_{i_0}X^{i_0}+\underbrace{\cdots}_{\text{divisible par } p} \]\[
    Q(X)=b_nX^n+\cdots +b_{j_0}X^{j_0}+\underbrace{\cdots}_{\text{divisible par } p}
\]
Alors, \[
    [PQ(X)]_{i_0+j_0}=a_{i_0}b_{j_0}+pk \quad \text{non divisible par $p$}
\]
Or c'est un coef de $PQ$ donc divisible par $p$, absurde.

Dans le cas général, \[
    C\left(\frac{PQ}{C(P)C(Q)}\right)=1=\frac{C(PQ)}{C(P)C(Q)}
\]

\begin{res}
    $P\in\mathbb Z[X]$ primitif est irréductible dans $\mathbb Z[X]$ si et seulement s'il l'est dans $\mathbb Q[X]$
\end{res}
\begin{proof}
Si le polynôme est irréductible dans $\mathbb Q[X]$ alors il l'est dans $\mathbb Z[X]$. On suppose $P$ irréductible dans $\mathbb Z[X]$ et $P=AB$ dans $\mathbb Q[X]$. Il existe $p, q$ minimaux positifs tels que $pqP=pAqB$, $pA$ et $qB$ dans $\mathbb Z[X]$ et \[
    C(pqP)=pqC(P)=pq=C(pA)C(qB) \implies P=\frac{pA}{C(pA)}\times \frac{qB}{C(qB)}
\]
absurde (les deux facteurs sont dans $\mathbb Z[X]$).
\end{proof}

\subsection{Critère d'Eisenstein}

\index{Eisenstein (critère de -- )}

\begin{res}[Critère d'Eisenstein]
On note $P=p_nX^n+\cdots + p_0\in\mathbb Z[X]$ non constant de degré $n$. Si $p$ est premier tel que \[
    \begin{cases}
    p\;|\;p_0,\cdots, p_{n-1}\\ p\;\not|\;p_n\\p^2\;\not|\;p_0
    \end{cases}
\]
alors $P$ est irréductible dans $\mathbb Q[X]$
\end{res}
\begin{proof}
On suppose par l'absurde que $P$ n'est pas irréductible dans $\mathbb Q[X]$ donc dans $\mathbb Z[X]$. On écrit alors $P=AB$ avec $A$ et $B$ non constants à coefficients entiers. On a \[
    p\;|\;p_0=a_0b_0
\]
donc par exemple $p$ divise $a_0$ mais pas $b_0$ (il ne peut pas diviser les deux par hypothèse). Si $p$ divise tous les $a_i$, alors $p\;|\;\dom A\times \dom B=p_n$ absurde. On note alors $i_0$ le plus petit $i$ tel que $p\;\not |\;a_i$. $B$ est non constant donc $i_0<n$. \[
    p\;|\; p_{i_0}=a_{i_0}b_0+\underbrace{a_{i_0-1}b_1+\cdots +a_0b_{i_0}}_{\text{divisible par }p} \qquad \text{non divisible par }p
\]
absurde d'où la conclusion.
\end{proof}

\begin{csq}
Pour $p$ premier, on considère \[
    \mu_p(X)=X^{p-1}+\cdots+X+1
\]
Alors, \[
    \mu_p\text{ irréductible }\iff \mu_p(X+1)\text{ irréductible } \iff \frac{(X+1)^p-1}{(X+1)-1}=\sum_{k=1}^p\binom pk X^{k-1}\text{ irréductible}
\]
vrai par le critère d'Eisenstein
\end{csq}

\begin{csq}
Pour tout $n\geq 1$, le polynôme \[
    S_n(X)=1+X+\frac{X}{2!}+\cdots +\frac{X^n}{n!}
\]
est irréductible. En effet, c'est équivalent à vérifier que \[
    n!S_n(X)=X^n+nX^{n-1}+\cdots+n!
\]
est irréductible dans $\mathbb Z[X]$.

Si $n=2m$, alors il existe $p$ premier tel que $m<p<2m$ donc $p<n<2p$ et on conclut avec Einsentein. Sinon, si $n=2m+1$ alors il existe $p$ premier tel que $m<p<2m$ et $p<2m<n$ et $2m\leq 2p-1$ donc $n\leq 2p-1<2p$ et on conclut de même.
\end{csq}

\subsection{{$\mathbb Z[X]$} vs {$\mathbb Z_p[X]$}}

On va montrer que $X^4+1$ est irréductible dans $\mathbb Z[X]$ et n'est jamais irréductible dans $\mathbb Z_p[X]$. Dans $\mathbb C[X]$, \[
    X^4+1=\left(X-e^{i\frac\pi4}\right)\left(X-ie^{i\frac\pi4}\right)\left(X+e^{i\frac\pi4}\right)\left(X+ie^{i\frac\pi4}\right)
\]
Aucun des facteurs ni des produits n'est dans $\mathbb Q[X]$ donc $X^4+1$ est irréductible dans $\mathbb Q[X]$ donc dans $\mathbb Z[X]$.

Pour $p=2$, $(X^4+1)=(X+1)^4$ n'est pas irréductible. On suppose donc $p$ premier impair et on observe que \begin{align*}
    X^4+1&=(X^2)^2-(-1)\\&=(X^2+1)^2-2X^2\\&=(X^2-1)^2-(-2X^2)
\end{align*}
Il est donc suffisant que $-1, 2$ ou $-2$ soit un carré de $\mathbb Z_p$. Si $-1$ et $2$ ne sont pas des carrés, alors $-2$ est un carré. En effet, si on note $\mathcal C^\star$ l'ensemble des carrés non nuls, alors \[
    \varphi: x\in\mathbb Z_p^\star \longmapsto x^2\in\mathcal C^\star
\]
est un morphisme surjectif de noyeau $\ker \varphi=\{\pm1\}$ donc $\#\mathcal C^\star=\frac{p-1}2$. \[
    X^{\frac{p-1}2}-1
\]
a au plus $\frac{p-1}2$ racines ($\mathbb Z_p$ est un corps) et s'annule sur $\mathcal C^\star$ donc $\mathcal C^\star$ est exactement l'ensemble des racines de ce polynôme (le petit théorème de fermat donne $x^{p-1}=1$ pour tout $x$ non nul). Si $-1$ et $2$ ne sont pas carrés, alors \[
    (-2)^{\frac{p-1}2}=(-1)^{\frac{p-1}2}\times (2)^{\frac{p-1}2}=-1\times -1=1
\]
donc $-2$ carré.

\section{Polynômes cyclotomiques (HP)}

\index{polynôme!cyclotomique}

\begin{dfn}
Soit $n\geq 1$. \begin{itemize}
    \item On appelle racine primitive $n$-ième de l'unité tout générateur de $\mathbb U_n$. On note $\mathbb P_n$ ces racines.
    \item Il y a $\varphi(n)$ générateurs de $\mathbb U_n$: les $w^l$ avec $l\land n=1$ et $w$ un générateur.
    \item On appelle $n$-ième polynôme cyclotomique le polynôme \[
        \mu_n(X)=\prod_{u\in\mathbb P_n}(X-u)
    \]
\end{itemize}
\end{dfn}

\begin{prop}
Pour tout $n\in\mathbb N^\star$, \[
    X^n-1=\prod_{d\;|\;n}\mu_d(X)
\]
\end{prop}

\begin{proof}
Pour $d\;|\;n$, on note $E_d=\{u\in\mathbb U_n, \ord u=d\}$, donc \[
    \mathbb U_n =\bigsqcup_{d\;|\;n}E_d
\]
ainsi \[
    X^n-1=\prod_{d\;|\; n}\prod_{u\in E_d}(X-u)=\prod_{d\;|\;n}\mu_d(X)
\]
car $E_d=\mathbb P_d$ (les éléments d'ordre $d$ sont au nombre de $\varphi(d)$ et engendrent tous $\mathbb U_d$).
\end{proof}

\begin{rem}
Si $n=p$ est premier alors \[
    \mu_n(X)=\frac{X^p-1}{X-1}=X^{p-1}+\cdots +1
\]
\end{rem}

\begin{prop}
Pour tout $n\geq 1$, $\mu_n$ est unitaire à coefficients entiers, de degré $\varphi(n)$.
\end{prop}

\begin{proof}
On raisonne par récurrence sur $n$ pour le caractère entier, le reste est immédiat.
\begin{itemize}
    \item $n=1$ immédiat
    \item On se donne $n\geq 1$ et on suppose la propriété vraie jusqu'au rang $n$. Si $n+1$ est premier, la conclusion est immédiate. Sinon, \[
        X^{n+1}-1=\prod_{d\;|\;n+1}\mu_d(X)=\mu_{n+1}(X)\times
        \underbrace{\prod_{\substack{d\;|\;n+1\\ d\neq n+1}}\mu_d(X)}_{\in\mathbb Z[X], \text{ unitaire }}
    \]
    Donc $\mu_{n+1}$ est le quotient dans la division euclidienne de $X^{n+1}-1$ par un polynôme \textbf{unitaire} de $\mathbb Z[X]$, d'où la conclusion.
\end{itemize}
\end{proof}

\subsection{Expression de $\mu_n$}

On introduit la fonction arithmétique de Möbius\index{Möbius!fonction arithmétique} \[
    \mu(n)=\begin{cases}
        1 &\text{si } n=1 \\
        (-1)^k &\text{si }n\text{ produit de $k$ premiers distincts}\\
        0 &\text{sinon}
    \end{cases}
\]
Cette fonction est multiplicative: pour $a, b$ premiers entre eux, $\mu(ab)=\mu(a)\mu(b)$. Puis, pour $n\geq 1$, \[
    \sum_{d\;|\;n}\mu(d)=\begin{cases}
        1&\text{si }n=1 \\
        0&\text{sinon}
    \end{cases}
\]
En effet, si $n=p_1^{\alpha_1}\cdots p_k^{\alpha_k}$, \[
    \sum_{d\;|\;n}\mu(d)=\sum_{\substack{d\;|\; n\\\mu(d)\neq 0}}=\sum_{I\subset \llbracket 1, k\rrbracket}(-1)^{\#I}=\sum_{i=0}^k\binom ki(-1)^i=(1-1)^k=0
\]

\begin{res}
\[
    \mu_n(X)=\prod_{d\;|\;n}\left(X^d-1\right)^{\mu\left(\frac nd\right)}
\]
\end{res}

\begin{proof}
\begin{align*}
    \prod_{d\;|\; n}\left(X^d-1\right)^{\mu\left(\frac nd\right)}&=\prod_{d\;|\;n}\;\prod_{d'\;|\;d}\mu_{d'}(X)^{\mu\left(\frac nd\right)} \\
    &= \prod_{d'\;|\; n}\;\prod_{u\;|\;\frac n{d'}}\mu_{d'}(X)^{\mu\left(\frac n{d'u}\right)} \\
    &= \prod_{d'\;|\;n}\mu_{d'}(X)^{\left(\displaystyle\sum_{u|\frac n{d'}}\mu\left(\frac n{d'u}\right)\right)}\\
    &=\mu_n(X)
\end{align*}
car \[
\sum_{u\;|\;\frac n{d'}}\mu\left(\frac n{d'u}\right)=\sum_{u\;|\;\frac n{d'}}\mu(u)=\begin{cases}
    1& \text{si }\frac n{d'}=1\iff d'=n \\
    0& \text{sinon}
\end{cases}
\]
\end{proof}

\begin{rem}
Cela donne une autre preuve de $\mu_n(X)\in\mathbb Z[X]$
\end{rem}

\subsection{Suites arithmétiques}

On note $p$ un nombre premier et $m\geq 1$ un entier non divisible par $p$. Supposons que $a$ est un entier et que \[
    \mu_m(a)= 0\pmod p
\]
On a dans ce cas \[
    p\;|\;\mu_m(a)\;|\;a^m-1 \text{ premier avec }a \implies p\;\not|\;a
\]

Si $d=\ord a$ alors $d\;|\; m$ donc $X^d-1|X^m-1$. Par l'absurde, si $d\neq n$ alors $X^d-1\land \mu_m=1$ (aucune racine commune dans $\mathbb C$) donc (Gauss) $(X^d-1)\times \mu_m(X)|X^m-1$. On en déduit ($d$ est l'ordre de $a$) \[
    (a^d-1)\mu_m(a)\;|\;a^m-1 \implies a^m-1=0\pmod {p^2}
\]
et comme $a+p$ vérifie les mêmes hypothèses ($a+p\equiv a\pmod p$ donc $\mu_m(a+p)=0\pmod {p}$ et l'ordre est le même donc $\neq m$), on a $(a+p)^m-1= 0\pmod {p^2}$, d'où \[
    p^2\;|\; (a+p)^m-a^m=p\left(a^{m-1}+\underbrace{\cdots}_{\text{multiple de }p}\right) \implies p\;|\; a
\]
c'est absurde, donc $\ord a = m$. Puis $m=\ord a\;|\;p-1$ donc $p$ est un terme de la suite $(1+km)_{k\in\mathbb N}$. Si cette suite n'a qu'un nombre fini de nombres premiers $p_1, \cdots, p_s$ alors pour $k$ assez grand, $\mu_m(kp_1\cdots p_s)>2$. On a aussi $\mu_m(0)=\pm 1$ (entier de module $1$) donc si $p$ est un premier qui divise $\mu_m(kp_1\cdots p_s)$ et $p$ est parmi les $p_i$ alors $p\;|\; \mu_m(0)$ absurde, donc $p$ n'est pas parmi les $p_i$ et la suite a une infinité de termes premiers.

\subsection{Irréductibilité}

\index{polynôme!irréductible}

Nous allons montrer que les polynômes cyclotomiques sont irréductibles par un argument de Landau.

On note $\alpha$ une racine primirive $n$-ième de l'unité et $P$ son polynôme minimal. La suite $(P(\alpha^k))_{k\in\mathbb N}$ est $n$-périodique à valeurs dasn $\mathbb Z[\alpha]$ et chacun des termes s'écrit  de manière unique $Q_k(\alpha)$ pour un polynôme $Q_k$ de degré $<\deg P$.

On note $A$ le plus grand coefficient de tous les $Q_k$ en valeurs absolues. Dans $\mathbb Z_p[X]$, $P(X^p)=P(X)^p$ (Frobenius) donc il existe $U\in\mathbb Z[X]$ tel que \[
    P(X^p)=P(X)^p+pU(X)
\]
Pour $p>A$, \[
    P(\alpha^p)=P(\alpha)^p+pU(\alpha)=pU(\alpha)=pV(\alpha) \quad \text{ avec }\deg V<\deg P
\]
donc \[
    Q_p(\alpha)=pV(\alpha) \qquad \text{ et }\qquad P\;|Q_p-pV
\]
et donc \[
    Q_p=pV
\]
d'où $p\;|\;$ les coefs de $Q_p$ qui sont de module $\leq A<p$ donc $Q_p=0$ et $P(\alpha^p)=0$. Ici, on s'est seulement servis du fait que $\alpha$ est une racine de l'unité, donc si $p_1, \cdots, p_r$ sont $>A$ alors $P(\alpha^{p_1\cdots p_r})=0$.

On note $m$ premier avec $n$, $N$ le produit des premiers $\leq A$. On a $m+nN=m\pmod n$ et si $p$ divise $m+nN$ alors $p>A$ donc $P(\alpha^{m+nN})=0=P(\alpha^m)$

Ainsi les $(\alpha^{m})$ avec $m\land n=1$ sont racines de $P$ donc $\mu_n(X) \;|\; P(X)$. $\mu_n(\alpha)=0$ donc $P$ divise $\mu_m$, donc $P=\mu_m$. C'est un polynôme minimal donc irréductible.

\section{Nombres algébriques}
\index{nombres algébriques}

Si $K\subset L$ sont deux corps alors on dit que $a\in L$ est algébrique sur $K$ s'il existe $P\in\mathbb K[X]$ non nul tel que $P(a)=0$.

Dans $\mathbb C$, on appelle nombre algébrique les nombres algébriques sur $\mathbb Q$ et on note $A$ l'ensemble de ces nombres.

\begin{rem}
Si $a$ est algébrique sur un corps $K$ alors \[
    I_a=\{P\in\mathbb K[X]\;/\; P(a)=0\}
\]
est un idéal non réduit à $\{0\}$ de $K[X]$ euclidien donc principal, donc il existe un unique polynôme unitaire $\mu_a(X)$ tel que $I_a=\mu_a K[X]$, que l'on appelle polynôme minimal de $a$. Dans ce cas \[
    \Vect_Kr((a^k)_{k\in\mathbb N})=\Vect_K(1, a, \cdots, a^{\deg \mu_a-1})=\{P(a), P\in K[X]\}
\]
\end{rem}

\subsection{Structure de corps}
\index{nombres algébriques!structure}

\begin{res}
$A$ est un corps
\end{res}

\begin{proof}
\begin{itemize}
    \item $A\subset \mathbb C$
    \item $1\in A$ en tant que racine de $X-1$
    \item Soient $x, y\in A$, $p=\deg \mu_x, q=\deg \mu_y$. Pour tout $n\in\mathbb N$, \[
        \begin{cases}
            x^n\in\Vect_{\mathbb Q}(1, x, \cdots, x^{p-1}) \\ y^n\in\Vect_{\mathbb Q}(1, y, \cdots, y^{q-1})
        \end{cases}\quad\text{ donc }\quad  x^ny^n, (x+y)^n\in\Vect_{\mathbb Q}(x^iy^j, i<q, j<q)
    \]

    Les familles $((x-y)^n)_n$ et $((xy)^n)_n$ sont liées donc $x-y, xy\in A$.
    \item Si $x\in A$ non nul et si on note $E=\Vect_{\mathbb Q}(1, x, \cdots, x^{p-1})$ alors \[
        \varphi: z\in E\longmapsto xz\in E
    \]
    est un endomorphisme injectif en dimension finie donc c'est un isomorphisme et $\varphi^{-1}(1)$ donne $x^{-1}\in E\subset A$
\end{itemize}
\end{proof}

\subsection{Règle de multiplicativité des degrés}

\index{multiplicativité des degrés}

On note $a$ algébrique sur $K$ de degré $d$. \[
    \varphi: P\in K[X] \longmapsto P(a)
\]
est un morphisme de $K$-algèbres donc $\Img\varphi=\Vect_K(1, \cdots, a^{p-1})$ est une $K$-algèbre, notée $K[a]$.

$K[d]$ est une algèbre de dimension finie $d$ sur $K$ et pour $x\in K[a]\setminus \{0\}$, \[
    \psi: y\in K[a]\longmapsto xy\in K[a]
\]
est un endomorphime injectif en dimension finie donc c'est un automorphisme et $x$ a un inverse dans $K[a]$, qui est donc un corps.

\begin{dfn}
Si $K\subset L$ alors $L$ est un $K$-ev et si $L$ est de dimension finie sur $K$ alors on note \[
    [L:K]=\dim_K(L)
\]
\end{dfn}

\begin{prop}
    \index{nombres algébriques!transitivité}
\begin{enumerate}
    \item Soit $K\subset L\subset N$ trois corps. Si $\dim_L N, \dim_K L<+\infty$ alors $\dim_K N<+\infty$ et \[
        [N:K]=[N:L]\times [L:K]
    \]
    \item Si $a\in\mathbb C$ est algébrique et si $b\in K$ est algébrique sur $\mathbb Q[a]$ alors $b$ est algébrique sur $\mathbb Q$.
\end{enumerate}
\end{prop}

\begin{proof} ~
\begin{enumerate}
    \item $(n_1, \cdots, n_p)$ base de $N$ sur $L$, $(\ell_1, \cdots, \ell_q)$ base de $L$ sur $K$. Si $n\in\mathbb N$ alors il existe $\lambda_1, \cdots, \lambda_p\in L$ tel que $n=\lambda_1n_1+\cdots+\lambda_p n_p$ et chaque $\lambda_i\in\Vect_{\mathbb Q}(\ell_1, \cdots, \ell_p)$ donc $n\in\Vect_{\mathbb Q}(\ell_in_j)$ (famille génératrice finie) donc \[
        [N:K]<+\infty
    \]
    Vérifions que $(\ell_in_j)_{\substack{1\leq i\leq q\\ 1\leq j\leq p}}$ est libre sur $\mathbb Q$. \[
        \sum_{i, j}\lambda_{i, j}\ell_in_j=0\implies \sum_{j=1}^pn_j\sum_{i=1}^q\lambda_{i, j}\ell_i=0 \implies \forall j\in\llbracket 1, p\rrbracket \sum_{i=1}^q\lambda_{i, j}\ell_i=0\implies \forall i, j, \lambda_{i, j}=0
    \]
    \item C'est une conséquence directe de $1.$
\end{enumerate}
\end{proof}

\begin{rem}
Si $b$ est racine d'une équation polynômiale à coefficients dans $\mathbb Q[a]$, il l'est aussi d'une équation polynômiale à coefficients dans $\mathbb Q$.
\end{rem}

\section{Polynômes orthogonaux}
\index{polynômes orthogonaux}

On note $p: I\longrightarrow \mathbb R_+^\star$ continue telle que pour tout entier naturel $k$, $t\longmapsto t^kp(t)$ est intégrable sur $I$. Dans ce cas, \[
    (P|Q)=\int_Ip(t)P(t)Q(t)\;\mathrm dt
\]
est bien définit et est un produit scalaire sur $\mathbb R[X]$

\begin{prop}
Il existe une unique famille $(P_k)_k$ de polynômes unitaires tels que
\begin{itemize}
    \item $\forall k\in\mathbb N, \quad \deg P_k=k$
    \item les $P_k$ sont deux à deux orthogonaux
\end{itemize}
\end{prop}

\begin{proof}
    Le procédé de Gram-Schmidt appliqué à la base canonique donne l'existence. Si $(P_k)$ convient alors
\begin{itemize}
    \item $P_0=1$ (unitaire de degré $0$)
    \item Si $P_0, \cdots, P_k$ conviennent, alors \[
        P_{k+1}=\underbrace{X^{k+1}}_{unitaire} + \underbrace{\lambda_0P_0+\cdots +\lambda_k P_k}_{\deg \leq k \text{ se décompose sur }P_0, \cdots, P_k}
    \]
et l'orthogonalité détermine de manière unique les $\lambda_i$
\end{itemize}
\end{proof}

\subsection{Récurrence d'ordre 2}

\begin{prop}
La famille des $(P_n)_n$ vérifie pour $n\geq 2$ \[
    P_n(X)=-(X+\lambda_{n-1})P_{n-1}(X)-\mu_{n-2}P_{n-2}(X)=0
\]
avec \[
    \lambda_{n-1}=-\frac{(XP_{n-1}|P_{n-1})}{\|P_{n-1}\|^2}, \quad \mu_{n-2}=-\frac{\|P_{n-1}\|^2}{\|P_{n-2}\|^2}
\]
\end{prop}

\begin{proof} On note $E_i=\mathbb R_i[X]$.
Pour $Q\in\Vect(P_0, \cdots, P_{n-2})=\mathbb R_{n-2}[X]$, $(XP_n|Q)=(P_n|XQ)=0$ car $XQ\in\mathbb R_{n-1}[X]$. On a alors \[
    P_n(X)-XP_{n-1}(X)\in E_{n-3}^\bot
\]
or $\deg (P_n-XP_{n-1})\leq n-1$ (simplification du coef dominant) donc \[
    P_n(X)-XP_{n-1}(X)=\lambda_{n-1}P_{n-1}(X)+\cdots +\lambda_0P_0(X)
\]
et on a $P_n-XP_{n-1}\bot P_0, \cdots, P_{n-3}$ donc $\lambda_0=\cdots =\lambda_{n-3}=0$. On en déduit, \[
    \lambda_{n-1}(P_{n-1}|P_{n-1})=(P_n|P_{n-1})-(XP_{n-1}|P_{n-1})
\]
et \[
    \lambda_{n-2}(P_{n-2}|P_{n-2})=\underbrace{(P_n|P_{n-2})}_{=0}-(XP_{n-1}|P_{n-2})=-(P_{n-1}|\underbrace{XP_{n-2}}_{P_{n-1}+\in E_{n-2}})=-(P_{n-1}|P_{n-1})
\]
d'où la conclusion avec $\mu_{n-2}=\lambda_{n-2}$
\end{proof}

\subsection{Changements de signes}

On appelle changement de signe d'un polynôme un point $x_0$ tel que $(x-x_0)P(x)$ est de signe constant au voisinage de $x_0$. On note $N_I(P)$ le nombre de changements de signes de $P$ sur l'intervalle $I$

Supposons $N_I(P_n)=p<n$. On note $x_1, \cdots, x_p$ les changements de signes et \[
    (P_n | \underbrace{(X-x_1)\cdots (X-x_p)}_{\in E_{n-1}})=0=\int_Ip(t)\underbrace{P_n(t)(t-x_1)\cdots (t-x_p)}_{\text{signe constant}}\;\mathrm dt
\]
donc \[
    \forall t\in I, p(t)P_n(t)(t-x_1)\cdots (t-x_p)=0
\]
et $P_n$=0 absurde. Donc $N_I(P_n)=n$ et $P_n$ est scindé à racines simples, toutes dans $I$.

\subsection{Entrelacement des racines}

On note \[
    w_n(x)=\left|\begin{matrix}P_n(x) & P_{n+1}(x)\\ P_n'(x) & P_{n+1}'(x)\end{matrix}\right|=P_n(x)P_{n+1}'(x)-P_n'(x)P_{n+1}(x)
\]
On a $w_0\equiv 1$ et après calcul avec la relation de récurrence, \[
    w_n(x)=P_n(x)^2\underbrace{-\mu_{n-1}}_{>0}w_{n-1}(x)
\]
donc comme $w_0>0$, on en déduit que pour tout $x, n$, $w_n(x)>0$. Notons $\alpha < \beta$ deux racines consécutives de $P_{n+1}$. \[
    P_{n+1}'(\alpha)P_n(\alpha)>0\qquad \text{ et }\qquad P_{n+1}'(\beta)P_n(\beta)>0
\] or $P_{n+1}$ de signe constant sur $]\alpha, \beta[$ et donc $P_{n+1}'(\alpha)P_{n+1}'(\beta)<0$ (car sinon il y aurait changement de signe). Ainsi, $P_n(\alpha)P_n(\beta)<0$ donc $P_n$ s'annule sur $]\alpha, \beta[$. On a ainsi trouvé $n$ racines de $P_n$, une entre chaque racine de $P_{n+1}$: il y a entrelacement des racines

\section{Formules d'interpolation}

\index{Lagrange!interpolation}

\subsection{Rappels}

On considère $x_1, \cdots, x_n\in\mathbb K$ deux à deux distincts, et on note \[
    \Delta_i(X)=\prod_{j\neq i}\frac{X-x_j}{x_i-x_j}
\]
de telle manière que \[
    \Delta_i(x_j)=\delta_{i, j}
\]

\begin{thm}
    On considère $x_1, \cdots, x_n, y_1, \cdots, y_n\in\mathbb K$, les $x_i$ deux à deux distincts. Dans ce cas, il existe un unique polynôme $P\in\mathbb K_{n-1}[X]$ tel que \[\forall i\in\llbracket 1, n\rrbracket, \quad P(x_i)=y_i\]
    De plus, \[
        P(X)=\sum_{i=1}^ny_i\Delta_i(X)
    \]
\end{thm}

\begin{proof}
Le polynôme donné convient. Si $P$ et $Q$ conviennent, $P-Q$ a $n$ racines et est de degré au plus $n-1$ donc $P=Q$.
\end{proof}

\begin{rem}
Si on note \[
    \omega (x)=\prod_{j=1}^n(x-x_j)
\]
alors \[
    \prod_{j\neq i}(x_i-x_j)=\lim_{x\to x_i}\frac{\omega(x)-\omega(x_i)}{x-x_i}=w'(x_i)
\]
et \[
    \Delta_i(x)=\frac{\omega (x)}{\omega(x_i)(x-x_i)}
\]
donc \[
    P(x)=\underbrace{\left(\sum_{i=1}^n\frac{y_i}{\omega'(x_i)(x-x_i)}\right)}_{\text{DSE de }\frac{P(x)}{\omega(x)}}\omega(x)
\]
\end{rem}

\subsection{Formules d'ordre supérieur}

On se donne $n_1, \cdots, n_r\in\mathbb N^\star$, $x_1, \cdots, x_r\in\mathbb K$ deux à deux distincts et $n=n_1+\cdots+n_r$. \[
    \varphi: P\in\mathbb K_{n-1}[X]\longmapsto(P(x_1), \cdots, P^{(n_1-1)}(x_1), \cdot, P(x_r), \cdots, P^{(n_r-1)})\in\mathbb K^n
\]
est une application linéaire injective, il y a égalité des dimensions donc c'est un isomorphisme et on en déduit que pour $(y_{i, k})_{i\leq r, j\leq n_i}$ donné, il existe un unique polynôme de degré au plus $n-1$ tel que \[
    \forall i, k, P^{(k)}(x_i)=y_{i, k}
\]

\subsection{Analyse de l'erreur}

On se donne $f\in\mathbb C^n([a, b], \mathbb R) ,x_1, \cdots, x_n\in [a, b]$ deux à deux distincts. \[
    P(x)=\sum_{i=1}^nf (x_i)\Delta_i(x)
\]
l'interpolation de Lagrange de $f$. On veut estimer $f(x)-P(x)$. On note $x\in [a, b]\setminus\{x_1, \cdots, x_r\}$ et on note \[
    g(t)=f(t)-P(t)-\lambda_x\omega(t)
\]
où $\lambda_x$ tel que $g(x)=0$ (ajout d'une racine supplémentaire). Rolle itéré donne:

\begin{itemize}
    \item
$g$ s'annule en $x_1, \cdots, x_n, x$ ($n+1$ points)
\item $\cdots$
\item
$g^{(n)}$ s'annule au moins une fois en $c_n$
\end{itemize}

\[
    0=f^{(n)}(c_n)-P^{(n)}(c_n)-\lambda_x n!
\]
donc $\lambda_x = \dfrac{f^{(n)}(c_n)}{n!}$ et donc \[
    |f(x)-P(x)|=|\lambda_x\omega (x)|\leq |\omega(x)|\sup_{[a, b]}\frac{|f^{(n)}|}{n!}
\]
cette inégalité est encore vraie pour $x\in\{x_1, \cdots, x_n\}$

\begin{rem}
On constate que l'erreur est meilleure si \[
    |\omega(x)|=|(x-x_1)\cdots (x-x_n)|
\]
est petit. On peut donc s'intéresser aux choix des $x_i$ minimisant $\|\omega\|_\infty$ sur $[a, b]$. C'est un problème difficile, mais le choix de $x_1, \cdots, x_n$ équirépartis est mauvais.
\end{rem}

\section{Suites de polynômes}

\begin{thm}[Weierstrass]
    \index{Weierstrass (théorème de -- )}
    Si $f$ est continue de $[a, b]$ dans $\mathbb C$ alors il existe une suite $(P_n)$ de fonctions polynômiales telle que \[
        P_n\xrightarrow[{[a, b]}]{CVU} f
    \]
\end{thm}

\begin{proof}[Idée de la preuve]
    Bernstein puis convolution
\end{proof}

\begin{rem}
Le résultat est faux sur un intervalle autre qu'un segment.
\end{rem}

\subsection{Inégalité de Bernstein}

\index{Bernstein!inégalité}

\begin{lmm}
Si $z_1, \cdots, z_n$ sont les racines du polynôme $z^n+1$ et $P$ polynôme réel de degré $n$, alors pour tout $t\in\mathbb C$, \[
    tP'(t)=\frac n2P(t)+\frac 1n\sum_{k=1}^nP(-z_k)\frac{2z_k}{(z_k-1)^2}
\]
\end{lmm}

\begin{proof}
On note $\varphi(z, t)=\frac{P(tz)-P(t)}{z-1}$ pour $z\neq 1$. C'est une fonction polynomiale de degré $\leq n-1$ pour $t$ fixé.

D'après la formule d'interpolation de Lagrange en $z_1, \cdots, z_n$, \begin{align*}
    \varphi(z, t)&=\sum_{k=1}^n\varphi(z_k, t)\frac{z^n+1}{z-z_k} \times \underbrace{\frac1{(z_k-z_1)\cdots \widehat{(z_k-z_k)}\cdots (z_k-z_n)}}_{=\lim_{z\to z_k}\limits \frac{z-z_k}{z^n+1}=\frac1{nz_k^{n-1}}=-\frac{z_k}n} \\ &= \sum_{k=1}^n-\varphi(z_k, t)\times \frac{z^n+1}{z-z_k}\frac{z_k}n
\end{align*}

On fait tendre $z$ vers $1$. \[
    tP'(t)=\sum_{k=1}^n\varphi(z_k, t)\frac{2}{z_k-1}\frac{z_k}n=
    \frac1n\sum_{k=1}^nP(-tz_k)\frac{2z_k}{(z_k-1)^2}-\frac{P(t)}{n}\sum_{k=1}^n\frac{2z_k}{(z_k-1)^2}
\]
Il nous reste à évaluer ce dernier terme. On reprend la formule précédente avec $P(t)=t^n$. \[
    nt^n=\frac1n\sum_{k=1}^n2t^n\frac{z_k^{n+1}}{(z_k-1)^2}-\frac{t^n}{n}\sum_{k=1}^n\frac{2z_k}{(z_k-1)^2}\quad \underset{z_k^{n+1}=-z_k}{=}\quad -\frac{2t^n}n\sum_{k=1}^n\frac{2z_k}{(z_k-1)^2}
\]
Et alors,
\[
\sum_{k=1}^n\frac{2z_k}{(z_k-1)^2}=-\frac{n^2}2
\]
ce qui termine la preuve de la formule.
\end{proof}

\begin{res}
On note pour $P\in\mathbb R[X]$, $\|P\|=\sup_{|z|=1}\limits|P(z)|$. Si $P$ est de degré $n$ alors \[
    \|P'\|\leq n\|P\|
\]
\end{res}

\begin{proof}
On prend $t$ de module $1$, \[
    |tP'(t)|=|P'(t)|\leq \frac n2|P(t)|+\frac1n\sum_{k=1}^n\left|\frac{2z_k}{(z_k-1)^2}\right|\times \|P\|
\]
Et $z_k\in\mathbb U$ donc s'écrit $z=e^{i\theta}$ et donc ($\theta\in ]0, 2\pi[$) \[
    \frac{2z_k}{(z_k-1)^2}=\frac{2}{(e^{i\frac \theta 2}-e^{-i\frac \theta2})^2}=\frac 2{-4\sin^2\frac\theta 2}<0
\]
et \[
    \sum_{k=1}^n\left|\frac{2z_k}{(z_k-1)^2}\right|=-\sum_{k=1}^n\frac{2z_k}{(z_k-1)^2}=\frac{n^2}2
\]
ce qui termine la preuve.
\end{proof}

\subsection{Théorème d'oscillation de Tchebychev}

\index{Tchebychev!théorème d'oscillation}
\index{Tchebychev!polynômes}

On note $C_n$ le $n$-ième polynôme de Tchebychev et $T_n=\dfrac{C_n}{2^{n-1}}$ de sorte que $T_n$ est unitaire.

\begin{res}
Pour tout $P$ unitaire de degré $n$, \[
    \frac 1{2^{n-1}}=\max_{[-1, 1]}|T_n|\leq \max_{[-1, 1]}|P|
\]
\end{res}

\begin{proof}
\[
    T_n(\cos \theta)=\frac1{2^{n-1}}\cos (n\theta) \implies \max_{[-1, 1]}|T_n|=\frac1{2^{n-1}}
\]
Supposons par l'absurde que $\displaystyle\max_{[-1, 1]}|P|<\frac1{2^{n-1}}$ et posons \[
    Q(X)=T_n(X)-P(X)
\]
Pour $x_k=\cos\left(\frac{k\pi}n\right), k\in\llbracket 0, n\rrbracket$, on a \[
    T_n(x_k)=\frac{(-1)^k}{2^{n-1}}
\]
d'où on tire le tableau de signes
\begin{center}
    \begin{tikzpicture}
        \tkzTabInit[espcl=2]{$x$ / 1 ,$Q(x)$ /1 }{$x_0$ , $x_1$, $x_2$, $\cdots$ , $x_n$ }%
        \tkzTabLine{+,0,-,0,+,0,\cdots ,0,(-1)^n\;}
    \end{tikzpicture}
\end{center}
On en déduit que $Q$ s'annule au moins $n$ fois et $\deg Q\leq n-1$ donc $Q=0$, donc $P=T_n$ absurde.
\end{proof}

% \todo{exo fermés}

\section{Aspects topologiques}

\subsection{Le théorème de Baire}

\begin{res}
    \index{Baire (théorème de -- )}
    \index{Cauchy!suites}
    \index{espace complet}
    $(E, \|\;\|)$ un e.v.n tel que toute suite de Cauchy converge (complet). Si $(O_n)_n$ est une suite d'ouverts denses de $E$ alors $\cap O_n$ est dense.
\end{res}

\begin{proof}
On note $U$ ouvert de $E$.
\begin{itemize}
    \item
$U\cap O_0$ est un ouvert non vide donc il existe $x_0\in U\cap O_0$ et $2a_0>0$ tel que $\mathcal B_o(x_o, 2a_o)\subset U\cap O_0$ et donc $\mathcal B_f(x_0, a_0)\subset U\cap O_0$. On note $r_0=\min(a_0, 1)$.
\item $\mathcal B_o(a_0, r_0)\cap O_1$ ouvert non vide donc il existe $x_1$ et $a_1>0$ tels que $\mathcal B_f(x_1, a_1)\subset \mathcal B_0(x_0, r_0)\cap O_1$ et on note $r_1=\min\left(a_1, \frac12\right)$

\item $\cdots$
\end{itemize}
On construit ainsi deux suites $(x_n)_n\in E^{\mathbb N}$, $(r_n)\in\mathbb {R_+^\star}^{\mathbb N}$ telles que \[
    \begin{cases}
        \mathcal B_f(x_n, r_n)\subset \mathcal B_f(x_{n-1}, r_{n-1})\subset \cdots \subset \mathcal B_f(x_0, r_0) \\
        \mathcal B_f(x_n, r_n)\subset O_n, \quad r_n\leq \dfrac 1n
    \end{cases}
\]
Pour $m>n$, $x_m\in\mathcal B_f(x_m, r_m)\subset\mathcal B_f(x_n, r_n)$ donc \[
    \|x_m-x_n\|\leq \frac 1n
\]
donc $(x_n)$ de Cauchy donc convergente vers $x$.

Pour $m$ fixé, APCR $x_n\in\mathcal B_f(x_m, r_m)$ fermé donc $x\in\mathcal B_f(x_m, r_m)\subset O_m$ donc $x\in O_m$ pour tout $m$. Ainsi, \[
    x\in\bigcap_{n\in\mathbb N}O_n
\]
et $x\in\mathcal B_f(x_0, r_0)\subset U$ donc $x\in U$ et donc \[
    U\cap\left(\bigcap_{n\in\mathbb N}O_n\right)
\]
pour tout ouvert $U$.

\begin{cor}
Si $(F_n)_n$ est une suite de fermés d'intérieurs vides alors \[
    \bigcup_{n\in\mathbb N}F_n
\] est d'intérieur vide.
\end{cor}
\end{proof}

\subsection{Incomplétude de {$\mathbb R[X]$}}

On supose qu'il existe une norme $\|\cdot \|$ telle que $(\mathbb R[X], \|\cdot \|)$ est complet (i.e. toutes les suites de Cauchy convergent)

\begin{itemize}
    \item $E_n=\Vect(1, \cdots, X^n)$ est un fermé (s.e.v. de dimension finie)
    \item \[
    \forall P\in E_n, \forall \varepsilon>0, \quad P+\frac{X^{n+1}}k\xrightarrow[k\to+\infty]{} P
    \] donc APCR \[
        P+\frac{X^{n+1}}k \in \mathcal B_0(P, \varepsilon)
    \]
    donc $\mathcal B_o(P, \varepsilon)\not\subset E_n$ et $E_n$ d'intérieur vide.
    \item Par le théorème de Baire, \[
        \underbrace{\bigcap_{n\in\mathbb N} E_n}_{\text{intérieur vide}}=\underbrace{\mathbb R[X]}_{\text{intérieur non vide}}
    \]
    absurde.
\end{itemize}

Il n'existe donc aucune norme pour laquelle $\mathbb R[X]$ est complet.

En particulier, $(X_k)_k$ est une base qu'on peut orthonormaliser par le procédé de Gram-Schmidt en une base $(P_k)_k$ qui est totale, sans que $\mathbb R[X]$ soit complet.
\endchapter


\newpage
\pagestyle{empty}

\printindex

\end{document}
